\subsection*{Textbook Formal Inventory}
This subsection records the exact formal material in \file{HAETAE_ROM_Programming.ec}.  It is generated from the proof-surface EasyCrypt file, so the line numbers and statements are meant to be read next to the checked source.

\noindent\textbf{Chapter role.} lazy ROM programming, counted bad events, and programming bounds.

\noindent\textbf{Inventory size.} 19 context declarations, 68 definitions/game objects, 99 lemmas or relational theorems, and 0 explicit Hoare/Phoare judgments were found.

\subsubsection*{Theory Context, Imports, and Quantified Modules}
\paragraph{Context 1 (line 1)}
\noindent\textbf{EasyCrypt name.}
\begin{lstlisting}[language=EasyCrypt,basicstyle=\ttfamily\scriptsize]
imported theories
\end{lstlisting}
\noindent\textbf{Checked EasyCrypt statement.}
\begin{lstlisting}[language=EasyCrypt,basicstyle=\ttfamily\scriptsize]
require import AllCore Real Distr List FSet PROM StdOrder Mu_mem.
\end{lstlisting}
\noindent\textbf{Formal mathematical reading.}
Mathematical context. This line imports or specializes earlier theories, making their carrier sets, functions, modules, and theorems available to the current file.

\noindent\textbf{Role in the proof.}
Use in this chapter. It fixes the proof context, imported mathematical language, or quantified module parameters for the following statements.

\paragraph{Context 2 (line 2)}
\noindent\textbf{EasyCrypt name.}
\begin{lstlisting}[language=EasyCrypt,basicstyle=\ttfamily\scriptsize]
imported theories
\end{lstlisting}
\noindent\textbf{Checked EasyCrypt statement.}
\begin{lstlisting}[language=EasyCrypt,basicstyle=\ttfamily\scriptsize]
require import HAETAE_Params HAETAE_Algebra HAETAE_Distributions.
\end{lstlisting}
\noindent\textbf{Formal mathematical reading.}
Mathematical context. This line imports or specializes earlier theories, making their carrier sets, functions, modules, and theorems available to the current file.

\noindent\textbf{Role in the proof.}
Use in this chapter. It fixes the proof context, imported mathematical language, or quantified module parameters for the following statements.

\paragraph{Context 3 (line 3)}
\noindent\textbf{EasyCrypt name.}
\begin{lstlisting}[language=EasyCrypt,basicstyle=\ttfamily\scriptsize]
imported theories
\end{lstlisting}
\noindent\textbf{Checked EasyCrypt statement.}
\begin{lstlisting}[language=EasyCrypt,basicstyle=\ttfamily\scriptsize]
require import HAETAE_ROM HAETAE_Transcript.
\end{lstlisting}
\noindent\textbf{Formal mathematical reading.}
Mathematical context. This line imports or specializes earlier theories, making their carrier sets, functions, modules, and theorems available to the current file.

\noindent\textbf{Role in the proof.}
Use in this chapter. It fixes the proof context, imported mathematical language, or quantified module parameters for the following statements.

\paragraph{Context 4 (line 4)}
\noindent\textbf{EasyCrypt name.}
\begin{lstlisting}[language=EasyCrypt,basicstyle=\ttfamily\scriptsize]
imported theories
\end{lstlisting}
\noindent\textbf{Checked EasyCrypt statement.}
\begin{lstlisting}[language=EasyCrypt,basicstyle=\ttfamily\scriptsize]
require import HAETAE_Rejection.
\end{lstlisting}
\noindent\textbf{Formal mathematical reading.}
Mathematical context. This line imports or specializes earlier theories, making their carrier sets, functions, modules, and theorems available to the current file.

\noindent\textbf{Role in the proof.}
Use in this chapter. It fixes the proof context, imported mathematical language, or quantified module parameters for the following statements.

\paragraph{Context 5 (line 5)}
\noindent\textbf{EasyCrypt name.}
\begin{lstlisting}[language=EasyCrypt,basicstyle=\ttfamily\scriptsize]
imported theories
\end{lstlisting}
\noindent\textbf{Checked EasyCrypt statement.}
\begin{lstlisting}[language=EasyCrypt,basicstyle=\ttfamily\scriptsize]
require import HAETAE_Reductions.
\end{lstlisting}
\noindent\textbf{Formal mathematical reading.}
Mathematical context. This line imports or specializes earlier theories, making their carrier sets, functions, modules, and theorems available to the current file.

\noindent\textbf{Role in the proof.}
Use in this chapter. It fixes the proof context, imported mathematical language, or quantified module parameters for the following statements.

\paragraph{Context 6 (line 7)}
\noindent\textbf{EasyCrypt name.}
\begin{lstlisting}[language=EasyCrypt,basicstyle=\ttfamily\scriptsize]
HAETAE_ROM_Programming
\end{lstlisting}
\noindent\textbf{Checked EasyCrypt statement.}
\begin{lstlisting}[language=EasyCrypt,basicstyle=\ttfamily\scriptsize]
theory HAETAE_ROM_Programming.
\end{lstlisting}
\noindent\textbf{Formal mathematical reading.}
Mathematical context. This begins a namespace for the formal development in this file.

\noindent\textbf{Role in the proof.}
Use in this chapter. It fixes the proof context, imported mathematical language, or quantified module parameters for the following statements.

\paragraph{Context 7 (line 9)}
\noindent\textbf{EasyCrypt name.}
\begin{lstlisting}[language=EasyCrypt,basicstyle=\ttfamily\scriptsize]
opened namespace
\end{lstlisting}
\noindent\textbf{Checked EasyCrypt statement.}
\begin{lstlisting}[language=EasyCrypt,basicstyle=\ttfamily\scriptsize]
import HAETAE_Params.
\end{lstlisting}
\noindent\textbf{Formal mathematical reading.}
Mathematical context. This line imports or specializes earlier theories, making their carrier sets, functions, modules, and theorems available to the current file.

\noindent\textbf{Role in the proof.}
Use in this chapter. It fixes the proof context, imported mathematical language, or quantified module parameters for the following statements.

\paragraph{Context 8 (line 10)}
\noindent\textbf{EasyCrypt name.}
\begin{lstlisting}[language=EasyCrypt,basicstyle=\ttfamily\scriptsize]
opened namespace
\end{lstlisting}
\noindent\textbf{Checked EasyCrypt statement.}
\begin{lstlisting}[language=EasyCrypt,basicstyle=\ttfamily\scriptsize]
import HAETAE_Algebra.
\end{lstlisting}
\noindent\textbf{Formal mathematical reading.}
Mathematical context. This line imports or specializes earlier theories, making their carrier sets, functions, modules, and theorems available to the current file.

\noindent\textbf{Role in the proof.}
Use in this chapter. It fixes the proof context, imported mathematical language, or quantified module parameters for the following statements.

\paragraph{Context 9 (line 11)}
\noindent\textbf{EasyCrypt name.}
\begin{lstlisting}[language=EasyCrypt,basicstyle=\ttfamily\scriptsize]
opened namespace
\end{lstlisting}
\noindent\textbf{Checked EasyCrypt statement.}
\begin{lstlisting}[language=EasyCrypt,basicstyle=\ttfamily\scriptsize]
import HAETAE_Distributions.
\end{lstlisting}
\noindent\textbf{Formal mathematical reading.}
Mathematical context. This line imports or specializes earlier theories, making their carrier sets, functions, modules, and theorems available to the current file.

\noindent\textbf{Role in the proof.}
Use in this chapter. It fixes the proof context, imported mathematical language, or quantified module parameters for the following statements.

\paragraph{Context 10 (line 12)}
\noindent\textbf{EasyCrypt name.}
\begin{lstlisting}[language=EasyCrypt,basicstyle=\ttfamily\scriptsize]
opened namespace
\end{lstlisting}
\noindent\textbf{Checked EasyCrypt statement.}
\begin{lstlisting}[language=EasyCrypt,basicstyle=\ttfamily\scriptsize]
import HAETAE_ROM.
\end{lstlisting}
\noindent\textbf{Formal mathematical reading.}
Mathematical context. This line imports or specializes earlier theories, making their carrier sets, functions, modules, and theorems available to the current file.

\noindent\textbf{Role in the proof.}
Use in this chapter. It fixes the proof context, imported mathematical language, or quantified module parameters for the following statements.

\paragraph{Context 11 (line 13)}
\noindent\textbf{EasyCrypt name.}
\begin{lstlisting}[language=EasyCrypt,basicstyle=\ttfamily\scriptsize]
opened namespace
\end{lstlisting}
\noindent\textbf{Checked EasyCrypt statement.}
\begin{lstlisting}[language=EasyCrypt,basicstyle=\ttfamily\scriptsize]
import HAETAE_Transcript.
\end{lstlisting}
\noindent\textbf{Formal mathematical reading.}
Mathematical context. This line imports or specializes earlier theories, making their carrier sets, functions, modules, and theorems available to the current file.

\noindent\textbf{Role in the proof.}
Use in this chapter. It fixes the proof context, imported mathematical language, or quantified module parameters for the following statements.

\paragraph{Context 12 (line 14)}
\noindent\textbf{EasyCrypt name.}
\begin{lstlisting}[language=EasyCrypt,basicstyle=\ttfamily\scriptsize]
opened namespace
\end{lstlisting}
\noindent\textbf{Checked EasyCrypt statement.}
\begin{lstlisting}[language=EasyCrypt,basicstyle=\ttfamily\scriptsize]
import HAETAE_Rejection.
\end{lstlisting}
\noindent\textbf{Formal mathematical reading.}
Mathematical context. This line imports or specializes earlier theories, making their carrier sets, functions, modules, and theorems available to the current file.

\noindent\textbf{Role in the proof.}
Use in this chapter. It fixes the proof context, imported mathematical language, or quantified module parameters for the following statements.

\paragraph{Context 13 (line 15)}
\noindent\textbf{EasyCrypt name.}
\begin{lstlisting}[language=EasyCrypt,basicstyle=\ttfamily\scriptsize]
opened namespace
\end{lstlisting}
\noindent\textbf{Checked EasyCrypt statement.}
\begin{lstlisting}[language=EasyCrypt,basicstyle=\ttfamily\scriptsize]
import HAETAE_Reductions.
\end{lstlisting}
\noindent\textbf{Formal mathematical reading.}
Mathematical context. This line imports or specializes earlier theories, making their carrier sets, functions, modules, and theorems available to the current file.

\noindent\textbf{Role in the proof.}
Use in this chapter. It fixes the proof context, imported mathematical language, or quantified module parameters for the following statements.

\paragraph{Context 14 (line 16)}
\noindent\textbf{EasyCrypt name.}
\begin{lstlisting}[language=EasyCrypt,basicstyle=\ttfamily\scriptsize]
opened namespace
\end{lstlisting}
\noindent\textbf{Checked EasyCrypt statement.}
\begin{lstlisting}[language=EasyCrypt,basicstyle=\ttfamily\scriptsize]
import RealOrder.
\end{lstlisting}
\noindent\textbf{Formal mathematical reading.}
Mathematical context. This line imports or specializes earlier theories, making their carrier sets, functions, modules, and theorems available to the current file.

\noindent\textbf{Role in the proof.}
Use in this chapter. It fixes the proof context, imported mathematical language, or quantified module parameters for the following statements.

\paragraph{Context 15 (line 315)}
\noindent\textbf{EasyCrypt name.}
\begin{lstlisting}[language=EasyCrypt,basicstyle=\ttfamily\scriptsize]
CountedLazyROMStateFacts
\end{lstlisting}
\noindent\textbf{Checked EasyCrypt statement.}
\begin{lstlisting}[language=EasyCrypt,basicstyle=\ttfamily\scriptsize]
section CountedLazyROMStateFacts.
\end{lstlisting}
\noindent\textbf{Formal mathematical reading.}
Mathematical context. This opens a scoped block of hypotheses, module parameters, and lemmas.

\noindent\textbf{Role in the proof.}
Use in this chapter. It fixes the proof context, imported mathematical language, or quantified module parameters for the following statements.

\paragraph{Context 16 (line 317)}
\noindent\textbf{EasyCrypt name.}
\begin{lstlisting}[language=EasyCrypt,basicstyle=\ttfamily\scriptsize]
O
\end{lstlisting}
\noindent\textbf{Checked EasyCrypt statement.}
\begin{lstlisting}[language=EasyCrypt,basicstyle=\ttfamily\scriptsize]
declare module O <: Oracle {-CountedLazyROM}.
\end{lstlisting}
\noindent\textbf{Formal mathematical reading.}
Mathematical context. This is universal quantification over an adversary, oracle, scheme, sampler, or reduction module satisfying the stated interface and memory restrictions.

\noindent\textbf{Role in the proof.}
Use in this chapter. It fixes the proof context, imported mathematical language, or quantified module parameters for the following statements.

\paragraph{Context 17 (line 1455)}
\noindent\textbf{EasyCrypt name.}
\begin{lstlisting}[language=EasyCrypt,basicstyle=\ttfamily\scriptsize]
CountedLazyROMBadFlags
\end{lstlisting}
\noindent\textbf{Checked EasyCrypt statement.}
\begin{lstlisting}[language=EasyCrypt,basicstyle=\ttfamily\scriptsize]
section CountedLazyROMBadFlags.
\end{lstlisting}
\noindent\textbf{Formal mathematical reading.}
Mathematical context. This opens a scoped block of hypotheses, module parameters, and lemmas.

\noindent\textbf{Role in the proof.}
Use in this chapter. It fixes the proof context, imported mathematical language, or quantified module parameters for the following statements.

\paragraph{Context 18 (line 1457)}
\noindent\textbf{EasyCrypt name.}
\begin{lstlisting}[language=EasyCrypt,basicstyle=\ttfamily\scriptsize]
O
\end{lstlisting}
\noindent\textbf{Checked EasyCrypt statement.}
\begin{lstlisting}[language=EasyCrypt,basicstyle=\ttfamily\scriptsize]
declare module O <: Oracle.
\end{lstlisting}
\noindent\textbf{Formal mathematical reading.}
Mathematical context. This is universal quantification over an adversary, oracle, scheme, sampler, or reduction module satisfying the stated interface and memory restrictions.

\noindent\textbf{Role in the proof.}
Use in this chapter. It fixes the proof context, imported mathematical language, or quantified module parameters for the following statements.

\paragraph{Context 19 (line 1458)}
\noindent\textbf{EasyCrypt name.}
\begin{lstlisting}[language=EasyCrypt,basicstyle=\ttfamily\scriptsize]
A
\end{lstlisting}
\noindent\textbf{Checked EasyCrypt statement.}
\begin{lstlisting}[language=EasyCrypt,basicstyle=\ttfamily\scriptsize]
declare module A <: CountedROMClient {-O, -CountedLazyROM, -CountedLazyROMBadGame}.
\end{lstlisting}
\noindent\textbf{Formal mathematical reading.}
Mathematical context. This is universal quantification over an adversary, oracle, scheme, sampler, or reduction module satisfying the stated interface and memory restrictions.

\noindent\textbf{Role in the proof.}
Use in this chapter. It fixes the proof context, imported mathematical language, or quantified module parameters for the following statements.

\subsubsection*{Definitions, Carrier Sets, Operations, Games, and Procedures}
\begin{formaldefinition}[EasyCrypt line 18]
\noindent\textbf{EasyCrypt name.}
\begin{lstlisting}[language=EasyCrypt,basicstyle=\ttfamily\scriptsize]
programming_site
\end{lstlisting}
\noindent\textbf{Checked EasyCrypt statement.}
\begin{lstlisting}[language=EasyCrypt,basicstyle=\ttfamily\scriptsize]
type programming_site = ro_query * ro_output.
\end{lstlisting}
\noindent\textbf{Formal mathematical reading.}
Mathematical definition. This is a definitional type abbreviation.  The exact displayed EasyCrypt statement gives the right-hand side; later proof steps may unfold it without adding an assumption.

\noindent\textbf{Role in the proof.}
Use in this chapter. It is part of the exact formal vocabulary used by subsequent lemmas and games.

\end{formaldefinition}

\begin{formaldefinition}[EasyCrypt line 20]
\noindent\textbf{EasyCrypt name.}
\begin{lstlisting}[language=EasyCrypt,basicstyle=\ttfamily\scriptsize]
programming_site_query
\end{lstlisting}
\noindent\textbf{Checked EasyCrypt statement.}
\begin{lstlisting}[language=EasyCrypt,basicstyle=\ttfamily\scriptsize]
op programming_site_query (site : programming_site) : ro_query = site.`1.
\end{lstlisting}
\noindent\textbf{Formal mathematical reading.}
Mathematical definition. This is a total EasyCrypt operation.  The exact displayed statement gives the domain, codomain, and defining equation when present.

\noindent\textbf{Role in the proof.}
Use in this chapter. It is part of the exact formal vocabulary used by subsequent lemmas and games.

\end{formaldefinition}

\begin{formaldefinition}[EasyCrypt line 21]
\noindent\textbf{EasyCrypt name.}
\begin{lstlisting}[language=EasyCrypt,basicstyle=\ttfamily\scriptsize]
programming_site_output
\end{lstlisting}
\noindent\textbf{Checked EasyCrypt statement.}
\begin{lstlisting}[language=EasyCrypt,basicstyle=\ttfamily\scriptsize]
op programming_site_output (site : programming_site) : ro_output = site.`2.
\end{lstlisting}
\noindent\textbf{Formal mathematical reading.}
Mathematical definition. This is a total EasyCrypt operation.  The exact displayed statement gives the domain, codomain, and defining equation when present.

\noindent\textbf{Role in the proof.}
Use in this chapter. It is part of the exact formal vocabulary used by subsequent lemmas and games.

\end{formaldefinition}

\begin{formaldefinition}[EasyCrypt line 23]
\noindent\textbf{EasyCrypt name.}
\begin{lstlisting}[language=EasyCrypt,basicstyle=\ttfamily\scriptsize]
programming_site_of_transcript
\end{lstlisting}
\noindent\textbf{Checked EasyCrypt statement.}
\begin{lstlisting}[language=EasyCrypt,basicstyle=\ttfamily\scriptsize]
op programming_site_of_transcript
   (md : mode) (tr : transcript) (y : ro_output) : programming_site =
  (transcript_challenge_query md tr, y).
\end{lstlisting}
\noindent\textbf{Formal mathematical reading.}
Mathematical definition. This is a total EasyCrypt operation.  The exact displayed statement gives the domain, codomain, and defining equation when present.

\noindent\textbf{Role in the proof.}
Use in this chapter. It defines transcript/extraction data for the Fiat-Shamir forking and special-soundness argument.

\end{formaldefinition}

\begin{formaldefinition}[EasyCrypt line 27]
\noindent\textbf{EasyCrypt name.}
\begin{lstlisting}[language=EasyCrypt,basicstyle=\ttfamily\scriptsize]
signature_programming_site_of_transcript
\end{lstlisting}
\noindent\textbf{Checked EasyCrypt statement.}
\begin{lstlisting}[language=EasyCrypt,basicstyle=\ttfamily\scriptsize]
op signature_programming_site_of_transcript
   (md : mode) (tr : transcript) (y : ro_output) : programming_site =
  (transcript_signature_challenge_query md tr, y).
\end{lstlisting}
\noindent\textbf{Formal mathematical reading.}
Mathematical definition. This is a total EasyCrypt operation.  The exact displayed statement gives the domain, codomain, and defining equation when present.

\noindent\textbf{Role in the proof.}
Use in this chapter. It defines transcript/extraction data for the Fiat-Shamir forking and special-soundness argument.

\end{formaldefinition}

\begin{formaldefinition}[EasyCrypt line 31]
\noindent\textbf{EasyCrypt name.}
\begin{lstlisting}[language=EasyCrypt,basicstyle=\ttfamily\scriptsize]
transcript_signature_challenge_output
\end{lstlisting}
\noindent\textbf{Checked EasyCrypt statement.}
\begin{lstlisting}[language=EasyCrypt,basicstyle=\ttfamily\scriptsize]
op transcript_signature_challenge_output (tr : transcript) : ro_output =
  ro_output_of_challenge (sig_challenge (transcript_signature tr)).
\end{lstlisting}
\noindent\textbf{Formal mathematical reading.}
Mathematical definition. This is a total EasyCrypt operation.  The exact displayed statement gives the domain, codomain, and defining equation when present.

\noindent\textbf{Role in the proof.}
Use in this chapter. It defines transcript/extraction data for the Fiat-Shamir forking and special-soundness argument.

\end{formaldefinition}

\begin{formaldefinition}[EasyCrypt line 34]
\noindent\textbf{EasyCrypt name.}
\begin{lstlisting}[language=EasyCrypt,basicstyle=\ttfamily\scriptsize]
actual_signature_programming_site
\end{lstlisting}
\noindent\textbf{Checked EasyCrypt statement.}
\begin{lstlisting}[language=EasyCrypt,basicstyle=\ttfamily\scriptsize]
op actual_signature_programming_site
   (md : mode) (tr : transcript) : programming_site =
  signature_programming_site_of_transcript md tr
    (transcript_signature_challenge_output tr).
\end{lstlisting}
\noindent\textbf{Formal mathematical reading.}
Mathematical definition. This is a total EasyCrypt operation.  The exact displayed statement gives the domain, codomain, and defining equation when present.

\noindent\textbf{Role in the proof.}
Use in this chapter. It is part of the exact formal vocabulary used by subsequent lemmas and games.

\end{formaldefinition}

\begin{formaldefinition}[EasyCrypt line 39]
\noindent\textbf{EasyCrypt name.}
\begin{lstlisting}[language=EasyCrypt,basicstyle=\ttfamily\scriptsize]
honest_signing_programming_site
\end{lstlisting}
\noindent\textbf{Checked EasyCrypt statement.}
\begin{lstlisting}[language=EasyCrypt,basicstyle=\ttfamily\scriptsize]
op honest_signing_programming_site
   (md : mode) (sk : skey) (m : message) (ctx : context)
   (coins : random_coins) : programming_site =
  actual_signature_programming_site md
    (transcript_from_honest_signing md sk m ctx coins).
\end{lstlisting}
\noindent\textbf{Formal mathematical reading.}
Mathematical definition. This is a total EasyCrypt operation.  The exact displayed statement gives the domain, codomain, and defining equation when present.

\noindent\textbf{Role in the proof.}
Use in this chapter. It is part of the exact formal vocabulary used by subsequent lemmas and games.

\end{formaldefinition}

\begin{formaldefinition}[EasyCrypt line 45]
\noindent\textbf{EasyCrypt name.}
\begin{lstlisting}[language=EasyCrypt,basicstyle=\ttfamily\scriptsize]
honest_signing_programming_site_query
\end{lstlisting}
\noindent\textbf{Checked EasyCrypt statement.}
\begin{lstlisting}[language=EasyCrypt,basicstyle=\ttfamily\scriptsize]
op honest_signing_programming_site_query
   (md : mode) (sk : skey) (m : message) (ctx : context)
   (coins : random_coins) : ro_query =
  programming_site_query
    (honest_signing_programming_site md sk m ctx coins).
\end{lstlisting}
\noindent\textbf{Formal mathematical reading.}
Mathematical definition. This is a total EasyCrypt operation.  The exact displayed statement gives the domain, codomain, and defining equation when present.

\noindent\textbf{Role in the proof.}
Use in this chapter. It is part of the exact formal vocabulary used by subsequent lemmas and games.

\end{formaldefinition}

\begin{formaldefinition}[EasyCrypt line 51]
\noindent\textbf{EasyCrypt name.}
\begin{lstlisting}[language=EasyCrypt,basicstyle=\ttfamily\scriptsize]
programmed_site_query_min_entropy_bound_for
\end{lstlisting}
\noindent\textbf{Checked EasyCrypt statement.}
\begin{lstlisting}[language=EasyCrypt,basicstyle=\ttfamily\scriptsize]
op programmed_site_query_min_entropy_bound_for
   (d : random_coins distr)
   (md : mode) (sk : skey) (m : message) (ctx : context)
   (qs : ro_query list) (bd : real) : bool =
  forall q, q \in qs =>
    mu d
      (fun coins =>
        honest_signing_programming_site_query md sk m ctx coins = q) <= bd.
\end{lstlisting}
\noindent\textbf{Formal mathematical reading.}
Mathematical definition. This is a Boolean predicate.  The exact displayed EasyCrypt statement gives its arguments; mathematically it is an event, invariant, support condition, or well-formedness predicate over those arguments.

\noindent\textbf{Role in the proof.}
Use in this chapter. It names a numerical term that appears in a game-hop or final security inequality.

\end{formaldefinition}

\begin{formaldefinition}[EasyCrypt line 60]
\noindent\textbf{EasyCrypt name.}
\begin{lstlisting}[language=EasyCrypt,basicstyle=\ttfamily\scriptsize]
programmed_site_query_min_entropy_bound
\end{lstlisting}
\noindent\textbf{Checked EasyCrypt statement.}
\begin{lstlisting}[language=EasyCrypt,basicstyle=\ttfamily\scriptsize]
op programmed_site_query_min_entropy_bound
   (md : mode) (sk : skey) (m : message) (ctx : context)
   (qs : ro_query list) (bd : real) : bool =
  programmed_site_query_min_entropy_bound_for drandom_coins md sk m ctx qs bd.
\end{lstlisting}
\noindent\textbf{Formal mathematical reading.}
Mathematical definition. This is a Boolean predicate.  The exact displayed EasyCrypt statement gives its arguments; mathematically it is an event, invariant, support condition, or well-formedness predicate over those arguments.

\noindent\textbf{Role in the proof.}
Use in this chapter. It names a numerical term that appears in a game-hop or final security inequality.

\end{formaldefinition}

\begin{formaldefinition}[EasyCrypt line 65]
\noindent\textbf{EasyCrypt name.}
\begin{lstlisting}[language=EasyCrypt,basicstyle=\ttfamily\scriptsize]
programming_site_matches_transcript
\end{lstlisting}
\noindent\textbf{Checked EasyCrypt statement.}
\begin{lstlisting}[language=EasyCrypt,basicstyle=\ttfamily\scriptsize]
op programming_site_matches_transcript
   (md : mode) (tr : transcript) (site : programming_site) : bool =
  programming_site_query site = transcript_challenge_query md tr /\
  transcript_challenge_matches tr (programming_site_output site).
\end{lstlisting}
\noindent\textbf{Formal mathematical reading.}
Mathematical definition. This is a Boolean predicate.  The exact displayed EasyCrypt statement gives its arguments; mathematically it is an event, invariant, support condition, or well-formedness predicate over those arguments.

\noindent\textbf{Role in the proof.}
Use in this chapter. It defines transcript/extraction data for the Fiat-Shamir forking and special-soundness argument.

\end{formaldefinition}

\begin{formaldefinition}[EasyCrypt line 70]
\noindent\textbf{EasyCrypt name.}
\begin{lstlisting}[language=EasyCrypt,basicstyle=\ttfamily\scriptsize]
reprogramming_failure
\end{lstlisting}
\noindent\textbf{Checked EasyCrypt statement.}
\begin{lstlisting}[language=EasyCrypt,basicstyle=\ttfamily\scriptsize]
op reprogramming_failure
   (md : mode) (tr : transcript) (site : programming_site) : bool =
  ! programming_site_matches_transcript md tr site.
\end{lstlisting}
\noindent\textbf{Formal mathematical reading.}
Mathematical definition. This is a Boolean predicate.  The exact displayed EasyCrypt statement gives its arguments; mathematically it is an event, invariant, support condition, or well-formedness predicate over those arguments.

\noindent\textbf{Role in the proof.}
Use in this chapter. It names a bad event or rejection condition whose probability must be zero or explicitly bounded.

\end{formaldefinition}

\begin{formaldefinition}[EasyCrypt line 74]
\noindent\textbf{EasyCrypt name.}
\begin{lstlisting}[language=EasyCrypt,basicstyle=\ttfamily\scriptsize]
challenge_entropy_support_ok
\end{lstlisting}
\noindent\textbf{Checked EasyCrypt statement.}
\begin{lstlisting}[language=EasyCrypt,basicstyle=\ttfamily\scriptsize]
op challenge_entropy_support_ok (md : mode) (tr : transcript) : bool =
  challenge_sparse md (transcript_challenge tr).
\end{lstlisting}
\noindent\textbf{Formal mathematical reading.}
Mathematical definition. This is a Boolean predicate.  The exact displayed EasyCrypt statement gives its arguments; mathematically it is an event, invariant, support condition, or well-formedness predicate over those arguments.

\noindent\textbf{Role in the proof.}
Use in this chapter. It names a well-formedness, support, or validity predicate needed by later preservation lemmas.

\end{formaldefinition}

\begin{formaldefinition}[EasyCrypt line 77]
\noindent\textbf{EasyCrypt name.}
\begin{lstlisting}[language=EasyCrypt,basicstyle=\ttfamily\scriptsize]
min_entropy_failure
\end{lstlisting}
\noindent\textbf{Checked EasyCrypt statement.}
\begin{lstlisting}[language=EasyCrypt,basicstyle=\ttfamily\scriptsize]
op min_entropy_failure (md : mode) (tr : transcript) : bool =
  ! challenge_entropy_support_ok md tr.
\end{lstlisting}
\noindent\textbf{Formal mathematical reading.}
Mathematical definition. This is a Boolean predicate.  The exact displayed EasyCrypt statement gives its arguments; mathematically it is an event, invariant, support condition, or well-formedness predicate over those arguments.

\noindent\textbf{Role in the proof.}
Use in this chapter. It names a bad event or rejection condition whose probability must be zero or explicitly bounded.

\end{formaldefinition}

\begin{formaldefinition}[EasyCrypt line 80]
\noindent\textbf{EasyCrypt name.}
\begin{lstlisting}[language=EasyCrypt,basicstyle=\ttfamily\scriptsize]
transcript_log_min_entropy_clear
\end{lstlisting}
\noindent\textbf{Checked EasyCrypt statement.}
\begin{lstlisting}[language=EasyCrypt,basicstyle=\ttfamily\scriptsize]
op transcript_log_min_entropy_clear (md : mode) (trs : transcript list) : bool =
  all (fun tr => ! min_entropy_failure md tr) trs.
\end{lstlisting}
\noindent\textbf{Formal mathematical reading.}
Mathematical definition. This is a Boolean predicate.  The exact displayed EasyCrypt statement gives its arguments; mathematically it is an event, invariant, support condition, or well-formedness predicate over those arguments.

\noindent\textbf{Role in the proof.}
Use in this chapter. It defines transcript/extraction data for the Fiat-Shamir forking and special-soundness argument.

\end{formaldefinition}

\begin{formaldefinition}[EasyCrypt line 83]
\noindent\textbf{EasyCrypt name.}
\begin{lstlisting}[language=EasyCrypt,basicstyle=\ttfamily\scriptsize]
signing_record_log_min_entropy_clear
\end{lstlisting}
\noindent\textbf{Checked EasyCrypt statement.}
\begin{lstlisting}[language=EasyCrypt,basicstyle=\ttfamily\scriptsize]
op signing_record_log_min_entropy_clear
   (md : mode) (rs : signing_transcript_record list) : bool =
  all (fun r => ! min_entropy_failure md (signing_record_transcript r)) rs.
\end{lstlisting}
\noindent\textbf{Formal mathematical reading.}
Mathematical definition. This is a Boolean predicate.  The exact displayed EasyCrypt statement gives its arguments; mathematically it is an event, invariant, support condition, or well-formedness predicate over those arguments.

\noindent\textbf{Role in the proof.}
Use in this chapter. It is part of the exact formal vocabulary used by subsequent lemmas and games.

\end{formaldefinition}

\begin{formaldefinition}[EasyCrypt line 87]
\noindent\textbf{EasyCrypt name.}
\begin{lstlisting}[language=EasyCrypt,basicstyle=\ttfamily\scriptsize]
fs_with_aborts_failure
\end{lstlisting}
\noindent\textbf{Checked EasyCrypt statement.}
\begin{lstlisting}[language=EasyCrypt,basicstyle=\ttfamily\scriptsize]
op fs_with_aborts_failure
   (md : mode) (tr : transcript) (site : programming_site) : bool =
  reprogramming_failure md tr site \/ min_entropy_failure md tr.
\end{lstlisting}
\noindent\textbf{Formal mathematical reading.}
Mathematical definition. This is a Boolean predicate.  The exact displayed EasyCrypt statement gives its arguments; mathematically it is an event, invariant, support condition, or well-formedness predicate over those arguments.

\noindent\textbf{Role in the proof.}
Use in this chapter. It names a bad event or rejection condition whose probability must be zero or explicitly bounded.

\end{formaldefinition}

\begin{formaldefinition}[EasyCrypt line 91]
\noindent\textbf{EasyCrypt name.}
\begin{lstlisting}[language=EasyCrypt,basicstyle=\ttfamily\scriptsize]
transcript_signature_programming_site_clear
\end{lstlisting}
\noindent\textbf{Checked EasyCrypt statement.}
\begin{lstlisting}[language=EasyCrypt,basicstyle=\ttfamily\scriptsize]
op transcript_signature_programming_site_clear
   (md : mode) (tr : transcript) : bool =
  ! fs_with_aborts_failure md tr (actual_signature_programming_site md tr).
\end{lstlisting}
\noindent\textbf{Formal mathematical reading.}
Mathematical definition. This is a Boolean predicate.  The exact displayed EasyCrypt statement gives its arguments; mathematically it is an event, invariant, support condition, or well-formedness predicate over those arguments.

\noindent\textbf{Role in the proof.}
Use in this chapter. It defines transcript/extraction data for the Fiat-Shamir forking and special-soundness argument.

\end{formaldefinition}

\begin{formaldefinition}[EasyCrypt line 95]
\noindent\textbf{EasyCrypt name.}
\begin{lstlisting}[language=EasyCrypt,basicstyle=\ttfamily\scriptsize]
transcript_log_signature_programming_sites_clear
\end{lstlisting}
\noindent\textbf{Checked EasyCrypt statement.}
\begin{lstlisting}[language=EasyCrypt,basicstyle=\ttfamily\scriptsize]
op transcript_log_signature_programming_sites_clear
   (md : mode) (trs : transcript list) : bool =
  all (transcript_signature_programming_site_clear md) trs.
\end{lstlisting}
\noindent\textbf{Formal mathematical reading.}
Mathematical definition. This is a Boolean predicate.  The exact displayed EasyCrypt statement gives its arguments; mathematically it is an event, invariant, support condition, or well-formedness predicate over those arguments.

\noindent\textbf{Role in the proof.}
Use in this chapter. It defines transcript/extraction data for the Fiat-Shamir forking and special-soundness argument.

\end{formaldefinition}

\begin{formaldefinition}[EasyCrypt line 99]
\noindent\textbf{EasyCrypt name.}
\begin{lstlisting}[language=EasyCrypt,basicstyle=\ttfamily\scriptsize]
fs_with_aborts_bound_obligation
\end{lstlisting}
\noindent\textbf{Checked EasyCrypt statement.}
\begin{lstlisting}[language=EasyCrypt,basicstyle=\ttfamily\scriptsize]
op fs_with_aborts_bound_obligation : bool =
  0%r <= fs_with_aborts_reprogramming_term /\
  0%r <= fs_with_aborts_min_entropy_term.
\end{lstlisting}
\noindent\textbf{Formal mathematical reading.}
Mathematical definition. This is a Boolean predicate.  The exact displayed EasyCrypt statement gives its arguments; mathematically it is an event, invariant, support condition, or well-formedness predicate over those arguments.

\noindent\textbf{Role in the proof.}
Use in this chapter. It names a numerical term that appears in a game-hop or final security inequality.

\end{formaldefinition}

\begin{formaldefinition}[EasyCrypt line 103]
\noindent\textbf{EasyCrypt name.}
\begin{lstlisting}[language=EasyCrypt,basicstyle=\ttfamily\scriptsize]
counted_rom_event
\end{lstlisting}
\noindent\textbf{Checked EasyCrypt statement.}
\begin{lstlisting}[language=EasyCrypt,basicstyle=\ttfamily\scriptsize]
type counted_rom_event =
  [ CountedHashQuery of ro_query
  | CountedSignatureSite of programming_site
  | CountedProgrammedSite of programming_site
  | CountedMinEntropyCheck of transcript ].
\end{lstlisting}
\noindent\textbf{Formal mathematical reading.}
Mathematical definition. This is a definitional type abbreviation.  The exact displayed EasyCrypt statement gives the right-hand side; later proof steps may unfold it without adding an assumption.

\noindent\textbf{Role in the proof.}
Use in this chapter. It is part of the exact formal vocabulary used by subsequent lemmas and games.

\end{formaldefinition}

\begin{formaldefinition}[EasyCrypt line 109]
\noindent\textbf{EasyCrypt name.}
\begin{lstlisting}[language=EasyCrypt,basicstyle=\ttfamily\scriptsize]
counted_event_is_hash_query
\end{lstlisting}
\noindent\textbf{Checked EasyCrypt statement.}
\begin{lstlisting}[language=EasyCrypt,basicstyle=\ttfamily\scriptsize]
op counted_event_is_hash_query (ev : counted_rom_event) : bool =
  with ev = CountedHashQuery _ => true
  with ev = CountedSignatureSite _ => false
  with ev = CountedProgrammedSite _ => false
  with ev = CountedMinEntropyCheck _ => false.
\end{lstlisting}
\noindent\textbf{Formal mathematical reading.}
Mathematical definition. This is a Boolean predicate.  The exact displayed EasyCrypt statement gives its arguments; mathematically it is an event, invariant, support condition, or well-formedness predicate over those arguments.

\noindent\textbf{Role in the proof.}
Use in this chapter. It is part of the exact formal vocabulary used by subsequent lemmas and games.

\end{formaldefinition}

\begin{formaldefinition}[EasyCrypt line 115]
\noindent\textbf{EasyCrypt name.}
\begin{lstlisting}[language=EasyCrypt,basicstyle=\ttfamily\scriptsize]
counted_event_is_signature_site
\end{lstlisting}
\noindent\textbf{Checked EasyCrypt statement.}
\begin{lstlisting}[language=EasyCrypt,basicstyle=\ttfamily\scriptsize]
op counted_event_is_signature_site (ev : counted_rom_event) : bool =
  with ev = CountedHashQuery _ => false
  with ev = CountedSignatureSite _ => true
  with ev = CountedProgrammedSite _ => false
  with ev = CountedMinEntropyCheck _ => false.
\end{lstlisting}
\noindent\textbf{Formal mathematical reading.}
Mathematical definition. This is a Boolean predicate.  The exact displayed EasyCrypt statement gives its arguments; mathematically it is an event, invariant, support condition, or well-formedness predicate over those arguments.

\noindent\textbf{Role in the proof.}
Use in this chapter. It is part of the exact formal vocabulary used by subsequent lemmas and games.

\end{formaldefinition}

\begin{formaldefinition}[EasyCrypt line 121]
\noindent\textbf{EasyCrypt name.}
\begin{lstlisting}[language=EasyCrypt,basicstyle=\ttfamily\scriptsize]
counted_event_is_programmed_site
\end{lstlisting}
\noindent\textbf{Checked EasyCrypt statement.}
\begin{lstlisting}[language=EasyCrypt,basicstyle=\ttfamily\scriptsize]
op counted_event_is_programmed_site (ev : counted_rom_event) : bool =
  with ev = CountedHashQuery _ => false
  with ev = CountedSignatureSite _ => false
  with ev = CountedProgrammedSite _ => true
  with ev = CountedMinEntropyCheck _ => false.
\end{lstlisting}
\noindent\textbf{Formal mathematical reading.}
Mathematical definition. This is a Boolean predicate.  The exact displayed EasyCrypt statement gives its arguments; mathematically it is an event, invariant, support condition, or well-formedness predicate over those arguments.

\noindent\textbf{Role in the proof.}
Use in this chapter. It is part of the exact formal vocabulary used by subsequent lemmas and games.

\end{formaldefinition}

\begin{formaldefinition}[EasyCrypt line 127]
\noindent\textbf{EasyCrypt name.}
\begin{lstlisting}[language=EasyCrypt,basicstyle=\ttfamily\scriptsize]
counted_event_is_min_entropy_check
\end{lstlisting}
\noindent\textbf{Checked EasyCrypt statement.}
\begin{lstlisting}[language=EasyCrypt,basicstyle=\ttfamily\scriptsize]
op counted_event_is_min_entropy_check (ev : counted_rom_event) : bool =
  with ev = CountedHashQuery _ => false
  with ev = CountedSignatureSite _ => false
  with ev = CountedProgrammedSite _ => false
  with ev = CountedMinEntropyCheck _ => true.
\end{lstlisting}
\noindent\textbf{Formal mathematical reading.}
Mathematical definition. This is a Boolean predicate.  The exact displayed EasyCrypt statement gives its arguments; mathematically it is an event, invariant, support condition, or well-formedness predicate over those arguments.

\noindent\textbf{Role in the proof.}
Use in this chapter. It is part of the exact formal vocabulary used by subsequent lemmas and games.

\end{formaldefinition}

\begin{formaldefinition}[EasyCrypt line 133]
\noindent\textbf{EasyCrypt name.}
\begin{lstlisting}[language=EasyCrypt,basicstyle=\ttfamily\scriptsize]
ro_query_is_challenge
\end{lstlisting}
\noindent\textbf{Checked EasyCrypt statement.}
\begin{lstlisting}[language=EasyCrypt,basicstyle=\ttfamily\scriptsize]
op ro_query_is_challenge (q : ro_query) : bool =
  with q = MessageHashQuery _ => false
  with q = ChallengeHashQuery _ => true
  with q = MatrixExpandQuery _ => false
  with q = SamplerExpandQuery _ => false.
\end{lstlisting}
\noindent\textbf{Formal mathematical reading.}
Mathematical definition. This is a Boolean predicate.  The exact displayed EasyCrypt statement gives its arguments; mathematically it is an event, invariant, support condition, or well-formedness predicate over those arguments.

\noindent\textbf{Role in the proof.}
Use in this chapter. It is part of the exact formal vocabulary used by subsequent lemmas and games.

\end{formaldefinition}

\begin{formaldefinition}[EasyCrypt line 139]
\noindent\textbf{EasyCrypt name.}
\begin{lstlisting}[language=EasyCrypt,basicstyle=\ttfamily\scriptsize]
challenge_query_log_count
\end{lstlisting}
\noindent\textbf{Checked EasyCrypt statement.}
\begin{lstlisting}[language=EasyCrypt,basicstyle=\ttfamily\scriptsize]
op challenge_query_log_count (qs : ro_query list) : int =
  card (oflist (filter ro_query_is_challenge qs)).
\end{lstlisting}
\noindent\textbf{Formal mathematical reading.}
Mathematical definition. This is a total EasyCrypt operation.  The exact displayed statement gives the domain, codomain, and defining equation when present.

\noindent\textbf{Role in the proof.}
Use in this chapter. It is part of the exact formal vocabulary used by subsequent lemmas and games.

\end{formaldefinition}

\begin{formaldefinition}[EasyCrypt line 193]
\noindent\textbf{EasyCrypt name.}
\begin{lstlisting}[language=EasyCrypt,basicstyle=\ttfamily\scriptsize]
counted_trace_hash_queries
\end{lstlisting}
\noindent\textbf{Checked EasyCrypt statement.}
\begin{lstlisting}[language=EasyCrypt,basicstyle=\ttfamily\scriptsize]
op counted_trace_hash_queries (evs : counted_rom_event list) : int =
  size (filter counted_event_is_hash_query evs).
\end{lstlisting}
\noindent\textbf{Formal mathematical reading.}
Mathematical definition. This is a total EasyCrypt operation.  The exact displayed statement gives the domain, codomain, and defining equation when present.

\noindent\textbf{Role in the proof.}
Use in this chapter. It is part of the exact formal vocabulary used by subsequent lemmas and games.

\end{formaldefinition}

\begin{formaldefinition}[EasyCrypt line 196]
\noindent\textbf{EasyCrypt name.}
\begin{lstlisting}[language=EasyCrypt,basicstyle=\ttfamily\scriptsize]
counted_trace_signature_sites
\end{lstlisting}
\noindent\textbf{Checked EasyCrypt statement.}
\begin{lstlisting}[language=EasyCrypt,basicstyle=\ttfamily\scriptsize]
op counted_trace_signature_sites (evs : counted_rom_event list) : int =
  size (filter counted_event_is_signature_site evs).
\end{lstlisting}
\noindent\textbf{Formal mathematical reading.}
Mathematical definition. This is a total EasyCrypt operation.  The exact displayed statement gives the domain, codomain, and defining equation when present.

\noindent\textbf{Role in the proof.}
Use in this chapter. It is part of the exact formal vocabulary used by subsequent lemmas and games.

\end{formaldefinition}

\begin{formaldefinition}[EasyCrypt line 199]
\noindent\textbf{EasyCrypt name.}
\begin{lstlisting}[language=EasyCrypt,basicstyle=\ttfamily\scriptsize]
counted_trace_programmed_sites
\end{lstlisting}
\noindent\textbf{Checked EasyCrypt statement.}
\begin{lstlisting}[language=EasyCrypt,basicstyle=\ttfamily\scriptsize]
op counted_trace_programmed_sites (evs : counted_rom_event list) : int =
  size (filter counted_event_is_programmed_site evs).
\end{lstlisting}
\noindent\textbf{Formal mathematical reading.}
Mathematical definition. This is a total EasyCrypt operation.  The exact displayed statement gives the domain, codomain, and defining equation when present.

\noindent\textbf{Role in the proof.}
Use in this chapter. It is part of the exact formal vocabulary used by subsequent lemmas and games.

\end{formaldefinition}

\begin{formaldefinition}[EasyCrypt line 202]
\noindent\textbf{EasyCrypt name.}
\begin{lstlisting}[language=EasyCrypt,basicstyle=\ttfamily\scriptsize]
counted_trace_min_entropy_checks
\end{lstlisting}
\noindent\textbf{Checked EasyCrypt statement.}
\begin{lstlisting}[language=EasyCrypt,basicstyle=\ttfamily\scriptsize]
op counted_trace_min_entropy_checks (evs : counted_rom_event list) : int =
  size (filter counted_event_is_min_entropy_check evs).
\end{lstlisting}
\noindent\textbf{Formal mathematical reading.}
Mathematical definition. This is a total EasyCrypt operation.  The exact displayed statement gives the domain, codomain, and defining equation when present.

\noindent\textbf{Role in the proof.}
Use in this chapter. It is part of the exact formal vocabulary used by subsequent lemmas and games.

\end{formaldefinition}

\begin{formaldefinition}[EasyCrypt line 205]
\noindent\textbf{EasyCrypt name.}
\begin{lstlisting}[language=EasyCrypt,basicstyle=\ttfamily\scriptsize]
programming_sites_same_query
\end{lstlisting}
\noindent\textbf{Checked EasyCrypt statement.}
\begin{lstlisting}[language=EasyCrypt,basicstyle=\ttfamily\scriptsize]
op programming_sites_same_query
   (left right : programming_site) : bool =
  programming_site_query left = programming_site_query right.
\end{lstlisting}
\noindent\textbf{Formal mathematical reading.}
Mathematical definition. This is a Boolean predicate.  The exact displayed EasyCrypt statement gives its arguments; mathematically it is an event, invariant, support condition, or well-formedness predicate over those arguments.

\noindent\textbf{Role in the proof.}
Use in this chapter. It is part of the exact formal vocabulary used by subsequent lemmas and games.

\end{formaldefinition}

\begin{formaldefinition}[EasyCrypt line 209]
\noindent\textbf{EasyCrypt name.}
\begin{lstlisting}[language=EasyCrypt,basicstyle=\ttfamily\scriptsize]
programming_sites_conflict
\end{lstlisting}
\noindent\textbf{Checked EasyCrypt statement.}
\begin{lstlisting}[language=EasyCrypt,basicstyle=\ttfamily\scriptsize]
op programming_sites_conflict
   (left right : programming_site) : bool =
  programming_sites_same_query left right /\
  programming_site_output left <> programming_site_output right.
\end{lstlisting}
\noindent\textbf{Formal mathematical reading.}
Mathematical definition. This is a Boolean predicate.  The exact displayed EasyCrypt statement gives its arguments; mathematically it is an event, invariant, support condition, or well-formedness predicate over those arguments.

\noindent\textbf{Role in the proof.}
Use in this chapter. It is part of the exact formal vocabulary used by subsequent lemmas and games.

\end{formaldefinition}

\begin{formaldefinition}[EasyCrypt line 214]
\noindent\textbf{EasyCrypt name.}
\begin{lstlisting}[language=EasyCrypt,basicstyle=\ttfamily\scriptsize]
programming_site_prequeried
\end{lstlisting}
\noindent\textbf{Checked EasyCrypt statement.}
\begin{lstlisting}[language=EasyCrypt,basicstyle=\ttfamily\scriptsize]
op programming_site_prequeried
   (queries : ro_query list) (site : programming_site) : bool =
  programming_site_query site \in queries.
\end{lstlisting}
\noindent\textbf{Formal mathematical reading.}
Mathematical definition. This is a Boolean predicate.  The exact displayed EasyCrypt statement gives its arguments; mathematically it is an event, invariant, support condition, or well-formedness predicate over those arguments.

\noindent\textbf{Role in the proof.}
Use in this chapter. It is part of the exact formal vocabulary used by subsequent lemmas and games.

\end{formaldefinition}

\begin{formaldefinition}[EasyCrypt line 218]
\noindent\textbf{EasyCrypt name.}
\begin{lstlisting}[language=EasyCrypt,basicstyle=\ttfamily\scriptsize]
programming_site_reprograms
\end{lstlisting}
\noindent\textbf{Checked EasyCrypt statement.}
\begin{lstlisting}[language=EasyCrypt,basicstyle=\ttfamily\scriptsize]
op programming_site_reprograms
   (programmed : programming_site list) (site : programming_site) : bool =
  has (fun old => programming_sites_conflict site old) programmed.
\end{lstlisting}
\noindent\textbf{Formal mathematical reading.}
Mathematical definition. This is a Boolean predicate.  The exact displayed EasyCrypt statement gives its arguments; mathematically it is an event, invariant, support condition, or well-formedness predicate over those arguments.

\noindent\textbf{Role in the proof.}
Use in this chapter. It is part of the exact formal vocabulary used by subsequent lemmas and games.

\end{formaldefinition}

\begin{formaldefinition}[EasyCrypt line 222]
\noindent\textbf{EasyCrypt name.}
\begin{lstlisting}[language=EasyCrypt,basicstyle=\ttfamily\scriptsize]
programming_site_fresh
\end{lstlisting}
\noindent\textbf{Checked EasyCrypt statement.}
\begin{lstlisting}[language=EasyCrypt,basicstyle=\ttfamily\scriptsize]
op programming_site_fresh
   (queries : ro_query list) (programmed : programming_site list)
   (site : programming_site) : bool =
  ! programming_site_prequeried queries site /\
  ! programming_site_reprograms programmed site.
\end{lstlisting}
\noindent\textbf{Formal mathematical reading.}
Mathematical definition. This is a Boolean predicate.  The exact displayed EasyCrypt statement gives its arguments; mathematically it is an event, invariant, support condition, or well-formedness predicate over those arguments.

\noindent\textbf{Role in the proof.}
Use in this chapter. It is part of the exact formal vocabulary used by subsequent lemmas and games.

\end{formaldefinition}

\begin{formaldefinition}[EasyCrypt line 228]
\noindent\textbf{EasyCrypt name.}
\begin{lstlisting}[language=EasyCrypt,basicstyle=\ttfamily\scriptsize]
programming_transcript_fresh
\end{lstlisting}
\noindent\textbf{Checked EasyCrypt statement.}
\begin{lstlisting}[language=EasyCrypt,basicstyle=\ttfamily\scriptsize]
op programming_transcript_fresh
   (md : mode) (queries : ro_query list)
   (programmed : programming_site list) (tr : transcript) : bool =
  programming_site_fresh queries programmed
    (actual_signature_programming_site md tr) /\
  ! min_entropy_failure md tr.
\end{lstlisting}
\noindent\textbf{Formal mathematical reading.}
Mathematical definition. This is a Boolean predicate.  The exact displayed EasyCrypt statement gives its arguments; mathematically it is an event, invariant, support condition, or well-formedness predicate over those arguments.

\noindent\textbf{Role in the proof.}
Use in this chapter. It defines transcript/extraction data for the Fiat-Shamir forking and special-soundness argument.

\end{formaldefinition}

\begin{formaldefinition}[EasyCrypt line 235]
\noindent\textbf{EasyCrypt name.}
\begin{lstlisting}[language=EasyCrypt,basicstyle=\ttfamily\scriptsize]
rom_counted_programming_bad
\end{lstlisting}
\noindent\textbf{Checked EasyCrypt statement.}
\begin{lstlisting}[language=EasyCrypt,basicstyle=\ttfamily\scriptsize]
op rom_counted_programming_bad
   (queries : ro_query list) (programmed : programming_site list)
   (entropy_bad : bool) : bool =
  entropy_bad \/
  has (fun site => programming_site_prequeried queries site) programmed \/
  has (fun site => programming_site_reprograms programmed site) programmed.
\end{lstlisting}
\noindent\textbf{Formal mathematical reading.}
Mathematical definition. This is a Boolean predicate.  The exact displayed EasyCrypt statement gives its arguments; mathematically it is an event, invariant, support condition, or well-formedness predicate over those arguments.

\noindent\textbf{Role in the proof.}
Use in this chapter. It names a bad event or rejection condition whose probability must be zero or explicitly bounded.

\end{formaldefinition}

\begin{formaldefinition}[EasyCrypt line 242]
\noindent\textbf{EasyCrypt name.}
\begin{lstlisting}[language=EasyCrypt,basicstyle=\ttfamily\scriptsize]
rom_counted_state_bad
\end{lstlisting}
\noindent\textbf{Checked EasyCrypt statement.}
\begin{lstlisting}[language=EasyCrypt,basicstyle=\ttfamily\scriptsize]
op rom_counted_state_bad
   (queries : ro_query list) (programmed : programming_site list)
   (bad_min_entropy : bool) : bool =
  rom_counted_programming_bad queries programmed bad_min_entropy.
\end{lstlisting}
\noindent\textbf{Formal mathematical reading.}
Mathematical definition. This is a Boolean predicate.  The exact displayed EasyCrypt statement gives its arguments; mathematically it is an event, invariant, support condition, or well-formedness predicate over those arguments.

\noindent\textbf{Role in the proof.}
Use in this chapter. It names a bad event or rejection condition whose probability must be zero or explicitly bounded.

\end{formaldefinition}

\begin{formaldefinition}[EasyCrypt line 247]
\noindent\textbf{EasyCrypt name.}
\begin{lstlisting}[language=EasyCrypt,basicstyle=\ttfamily\scriptsize]
counted_rom_programming_bound_obligation
\end{lstlisting}
\noindent\textbf{Checked EasyCrypt statement.}
\begin{lstlisting}[language=EasyCrypt,basicstyle=\ttfamily\scriptsize]
op counted_rom_programming_bound_obligation : bool =
  0%r <= counted_rom_programming_loss_term.
\end{lstlisting}
\noindent\textbf{Formal mathematical reading.}
Mathematical definition. This is a Boolean predicate.  The exact displayed EasyCrypt statement gives its arguments; mathematically it is an event, invariant, support condition, or well-formedness predicate over those arguments.

\noindent\textbf{Role in the proof.}
Use in this chapter. It names a numerical term that appears in a game-hop or final security inequality.

\end{formaldefinition}

\begin{formaldefinition}[EasyCrypt line 250]
\noindent\textbf{EasyCrypt name.}
\begin{lstlisting}[language=EasyCrypt,basicstyle=\ttfamily\scriptsize]
CountedLazyROM
\end{lstlisting}
\noindent\textbf{Checked EasyCrypt statement.}
\begin{lstlisting}[language=EasyCrypt,basicstyle=\ttfamily\scriptsize]
module CountedLazyROM(O : Oracle) = {
\end{lstlisting}
\noindent\textbf{Formal mathematical reading.}
Mathematical definition. This is a probabilistic program/game or a module adapter.  Its procedures induce the probability space used by the game-based proof.

\noindent\textbf{Role in the proof.}
Use in this chapter. It defines the oracle boundary through which adversaries and simulations interact.

\end{formaldefinition}

\begin{formaldefinition}[EasyCrypt line 258]
\noindent\textbf{EasyCrypt name.}
\begin{lstlisting}[language=EasyCrypt,basicstyle=\ttfamily\scriptsize]
init
\end{lstlisting}
\noindent\textbf{Checked EasyCrypt statement.}
\begin{lstlisting}[language=EasyCrypt,basicstyle=\ttfamily\scriptsize]
proc init() : unit = {
\end{lstlisting}
\noindent\textbf{Formal mathematical reading.}
Mathematical definition. This procedure is a command in the probabilistic program.  Its return value, state updates, and oracle calls are the operational content of the game.

\noindent\textbf{Role in the proof.}
Use in this chapter. It supplies an executable component of the formal probabilistic model.

\end{formaldefinition}

\begin{formaldefinition}[EasyCrypt line 268]
\noindent\textbf{EasyCrypt name.}
\begin{lstlisting}[language=EasyCrypt,basicstyle=\ttfamily\scriptsize]
get
\end{lstlisting}
\noindent\textbf{Checked EasyCrypt statement.}
\begin{lstlisting}[language=EasyCrypt,basicstyle=\ttfamily\scriptsize]
proc get(q : ro_query) : ro_output = {
\end{lstlisting}
\noindent\textbf{Formal mathematical reading.}
Mathematical definition. This procedure is a command in the probabilistic program.  Its return value, state updates, and oracle calls are the operational content of the game.

\noindent\textbf{Role in the proof.}
Use in this chapter. It supplies an executable component of the formal probabilistic model.

\end{formaldefinition}

\begin{formaldefinition}[EasyCrypt line 276]
\noindent\textbf{EasyCrypt name.}
\begin{lstlisting}[language=EasyCrypt,basicstyle=\ttfamily\scriptsize]
program
\end{lstlisting}
\noindent\textbf{Checked EasyCrypt statement.}
\begin{lstlisting}[language=EasyCrypt,basicstyle=\ttfamily\scriptsize]
proc program(site : programming_site) : unit = {
\end{lstlisting}
\noindent\textbf{Formal mathematical reading.}
Mathematical definition. This procedure is a command in the probabilistic program.  Its return value, state updates, and oracle calls are the operational content of the game.

\noindent\textbf{Role in the proof.}
Use in this chapter. It supplies an executable component of the formal probabilistic model.

\end{formaldefinition}

\begin{formaldefinition}[EasyCrypt line 285]
\noindent\textbf{EasyCrypt name.}
\begin{lstlisting}[language=EasyCrypt,basicstyle=\ttfamily\scriptsize]
observe_signature_site
\end{lstlisting}
\noindent\textbf{Checked EasyCrypt statement.}
\begin{lstlisting}[language=EasyCrypt,basicstyle=\ttfamily\scriptsize]
proc observe_signature_site(site : programming_site) : unit = {
\end{lstlisting}
\noindent\textbf{Formal mathematical reading.}
Mathematical definition. This procedure is a command in the probabilistic program.  Its return value, state updates, and oracle calls are the operational content of the game.

\noindent\textbf{Role in the proof.}
Use in this chapter. It supplies an executable component of the formal probabilistic model.

\end{formaldefinition}

\begin{formaldefinition}[EasyCrypt line 289]
\noindent\textbf{EasyCrypt name.}
\begin{lstlisting}[language=EasyCrypt,basicstyle=\ttfamily\scriptsize]
observe_transcript
\end{lstlisting}
\noindent\textbf{Checked EasyCrypt statement.}
\begin{lstlisting}[language=EasyCrypt,basicstyle=\ttfamily\scriptsize]
proc observe_transcript(md : mode, tr : transcript) : unit = {
\end{lstlisting}
\noindent\textbf{Formal mathematical reading.}
Mathematical definition. This procedure is a command in the probabilistic program.  Its return value, state updates, and oracle calls are the operational content of the game.

\noindent\textbf{Role in the proof.}
Use in this chapter. It supplies an executable component of the formal probabilistic model.

\end{formaldefinition}

\begin{formaldefinition}[EasyCrypt line 297]
\noindent\textbf{EasyCrypt name.}
\begin{lstlisting}[language=EasyCrypt,basicstyle=\ttfamily\scriptsize]
bad
\end{lstlisting}
\noindent\textbf{Checked EasyCrypt statement.}
\begin{lstlisting}[language=EasyCrypt,basicstyle=\ttfamily\scriptsize]
proc bad() : bool = {
\end{lstlisting}
\noindent\textbf{Formal mathematical reading.}
Mathematical definition. This procedure is a command in the probabilistic program.  Its return value, state updates, and oracle calls are the operational content of the game.

\noindent\textbf{Role in the proof.}
Use in this chapter. It supplies an executable component of the formal probabilistic model.

\end{formaldefinition}

\begin{formaldefinition}[EasyCrypt line 302]
\noindent\textbf{EasyCrypt name.}
\begin{lstlisting}[language=EasyCrypt,basicstyle=\ttfamily\scriptsize]
SigningCoinSampler
\end{lstlisting}
\noindent\textbf{Checked EasyCrypt statement.}
\begin{lstlisting}[language=EasyCrypt,basicstyle=\ttfamily\scriptsize]
module type SigningCoinSampler = {
\end{lstlisting}
\noindent\textbf{Formal mathematical reading.}
Mathematical definition. This is an interface for an oracle, adversary, scheme, sampler, solver, or reduction.  A later theorem quantified over this interface is universally quantified over all modules implementing these procedures.

\noindent\textbf{Role in the proof.}
Use in this chapter. It defines the sampling procedure whose distributional behavior is justified by the later loss lemmas.

\end{formaldefinition}

\begin{formaldefinition}[EasyCrypt line 303]
\noindent\textbf{EasyCrypt name.}
\begin{lstlisting}[language=EasyCrypt,basicstyle=\ttfamily\scriptsize]
sample
\end{lstlisting}
\noindent\textbf{Checked EasyCrypt statement.}
\begin{lstlisting}[language=EasyCrypt,basicstyle=\ttfamily\scriptsize]
proc sample() : random_coins
}.
\end{lstlisting}
\noindent\textbf{Formal mathematical reading.}
Mathematical definition. This procedure is a command in the probabilistic program.  Its return value, state updates, and oracle calls are the operational content of the game.

\noindent\textbf{Role in the proof.}
Use in this chapter. It defines the sampling procedure whose distributional behavior is justified by the later loss lemmas.

\end{formaldefinition}

\begin{formaldefinition}[EasyCrypt line 306]
\noindent\textbf{EasyCrypt name.}
\begin{lstlisting}[language=EasyCrypt,basicstyle=\ttfamily\scriptsize]
DirectSigningCoinSampler
\end{lstlisting}
\noindent\textbf{Checked EasyCrypt statement.}
\begin{lstlisting}[language=EasyCrypt,basicstyle=\ttfamily\scriptsize]
module DirectSigningCoinSampler : SigningCoinSampler = {
\end{lstlisting}
\noindent\textbf{Formal mathematical reading.}
Mathematical definition. This is a probabilistic program/game or a module adapter.  Its procedures induce the probability space used by the game-based proof.

\noindent\textbf{Role in the proof.}
Use in this chapter. It defines the sampling procedure whose distributional behavior is justified by the later loss lemmas.

\end{formaldefinition}

\begin{formaldefinition}[EasyCrypt line 307]
\noindent\textbf{EasyCrypt name.}
\begin{lstlisting}[language=EasyCrypt,basicstyle=\ttfamily\scriptsize]
sample
\end{lstlisting}
\noindent\textbf{Checked EasyCrypt statement.}
\begin{lstlisting}[language=EasyCrypt,basicstyle=\ttfamily\scriptsize]
proc sample() : random_coins = {
\end{lstlisting}
\noindent\textbf{Formal mathematical reading.}
Mathematical definition. This procedure is a command in the probabilistic program.  Its return value, state updates, and oracle calls are the operational content of the game.

\noindent\textbf{Role in the proof.}
Use in this chapter. It defines the sampling procedure whose distributional behavior is justified by the later loss lemmas.

\end{formaldefinition}

\begin{formaldefinition}[EasyCrypt line 501]
\noindent\textbf{EasyCrypt name.}
\begin{lstlisting}[language=EasyCrypt,basicstyle=\ttfamily\scriptsize]
CountedROMInterface
\end{lstlisting}
\noindent\textbf{Checked EasyCrypt statement.}
\begin{lstlisting}[language=EasyCrypt,basicstyle=\ttfamily\scriptsize]
module type CountedROMInterface = {
\end{lstlisting}
\noindent\textbf{Formal mathematical reading.}
Mathematical definition. This is an interface for an oracle, adversary, scheme, sampler, solver, or reduction.  A later theorem quantified over this interface is universally quantified over all modules implementing these procedures.

\noindent\textbf{Role in the proof.}
Use in this chapter. It defines the oracle boundary through which adversaries and simulations interact.

\end{formaldefinition}

\begin{formaldefinition}[EasyCrypt line 502]
\noindent\textbf{EasyCrypt name.}
\begin{lstlisting}[language=EasyCrypt,basicstyle=\ttfamily\scriptsize]
get
\end{lstlisting}
\noindent\textbf{Checked EasyCrypt statement.}
\begin{lstlisting}[language=EasyCrypt,basicstyle=\ttfamily\scriptsize]
proc get(q : ro_query) : ro_output
  proc program(site : programming_site) : unit
  proc observe_signature_site(site : programming_site) : unit
  proc observe_transcript(md : mode, tr : transcript) : unit
}.
\end{lstlisting}
\noindent\textbf{Formal mathematical reading.}
Mathematical definition. This procedure is a command in the probabilistic program.  Its return value, state updates, and oracle calls are the operational content of the game.

\noindent\textbf{Role in the proof.}
Use in this chapter. It supplies an executable component of the formal probabilistic model.

\end{formaldefinition}

\begin{formaldefinition}[EasyCrypt line 508]
\noindent\textbf{EasyCrypt name.}
\begin{lstlisting}[language=EasyCrypt,basicstyle=\ttfamily\scriptsize]
CountedROMClient
\end{lstlisting}
\noindent\textbf{Checked EasyCrypt statement.}
\begin{lstlisting}[language=EasyCrypt,basicstyle=\ttfamily\scriptsize]
module type CountedROMClient(C : CountedROMInterface) = {
\end{lstlisting}
\noindent\textbf{Formal mathematical reading.}
Mathematical definition. This is an interface for an oracle, adversary, scheme, sampler, solver, or reduction.  A later theorem quantified over this interface is universally quantified over all modules implementing these procedures.

\noindent\textbf{Role in the proof.}
Use in this chapter. It defines the oracle boundary through which adversaries and simulations interact.

\end{formaldefinition}

\begin{formaldefinition}[EasyCrypt line 509]
\noindent\textbf{EasyCrypt name.}
\begin{lstlisting}[language=EasyCrypt,basicstyle=\ttfamily\scriptsize]
run
\end{lstlisting}
\noindent\textbf{Checked EasyCrypt statement.}
\begin{lstlisting}[language=EasyCrypt,basicstyle=\ttfamily\scriptsize]
proc run() : unit {
\end{lstlisting}
\noindent\textbf{Formal mathematical reading.}
Mathematical definition. This procedure is a command in the probabilistic program.  Its return value, state updates, and oracle calls are the operational content of the game.

\noindent\textbf{Role in the proof.}
Use in this chapter. It supplies an executable component of the formal probabilistic model.

\end{formaldefinition}

\begin{formaldefinition}[EasyCrypt line 517]
\noindent\textbf{EasyCrypt name.}
\begin{lstlisting}[language=EasyCrypt,basicstyle=\ttfamily\scriptsize]
CountedLazyROMBadGame
\end{lstlisting}
\noindent\textbf{Checked EasyCrypt statement.}
\begin{lstlisting}[language=EasyCrypt,basicstyle=\ttfamily\scriptsize]
module CountedLazyROMBadGame(
  O : Oracle,
  A : CountedROMClient
) = {
\end{lstlisting}
\noindent\textbf{Formal mathematical reading.}
Mathematical definition. This is a probabilistic program/game or a module adapter.  Its procedures induce the probability space used by the game-based proof.

\noindent\textbf{Role in the proof.}
Use in this chapter. It defines the game or game interface whose success event is bounded by later theorems.

\end{formaldefinition}

\begin{formaldefinition}[EasyCrypt line 521]
\noindent\textbf{EasyCrypt name.}
\begin{lstlisting}[language=EasyCrypt,basicstyle=\ttfamily\scriptsize]
C
\end{lstlisting}
\noindent\textbf{Checked EasyCrypt statement.}
\begin{lstlisting}[language=EasyCrypt,basicstyle=\ttfamily\scriptsize]
module C = CountedLazyROM(O)
  module A = A(C)

  proc main() : bool = {
\end{lstlisting}
\noindent\textbf{Formal mathematical reading.}
Mathematical definition. This is a probabilistic program/game or a module adapter.  Its procedures induce the probability space used by the game-based proof.

\noindent\textbf{Role in the proof.}
Use in this chapter. It supplies an executable component of the formal probabilistic model.

\end{formaldefinition}

\begin{formaldefinition}[EasyCrypt line 534]
\noindent\textbf{EasyCrypt name.}
\begin{lstlisting}[language=EasyCrypt,basicstyle=\ttfamily\scriptsize]
ProgramChallenge
\end{lstlisting}
\noindent\textbf{Checked EasyCrypt statement.}
\begin{lstlisting}[language=EasyCrypt,basicstyle=\ttfamily\scriptsize]
module ProgramChallenge(O : Oracle) = {
\end{lstlisting}
\noindent\textbf{Formal mathematical reading.}
Mathematical definition. This is a probabilistic program/game or a module adapter.  Its procedures induce the probability space used by the game-based proof.

\noindent\textbf{Role in the proof.}
Use in this chapter. It supplies an executable component of the formal probabilistic model.

\end{formaldefinition}

\begin{formaldefinition}[EasyCrypt line 535]
\noindent\textbf{EasyCrypt name.}
\begin{lstlisting}[language=EasyCrypt,basicstyle=\ttfamily\scriptsize]
program
\end{lstlisting}
\noindent\textbf{Checked EasyCrypt statement.}
\begin{lstlisting}[language=EasyCrypt,basicstyle=\ttfamily\scriptsize]
proc program(site : programming_site) : unit = {
\end{lstlisting}
\noindent\textbf{Formal mathematical reading.}
Mathematical definition. This procedure is a command in the probabilistic program.  Its return value, state updates, and oracle calls are the operational content of the game.

\noindent\textbf{Role in the proof.}
Use in this chapter. It supplies an executable component of the formal probabilistic model.

\end{formaldefinition}

\begin{formaldefinition}[EasyCrypt line 540]
\noindent\textbf{EasyCrypt name.}
\begin{lstlisting}[language=EasyCrypt,basicstyle=\ttfamily\scriptsize]
TranscriptProgrammer
\end{lstlisting}
\noindent\textbf{Checked EasyCrypt statement.}
\begin{lstlisting}[language=EasyCrypt,basicstyle=\ttfamily\scriptsize]
module type TranscriptProgrammer(O : Oracle) = {
\end{lstlisting}
\noindent\textbf{Formal mathematical reading.}
Mathematical definition. This is an interface for an oracle, adversary, scheme, sampler, solver, or reduction.  A later theorem quantified over this interface is universally quantified over all modules implementing these procedures.

\noindent\textbf{Role in the proof.}
Use in this chapter. It supplies an executable component of the formal probabilistic model.

\end{formaldefinition}

\begin{formaldefinition}[EasyCrypt line 541]
\noindent\textbf{EasyCrypt name.}
\begin{lstlisting}[language=EasyCrypt,basicstyle=\ttfamily\scriptsize]
program
\end{lstlisting}
\noindent\textbf{Checked EasyCrypt statement.}
\begin{lstlisting}[language=EasyCrypt,basicstyle=\ttfamily\scriptsize]
proc program(md : mode, tr : transcript, y : ro_output) : unit
}.
\end{lstlisting}
\noindent\textbf{Formal mathematical reading.}
Mathematical definition. This procedure is a command in the probabilistic program.  Its return value, state updates, and oracle calls are the operational content of the game.

\noindent\textbf{Role in the proof.}
Use in this chapter. It supplies an executable component of the formal probabilistic model.

\end{formaldefinition}

\begin{formaldefinition}[EasyCrypt line 544]
\noindent\textbf{EasyCrypt name.}
\begin{lstlisting}[language=EasyCrypt,basicstyle=\ttfamily\scriptsize]
TranscriptProgrammerFromSite
\end{lstlisting}
\noindent\textbf{Checked EasyCrypt statement.}
\begin{lstlisting}[language=EasyCrypt,basicstyle=\ttfamily\scriptsize]
module TranscriptProgrammerFromSite(O : Oracle) = {
\end{lstlisting}
\noindent\textbf{Formal mathematical reading.}
Mathematical definition. This is a probabilistic program/game or a module adapter.  Its procedures induce the probability space used by the game-based proof.

\noindent\textbf{Role in the proof.}
Use in this chapter. It defines the oracle boundary through which adversaries and simulations interact.

\end{formaldefinition}

\begin{formaldefinition}[EasyCrypt line 545]
\noindent\textbf{EasyCrypt name.}
\begin{lstlisting}[language=EasyCrypt,basicstyle=\ttfamily\scriptsize]
program
\end{lstlisting}
\noindent\textbf{Checked EasyCrypt statement.}
\begin{lstlisting}[language=EasyCrypt,basicstyle=\ttfamily\scriptsize]
proc program(md : mode, tr : transcript, y : ro_output) : unit = {
\end{lstlisting}
\noindent\textbf{Formal mathematical reading.}
Mathematical definition. This procedure is a command in the probabilistic program.  Its return value, state updates, and oracle calls are the operational content of the game.

\noindent\textbf{Role in the proof.}
Use in this chapter. It supplies an executable component of the formal probabilistic model.

\end{formaldefinition}

\begin{formaldefinition}[EasyCrypt line 611]
\noindent\textbf{EasyCrypt name.}
\begin{lstlisting}[language=EasyCrypt,basicstyle=\ttfamily\scriptsize]
programmed_site_entropy_token_from_query
\end{lstlisting}
\noindent\textbf{Checked EasyCrypt statement.}
\begin{lstlisting}[language=EasyCrypt,basicstyle=\ttfamily\scriptsize]
op programmed_site_entropy_token_from_query (q : ro_query) : int =
  with q = ChallengeHashQuery x => nth 0 x.`3 0
  with q = MessageHashQuery _ => 0
  with q = MatrixExpandQuery _ => 0
  with q = SamplerExpandQuery _ => 0.
\end{lstlisting}
\noindent\textbf{Formal mathematical reading.}
Mathematical definition. This is a total EasyCrypt operation.  The exact displayed statement gives the domain, codomain, and defining equation when present.

\noindent\textbf{Role in the proof.}
Use in this chapter. It names a well-formedness, support, or validity predicate needed by later preservation lemmas.

\end{formaldefinition}

\begin{formaldefinition}[EasyCrypt line 884]
\noindent\textbf{EasyCrypt name.}
\begin{lstlisting}[language=EasyCrypt,basicstyle=\ttfamily\scriptsize]
programming_site_conflict_free
\end{lstlisting}
\noindent\textbf{Checked EasyCrypt statement.}
\begin{lstlisting}[language=EasyCrypt,basicstyle=\ttfamily\scriptsize]
op programming_site_conflict_free
   (site : programming_site) (sites : programming_site list) : bool =
  all (fun old => ! programming_sites_conflict site old) sites.
\end{lstlisting}
\noindent\textbf{Formal mathematical reading.}
Mathematical definition. This is a Boolean predicate.  The exact displayed EasyCrypt statement gives its arguments; mathematically it is an event, invariant, support condition, or well-formedness predicate over those arguments.

\noindent\textbf{Role in the proof.}
Use in this chapter. It is part of the exact formal vocabulary used by subsequent lemmas and games.

\end{formaldefinition}

\begin{formaldefinition}[EasyCrypt line 888]
\noindent\textbf{EasyCrypt name.}
\begin{lstlisting}[language=EasyCrypt,basicstyle=\ttfamily\scriptsize]
programming_site_log_conflict_free_for_honest
\end{lstlisting}
\noindent\textbf{Checked EasyCrypt statement.}
\begin{lstlisting}[language=EasyCrypt,basicstyle=\ttfamily\scriptsize]
op programming_site_log_conflict_free_for_honest
   (md : mode) (sk : skey) (sites : programming_site list) : bool =
  forall m ctx coins,
    programming_site_conflict_free
      (honest_signing_programming_site md sk m ctx coins) sites.
\end{lstlisting}
\noindent\textbf{Formal mathematical reading.}
Mathematical definition. This is a Boolean predicate.  The exact displayed EasyCrypt statement gives its arguments; mathematically it is an event, invariant, support condition, or well-formedness predicate over those arguments.

\noindent\textbf{Role in the proof.}
Use in this chapter. It is part of the exact formal vocabulary used by subsequent lemmas and games.

\end{formaldefinition}

\begin{formaldefinition}[EasyCrypt line 1460]
\noindent\textbf{EasyCrypt name.}
\begin{lstlisting}[language=EasyCrypt,basicstyle=\ttfamily\scriptsize]
counted_lazy_rom_bad_flags_budget_sound
\end{lstlisting}
\noindent\textbf{Checked EasyCrypt statement.}
\begin{lstlisting}[language=EasyCrypt,basicstyle=\ttfamily\scriptsize]
op counted_lazy_rom_bad_flags_budget_sound (p : real) : bool =
  p <= counted_rom_programming_loss_term.
\end{lstlisting}
\noindent\textbf{Formal mathematical reading.}
Mathematical definition. This is a Boolean predicate.  The exact displayed EasyCrypt statement gives its arguments; mathematically it is an event, invariant, support condition, or well-formedness predicate over those arguments.

\noindent\textbf{Role in the proof.}
Use in this chapter. It names a bad event or rejection condition whose probability must be zero or explicitly bounded.

\end{formaldefinition}

\subsubsection*{Lemmas, Theorems, Relational Equivalences, and Their Proof Dependencies}
\begin{formaltheorem}[EasyCrypt line 142]
\noindent\textbf{EasyCrypt name.}
\begin{lstlisting}[language=EasyCrypt,basicstyle=\ttfamily\scriptsize]
challenge_query_log_count_cons_le
\end{lstlisting}
\noindent\textbf{Checked EasyCrypt statement.}
\begin{lstlisting}[language=EasyCrypt,basicstyle=\ttfamily\scriptsize]
lemma challenge_query_log_count_cons_le q qs :
  challenge_query_log_count (q :: qs) <=
  challenge_query_log_count qs + 1.
\end{lstlisting}
\noindent\textbf{Formal mathematical reading.}
Mathematical theorem. This is an inequality, usually an arithmetic loss bound, event inclusion bound, or monotonicity fact used to compose probabilities.

\noindent\textbf{Role in the proof.}
Use in this chapter. It contributes directly to an advantage bound, bad-event bound, or real-arithmetic composition step.

\noindent\textbf{Proof-script interpretation.}
Proof-script reading. The machine proof decomposes the theorem into rewrites, program-logic obligations, relational equivalences, and arithmetic side conditions.
\emph{Main tactics visible in the proof script:} \ec{rewrite}, \ec{case}, \ec{smt}.
\emph{Named identifiers syntactically cited in the proof block:}
\begin{lstlisting}[language=EasyCrypt,basicstyle=\ttfamily\scriptsize]
challenge_query_log_count
fcardU1
fsetUC
oflist_cons
site
\end{lstlisting}

\end{formaltheorem}

\begin{formaltheorem}[EasyCrypt line 156]
\noindent\textbf{EasyCrypt name.}
\begin{lstlisting}[language=EasyCrypt,basicstyle=\ttfamily\scriptsize]
challenge_query_log_count_nonchallenge_cons
\end{lstlisting}
\noindent\textbf{Checked EasyCrypt statement.}
\begin{lstlisting}[language=EasyCrypt,basicstyle=\ttfamily\scriptsize]
lemma challenge_query_log_count_nonchallenge_cons q qs :
  ! ro_query_is_challenge q =>
  challenge_query_log_count (q :: qs) =
  challenge_query_log_count qs.
\end{lstlisting}
\noindent\textbf{Formal mathematical reading.}
Mathematical theorem. This is an implication, normally turning one invariant, validity predicate, or proof obligation into another.

\noindent\textbf{Role in the proof.}
Use in this chapter. It contributes directly to an advantage bound, bad-event bound, or real-arithmetic composition step.

\noindent\textbf{Proof-script interpretation.}
Proof-script reading. The machine proof decomposes the theorem into rewrites, program-logic obligations, relational equivalences, and arithmetic side conditions.
\emph{Main tactics visible in the proof script:} \ec{rewrite}, \ec{case}.
\emph{Named identifiers syntactically cited in the proof block:}
\begin{lstlisting}[language=EasyCrypt,basicstyle=\ttfamily\scriptsize]
challenge_query_log_count
\end{lstlisting}

\end{formaltheorem}

\begin{formaltheorem}[EasyCrypt line 169]
\noindent\textbf{EasyCrypt name.}
\begin{lstlisting}[language=EasyCrypt,basicstyle=\ttfamily\scriptsize]
challenge_query_log_count_message_cons
\end{lstlisting}
\noindent\textbf{Checked EasyCrypt statement.}
\begin{lstlisting}[language=EasyCrypt,basicstyle=\ttfamily\scriptsize]
lemma challenge_query_log_count_message_cons pk ctx m qs :
  challenge_query_log_count (message_hash_query pk ctx m :: qs) =
  challenge_query_log_count qs.
\end{lstlisting}
\noindent\textbf{Formal mathematical reading.}
Mathematical theorem. This is an equality or definitional expansion used for rewriting and normalization.

\noindent\textbf{Role in the proof.}
Use in this chapter. It contributes directly to an advantage bound, bad-event bound, or real-arithmetic composition step.

\noindent\textbf{Proof-script interpretation.}
Proof-script reading. The machine proof decomposes the theorem into rewrites, program-logic obligations, relational equivalences, and arithmetic side conditions.
\emph{Main tactics visible in the proof script:} \ec{apply}, \ec{rewrite}.
\emph{Named identifiers syntactically cited in the proof block:}
\begin{lstlisting}[language=EasyCrypt,basicstyle=\ttfamily\scriptsize]
challenge_query_log_count_nonchallenge_cons
message_hash_query
ro_query_is_challenge
\end{lstlisting}

\end{formaltheorem}

\begin{formaltheorem}[EasyCrypt line 177]
\noindent\textbf{EasyCrypt name.}
\begin{lstlisting}[language=EasyCrypt,basicstyle=\ttfamily\scriptsize]
challenge_query_log_count_matrix_cons
\end{lstlisting}
\noindent\textbf{Checked EasyCrypt statement.}
\begin{lstlisting}[language=EasyCrypt,basicstyle=\ttfamily\scriptsize]
lemma challenge_query_log_count_matrix_cons md sd qs :
  challenge_query_log_count (matrix_expand_query md sd :: qs) =
  challenge_query_log_count qs.
\end{lstlisting}
\noindent\textbf{Formal mathematical reading.}
Mathematical theorem. This is an equality or definitional expansion used for rewriting and normalization.

\noindent\textbf{Role in the proof.}
Use in this chapter. It contributes directly to an advantage bound, bad-event bound, or real-arithmetic composition step.

\noindent\textbf{Proof-script interpretation.}
Proof-script reading. The machine proof decomposes the theorem into rewrites, program-logic obligations, relational equivalences, and arithmetic side conditions.
\emph{Main tactics visible in the proof script:} \ec{apply}, \ec{rewrite}.
\emph{Named identifiers syntactically cited in the proof block:}
\begin{lstlisting}[language=EasyCrypt,basicstyle=\ttfamily\scriptsize]
challenge_query_log_count_nonchallenge_cons
matrix_expand_query
ro_query_is_challenge
\end{lstlisting}

\end{formaltheorem}

\begin{formaltheorem}[EasyCrypt line 185]
\noindent\textbf{EasyCrypt name.}
\begin{lstlisting}[language=EasyCrypt,basicstyle=\ttfamily\scriptsize]
challenge_query_log_count_sampler_cons
\end{lstlisting}
\noindent\textbf{Checked EasyCrypt statement.}
\begin{lstlisting}[language=EasyCrypt,basicstyle=\ttfamily\scriptsize]
lemma challenge_query_log_count_sampler_cons coins qs :
  challenge_query_log_count (sampler_expand_query coins :: qs) =
  challenge_query_log_count qs.
\end{lstlisting}
\noindent\textbf{Formal mathematical reading.}
Mathematical theorem. This is an equality or definitional expansion used for rewriting and normalization.

\noindent\textbf{Role in the proof.}
Use in this chapter. It contributes directly to an advantage bound, bad-event bound, or real-arithmetic composition step.

\noindent\textbf{Proof-script interpretation.}
Proof-script reading. The machine proof decomposes the theorem into rewrites, program-logic obligations, relational equivalences, and arithmetic side conditions.
\emph{Main tactics visible in the proof script:} \ec{apply}, \ec{rewrite}.
\emph{Named identifiers syntactically cited in the proof block:}
\begin{lstlisting}[language=EasyCrypt,basicstyle=\ttfamily\scriptsize]
challenge_query_log_count_nonchallenge_cons
ro_query_is_challenge
sampler_expand_query
\end{lstlisting}

\end{formaltheorem}

\begin{formaltheorem}[EasyCrypt line 319]
\noindent\textbf{EasyCrypt name.}
\begin{lstlisting}[language=EasyCrypt,basicstyle=\ttfamily\scriptsize]
counted_lazy_rom_init_clears
\end{lstlisting}
\noindent\textbf{Checked EasyCrypt statement.}
\begin{lstlisting}[language=EasyCrypt,basicstyle=\ttfamily\scriptsize]
lemma counted_lazy_rom_init_clears :
  hoare[CountedLazyROM(O).init :
    true ==>
    ! CountedLazyROM.bad_prequery /\
    ! CountedLazyROM.bad_reprogram /\
    ! CountedLazyROM.bad_min_entropy /\
    CountedLazyROM.programmed_sites = []].
\end{lstlisting}
\noindent\textbf{Formal mathematical reading.}
Mathematical theorem. This is a relational program theorem: two games or procedures are coupled so that a precondition on their initial states implies the stated postcondition on final states and return values.

\noindent\textbf{Role in the proof.}
Use in this chapter. It contributes directly to an advantage bound, bad-event bound, or real-arithmetic composition step.

\noindent\textbf{Proof-script interpretation.}
Proof-script reading. The machine proof decomposes the theorem into rewrites, program-logic obligations, relational equivalences, and arithmetic side conditions.
\emph{Main tactics visible in the proof script:} \ec{proc}, \ec{wp}, \ec{call}.

\end{formaltheorem}

\begin{formaltheorem}[EasyCrypt line 333]
\noindent\textbf{EasyCrypt name.}
\begin{lstlisting}[language=EasyCrypt,basicstyle=\ttfamily\scriptsize]
counted_lazy_rom_init_clears_hash_queries
\end{lstlisting}
\noindent\textbf{Checked EasyCrypt statement.}
\begin{lstlisting}[language=EasyCrypt,basicstyle=\ttfamily\scriptsize]
lemma counted_lazy_rom_init_clears_hash_queries :
  hoare[CountedLazyROM(O).init :
    true ==> CountedLazyROM.hash_queries = []].
\end{lstlisting}
\noindent\textbf{Formal mathematical reading.}
Mathematical theorem. This is a relational program theorem: two games or procedures are coupled so that a precondition on their initial states implies the stated postcondition on final states and return values.

\noindent\textbf{Role in the proof.}
Use in this chapter. It contributes directly to an advantage bound, bad-event bound, or real-arithmetic composition step.

\noindent\textbf{Proof-script interpretation.}
Proof-script reading. The machine proof decomposes the theorem into rewrites, program-logic obligations, relational equivalences, and arithmetic side conditions.
\emph{Main tactics visible in the proof script:} \ec{proc}, \ec{wp}, \ec{call}.

\end{formaltheorem}

\begin{formaltheorem}[EasyCrypt line 343]
\noindent\textbf{EasyCrypt name.}
\begin{lstlisting}[language=EasyCrypt,basicstyle=\ttfamily\scriptsize]
counted_lazy_rom_get_preserves_clear
\end{lstlisting}
\noindent\textbf{Checked EasyCrypt statement.}
\begin{lstlisting}[language=EasyCrypt,basicstyle=\ttfamily\scriptsize]
lemma counted_lazy_rom_get_preserves_clear :
  hoare[CountedLazyROM(O).get :
    ! CountedLazyROM.bad_prequery /\
    ! CountedLazyROM.bad_reprogram /\
    ! CountedLazyROM.bad_min_entropy /\
    CountedLazyROM.programmed_sites = [] ==>
    ! CountedLazyROM.bad_prequery /\
    ! CountedLazyROM.bad_reprogram /\
    ! CountedLazyROM.bad_min_entropy /\
    CountedLazyROM.programmed_sites = []].
\end{lstlisting}
\noindent\textbf{Formal mathematical reading.}
Mathematical theorem. This is a relational program theorem: two games or procedures are coupled so that a precondition on their initial states implies the stated postcondition on final states and return values.

\noindent\textbf{Role in the proof.}
Use in this chapter. It contributes directly to an advantage bound, bad-event bound, or real-arithmetic composition step.

\noindent\textbf{Proof-script interpretation.}
Proof-script reading. The machine proof decomposes the theorem into rewrites, program-logic obligations, relational equivalences, and arithmetic side conditions.
\emph{Main tactics visible in the proof script:} \ec{proc}, \ec{wp}, \ec{call}.

\end{formaltheorem}

\begin{formaltheorem}[EasyCrypt line 360]
\noindent\textbf{EasyCrypt name.}
\begin{lstlisting}[language=EasyCrypt,basicstyle=\ttfamily\scriptsize]
counted_lazy_rom_get_nonchallenge_preserves_challenge_query_count
\end{lstlisting}
\noindent\textbf{Checked EasyCrypt statement.}
\begin{lstlisting}[language=EasyCrypt,basicstyle=\ttfamily\scriptsize]
lemma counted_lazy_rom_get_nonchallenge_preserves_challenge_query_count n :
  hoare[CountedLazyROM(O).get :
    ! ro_query_is_challenge q /\
    challenge_query_log_count CountedLazyROM.hash_queries <= n ==>
    challenge_query_log_count CountedLazyROM.hash_queries <= n].
\end{lstlisting}
\noindent\textbf{Formal mathematical reading.}
Mathematical theorem. This is a relational program theorem: two games or procedures are coupled so that a precondition on their initial states implies the stated postcondition on final states and return values.

\noindent\textbf{Role in the proof.}
Use in this chapter. It contributes directly to an advantage bound, bad-event bound, or real-arithmetic composition step.

\noindent\textbf{Proof-script interpretation.}
Proof-script reading. The machine proof decomposes the theorem into rewrites, program-logic obligations, relational equivalences, and arithmetic side conditions.
\emph{Main tactics visible in the proof script:} \ec{proc}, \ec{wp}, \ec{call}, \ec{smt}.

\end{formaltheorem}

\begin{formaltheorem}[EasyCrypt line 372]
\noindent\textbf{EasyCrypt name.}
\begin{lstlisting}[language=EasyCrypt,basicstyle=\ttfamily\scriptsize]
counted_lazy_rom_get_preserves_bad_clear
\end{lstlisting}
\noindent\textbf{Checked EasyCrypt statement.}
\begin{lstlisting}[language=EasyCrypt,basicstyle=\ttfamily\scriptsize]
lemma counted_lazy_rom_get_preserves_bad_clear :
  hoare[CountedLazyROM(O).get :
    ! CountedLazyROM.bad_prequery /\
    ! CountedLazyROM.bad_reprogram /\
    ! CountedLazyROM.bad_min_entropy ==>
    ! CountedLazyROM.bad_prequery /\
    ! CountedLazyROM.bad_reprogram /\
    ! CountedLazyROM.bad_min_entropy].
\end{lstlisting}
\noindent\textbf{Formal mathematical reading.}
Mathematical theorem. This is a relational program theorem: two games or procedures are coupled so that a precondition on their initial states implies the stated postcondition on final states and return values.

\noindent\textbf{Role in the proof.}
Use in this chapter. It contributes directly to an advantage bound, bad-event bound, or real-arithmetic composition step.

\noindent\textbf{Proof-script interpretation.}
Proof-script reading. The machine proof decomposes the theorem into rewrites, program-logic obligations, relational equivalences, and arithmetic side conditions.
\emph{Main tactics visible in the proof script:} \ec{proc}, \ec{wp}, \ec{call}.

\end{formaltheorem}

\begin{formaltheorem}[EasyCrypt line 387]
\noindent\textbf{EasyCrypt name.}
\begin{lstlisting}[language=EasyCrypt,basicstyle=\ttfamily\scriptsize]
counted_lazy_rom_get_preserves_min_entropy_clear
\end{lstlisting}
\noindent\textbf{Checked EasyCrypt statement.}
\begin{lstlisting}[language=EasyCrypt,basicstyle=\ttfamily\scriptsize]
lemma counted_lazy_rom_get_preserves_min_entropy_clear :
  hoare[CountedLazyROM(O).get :
    ! CountedLazyROM.bad_min_entropy ==>
    ! CountedLazyROM.bad_min_entropy].
\end{lstlisting}
\noindent\textbf{Formal mathematical reading.}
Mathematical theorem. This is a relational program theorem: two games or procedures are coupled so that a precondition on their initial states implies the stated postcondition on final states and return values.

\noindent\textbf{Role in the proof.}
Use in this chapter. It contributes directly to an advantage bound, bad-event bound, or real-arithmetic composition step.

\noindent\textbf{Proof-script interpretation.}
Proof-script reading. The machine proof decomposes the theorem into rewrites, program-logic obligations, relational equivalences, and arithmetic side conditions.
\emph{Main tactics visible in the proof script:} \ec{proc}, \ec{wp}, \ec{call}.

\end{formaltheorem}

\begin{formaltheorem}[EasyCrypt line 398]
\noindent\textbf{EasyCrypt name.}
\begin{lstlisting}[language=EasyCrypt,basicstyle=\ttfamily\scriptsize]
counted_lazy_rom_observe_transcript_preserves_clear
\end{lstlisting}
\noindent\textbf{Checked EasyCrypt statement.}
\begin{lstlisting}[language=EasyCrypt,basicstyle=\ttfamily\scriptsize]
lemma counted_lazy_rom_observe_transcript_preserves_clear :
  hoare[CountedLazyROM(O).observe_transcript :
    ! min_entropy_failure md tr /\
    ! CountedLazyROM.bad_prequery /\
    ! CountedLazyROM.bad_reprogram /\
    ! CountedLazyROM.bad_min_entropy /\
    CountedLazyROM.programmed_sites = [] ==>
    ! CountedLazyROM.bad_prequery /\
    ! CountedLazyROM.bad_reprogram /\
    ! CountedLazyROM.bad_min_entropy /\
    CountedLazyROM.programmed_sites = []].
\end{lstlisting}
\noindent\textbf{Formal mathematical reading.}
Mathematical theorem. This is a relational program theorem: two games or procedures are coupled so that a precondition on their initial states implies the stated postcondition on final states and return values.

\noindent\textbf{Role in the proof.}
Use in this chapter. It contributes directly to an advantage bound, bad-event bound, or real-arithmetic composition step.

\noindent\textbf{Proof-script interpretation.}
Proof-script reading. The machine proof decomposes the theorem into rewrites, program-logic obligations, relational equivalences, and arithmetic side conditions.
\emph{Main tactics visible in the proof script:} \ec{proc}, \ec{wp}, \ec{rewrite}.
\emph{Named identifiers syntactically cited in the proof block:}
\begin{lstlisting}[language=EasyCrypt,basicstyle=\ttfamily\scriptsize]
entropy_ok
hr
min_ok
preq_ok
programmed_ok
reprog_ok
\end{lstlisting}

\end{formaltheorem}

\begin{formaltheorem}[EasyCrypt line 417]
\noindent\textbf{EasyCrypt name.}
\begin{lstlisting}[language=EasyCrypt,basicstyle=\ttfamily\scriptsize]
counted_lazy_rom_observe_transcript_preserves_min_entropy_clear
\end{lstlisting}
\noindent\textbf{Checked EasyCrypt statement.}
\begin{lstlisting}[language=EasyCrypt,basicstyle=\ttfamily\scriptsize]
lemma counted_lazy_rom_observe_transcript_preserves_min_entropy_clear :
  hoare[CountedLazyROM(O).observe_transcript :
    ! min_entropy_failure md tr /\
    ! CountedLazyROM.bad_min_entropy ==>
    ! CountedLazyROM.bad_min_entropy].
\end{lstlisting}
\noindent\textbf{Formal mathematical reading.}
Mathematical theorem. This is a relational program theorem: two games or procedures are coupled so that a precondition on their initial states implies the stated postcondition on final states and return values.

\noindent\textbf{Role in the proof.}
Use in this chapter. It contributes directly to an advantage bound, bad-event bound, or real-arithmetic composition step.

\noindent\textbf{Proof-script interpretation.}
Proof-script reading. The machine proof decomposes the theorem into rewrites, program-logic obligations, relational equivalences, and arithmetic side conditions.
\emph{Main tactics visible in the proof script:} \ec{proc}, \ec{wp}, \ec{rewrite}.
\emph{Named identifiers syntactically cited in the proof block:}
\begin{lstlisting}[language=EasyCrypt,basicstyle=\ttfamily\scriptsize]
entropy_ok
hr
min_ok
\end{lstlisting}

\end{formaltheorem}

\begin{formaltheorem}[EasyCrypt line 430]
\noindent\textbf{EasyCrypt name.}
\begin{lstlisting}[language=EasyCrypt,basicstyle=\ttfamily\scriptsize]
counted_lazy_rom_observe_transcript_preserves_bad_clear
\end{lstlisting}
\noindent\textbf{Checked EasyCrypt statement.}
\begin{lstlisting}[language=EasyCrypt,basicstyle=\ttfamily\scriptsize]
lemma counted_lazy_rom_observe_transcript_preserves_bad_clear :
  hoare[CountedLazyROM(O).observe_transcript :
    ! min_entropy_failure md tr /\
    ! CountedLazyROM.bad_prequery /\
    ! CountedLazyROM.bad_reprogram /\
    ! CountedLazyROM.bad_min_entropy ==>
    ! CountedLazyROM.bad_prequery /\
    ! CountedLazyROM.bad_reprogram /\
    ! CountedLazyROM.bad_min_entropy].
\end{lstlisting}
\noindent\textbf{Formal mathematical reading.}
Mathematical theorem. This is a relational program theorem: two games or procedures are coupled so that a precondition on their initial states implies the stated postcondition on final states and return values.

\noindent\textbf{Role in the proof.}
Use in this chapter. It contributes directly to an advantage bound, bad-event bound, or real-arithmetic composition step.

\noindent\textbf{Proof-script interpretation.}
Proof-script reading. The machine proof decomposes the theorem into rewrites, program-logic obligations, relational equivalences, and arithmetic side conditions.
\emph{Main tactics visible in the proof script:} \ec{proc}, \ec{wp}, \ec{rewrite}.
\emph{Named identifiers syntactically cited in the proof block:}
\begin{lstlisting}[language=EasyCrypt,basicstyle=\ttfamily\scriptsize]
entropy_ok
hr
min_ok
preq_ok
reprog_ok
\end{lstlisting}

\end{formaltheorem}

\begin{formaltheorem}[EasyCrypt line 447]
\noindent\textbf{EasyCrypt name.}
\begin{lstlisting}[language=EasyCrypt,basicstyle=\ttfamily\scriptsize]
counted_lazy_rom_program_preserves_bad_clear_if_fresh
\end{lstlisting}
\noindent\textbf{Checked EasyCrypt statement.}
\begin{lstlisting}[language=EasyCrypt,basicstyle=\ttfamily\scriptsize]
lemma counted_lazy_rom_program_preserves_bad_clear_if_fresh :
  hoare[CountedLazyROM(O).program :
    ! CountedLazyROM.bad_prequery /\
    ! CountedLazyROM.bad_reprogram /\
    ! CountedLazyROM.bad_min_entropy /\
    programming_site_fresh
      CountedLazyROM.hash_queries
      CountedLazyROM.programmed_sites site ==>
    ! CountedLazyROM.bad_prequery /\
    ! CountedLazyROM.bad_reprogram /\
    ! CountedLazyROM.bad_min_entropy].
\end{lstlisting}
\noindent\textbf{Formal mathematical reading.}
Mathematical theorem. This is a relational program theorem: two games or procedures are coupled so that a precondition on their initial states implies the stated postcondition on final states and return values.

\noindent\textbf{Role in the proof.}
Use in this chapter. It contributes directly to an advantage bound, bad-event bound, or real-arithmetic composition step.

\noindent\textbf{Proof-script interpretation.}
Proof-script reading. The machine proof decomposes the theorem into rewrites, program-logic obligations, relational equivalences, and arithmetic side conditions.
\emph{Main tactics visible in the proof script:} \ec{proc}, \ec{wp}, \ec{call}.

\end{formaltheorem}

\begin{formaltheorem}[EasyCrypt line 465]
\noindent\textbf{EasyCrypt name.}
\begin{lstlisting}[language=EasyCrypt,basicstyle=\ttfamily\scriptsize]
counted_lazy_rom_program_preserves_min_entropy_clear
\end{lstlisting}
\noindent\textbf{Checked EasyCrypt statement.}
\begin{lstlisting}[language=EasyCrypt,basicstyle=\ttfamily\scriptsize]
lemma counted_lazy_rom_program_preserves_min_entropy_clear :
  hoare[CountedLazyROM(O).program :
    ! CountedLazyROM.bad_min_entropy ==>
    ! CountedLazyROM.bad_min_entropy].
\end{lstlisting}
\noindent\textbf{Formal mathematical reading.}
Mathematical theorem. This is a relational program theorem: two games or procedures are coupled so that a precondition on their initial states implies the stated postcondition on final states and return values.

\noindent\textbf{Role in the proof.}
Use in this chapter. It contributes directly to an advantage bound, bad-event bound, or real-arithmetic composition step.

\noindent\textbf{Proof-script interpretation.}
Proof-script reading. The machine proof decomposes the theorem into rewrites, program-logic obligations, relational equivalences, and arithmetic side conditions.
\emph{Main tactics visible in the proof script:} \ec{proc}, \ec{wp}, \ec{call}.

\end{formaltheorem}

\begin{formaltheorem}[EasyCrypt line 476]
\noindent\textbf{EasyCrypt name.}
\begin{lstlisting}[language=EasyCrypt,basicstyle=\ttfamily\scriptsize]
counted_lazy_rom_bad_preserves_min_entropy_clear
\end{lstlisting}
\noindent\textbf{Checked EasyCrypt statement.}
\begin{lstlisting}[language=EasyCrypt,basicstyle=\ttfamily\scriptsize]
lemma counted_lazy_rom_bad_preserves_min_entropy_clear :
  hoare[CountedLazyROM(O).bad :
    ! CountedLazyROM.bad_min_entropy ==>
    ! CountedLazyROM.bad_min_entropy].
\end{lstlisting}
\noindent\textbf{Formal mathematical reading.}
Mathematical theorem. This is a relational program theorem: two games or procedures are coupled so that a precondition on their initial states implies the stated postcondition on final states and return values.

\noindent\textbf{Role in the proof.}
Use in this chapter. It contributes directly to an advantage bound, bad-event bound, or real-arithmetic composition step.

\noindent\textbf{Proof-script interpretation.}
Proof-script reading. The machine proof decomposes the theorem into rewrites, program-logic obligations, relational equivalences, and arithmetic side conditions.
\emph{Main tactics visible in the proof script:} \ec{proc}.

\end{formaltheorem}

\begin{formaltheorem}[EasyCrypt line 486]
\noindent\textbf{EasyCrypt name.}
\begin{lstlisting}[language=EasyCrypt,basicstyle=\ttfamily\scriptsize]
counted_lazy_rom_bad_false_when_clear
\end{lstlisting}
\noindent\textbf{Checked EasyCrypt statement.}
\begin{lstlisting}[language=EasyCrypt,basicstyle=\ttfamily\scriptsize]
lemma counted_lazy_rom_bad_false_when_clear :
  hoare[CountedLazyROM(O).bad :
    ! CountedLazyROM.bad_prequery /\
    ! CountedLazyROM.bad_reprogram /\
    ! CountedLazyROM.bad_min_entropy ==>
    ! res].
\end{lstlisting}
\noindent\textbf{Formal mathematical reading.}
Mathematical theorem. This is a relational program theorem: two games or procedures are coupled so that a precondition on their initial states implies the stated postcondition on final states and return values.

\noindent\textbf{Role in the proof.}
Use in this chapter. It contributes directly to an advantage bound, bad-event bound, or real-arithmetic composition step.

\noindent\textbf{Proof-script interpretation.}
Proof-script reading. The machine proof decomposes the theorem into rewrites, program-logic obligations, relational equivalences, and arithmetic side conditions.
\emph{Main tactics visible in the proof script:} \ec{proc}, \ec{rewrite}.
\emph{Named identifiers syntactically cited in the proof block:}
\begin{lstlisting}[language=EasyCrypt,basicstyle=\ttfamily\scriptsize]
entropy_ok
hr
preq_ok
reprog_ok
\end{lstlisting}

\end{formaltheorem}

\begin{formaltheorem}[EasyCrypt line 550]
\noindent\textbf{EasyCrypt name.}
\begin{lstlisting}[language=EasyCrypt,basicstyle=\ttfamily\scriptsize]
programming_site_self_matches
\end{lstlisting}
\noindent\textbf{Checked EasyCrypt statement.}
\begin{lstlisting}[language=EasyCrypt,basicstyle=\ttfamily\scriptsize]
lemma programming_site_self_matches md tr y :
  transcript_challenge_matches tr y =>
  programming_site_matches_transcript md tr
    (programming_site_of_transcript md tr y).
\end{lstlisting}
\noindent\textbf{Formal mathematical reading.}
Mathematical theorem. This is an implication, normally turning one invariant, validity predicate, or proof obligation into another.

\noindent\textbf{Role in the proof.}
Use in this chapter. It is a reusable checked theorem cited by later proof scripts.

\noindent\textbf{Proof-script interpretation.}
Proof-script reading. The machine proof decomposes the theorem into rewrites, program-logic obligations, relational equivalences, and arithmetic side conditions.
\emph{Main tactics visible in the proof script:} \ec{rewrite}.
\emph{Named identifiers syntactically cited in the proof block:}
\begin{lstlisting}[language=EasyCrypt,basicstyle=\ttfamily\scriptsize]
programming_site_matches_transcript
programming_site_of_transcript
programming_site_output
programming_site_query
\end{lstlisting}

\end{formaltheorem}

\begin{formaltheorem}[EasyCrypt line 560]
\noindent\textbf{EasyCrypt name.}
\begin{lstlisting}[language=EasyCrypt,basicstyle=\ttfamily\scriptsize]
programming_site_self_no_reprogramming_failure
\end{lstlisting}
\noindent\textbf{Checked EasyCrypt statement.}
\begin{lstlisting}[language=EasyCrypt,basicstyle=\ttfamily\scriptsize]
lemma programming_site_self_no_reprogramming_failure md tr y :
  transcript_challenge_matches tr y =>
  ! reprogramming_failure md tr (programming_site_of_transcript md tr y).
\end{lstlisting}
\noindent\textbf{Formal mathematical reading.}
Mathematical theorem. This is an implication, normally turning one invariant, validity predicate, or proof obligation into another.

\noindent\textbf{Role in the proof.}
Use in this chapter. It is a reusable checked theorem cited by later proof scripts.

\noindent\textbf{Proof-script interpretation.}
Proof-script reading. The machine proof decomposes the theorem into rewrites, program-logic obligations, relational equivalences, and arithmetic side conditions.
\emph{Main tactics visible in the proof script:} \ec{rewrite}.
\emph{Named identifiers syntactically cited in the proof block:}
\begin{lstlisting}[language=EasyCrypt,basicstyle=\ttfamily\scriptsize]
match_y
md
programming_site_self_matches
reprogramming_failure
tr
\end{lstlisting}

\end{formaltheorem}

\begin{formaltheorem}[EasyCrypt line 569]
\noindent\textbf{EasyCrypt name.}
\begin{lstlisting}[language=EasyCrypt,basicstyle=\ttfamily\scriptsize]
fs_with_aborts_failureE
\end{lstlisting}
\noindent\textbf{Checked EasyCrypt statement.}
\begin{lstlisting}[language=EasyCrypt,basicstyle=\ttfamily\scriptsize]
lemma fs_with_aborts_failureE md tr site :
  fs_with_aborts_failure md tr site =
  (reprogramming_failure md tr site \/ min_entropy_failure md tr).
\end{lstlisting}
\noindent\textbf{Formal mathematical reading.}
Mathematical theorem. This is an equality or definitional expansion used for rewriting and normalization.

\noindent\textbf{Role in the proof.}
Use in this chapter. It is a reusable checked theorem cited by later proof scripts.

\noindent\textbf{Proof-script interpretation.}
No proof block was collected for this item.

\end{formaltheorem}

\begin{formaltheorem}[EasyCrypt line 574]
\noindent\textbf{EasyCrypt name.}
\begin{lstlisting}[language=EasyCrypt,basicstyle=\ttfamily\scriptsize]
programming_site_prequeried_cons
\end{lstlisting}
\noindent\textbf{Checked EasyCrypt statement.}
\begin{lstlisting}[language=EasyCrypt,basicstyle=\ttfamily\scriptsize]
lemma programming_site_prequeried_cons qs site :
  programming_site_prequeried (programming_site_query site :: qs) site.
\end{lstlisting}
\noindent\textbf{Formal mathematical reading.}
Mathematical theorem. This lemma is a universally quantified proposition over the variables shown in the EasyCrypt statement.

\noindent\textbf{Role in the proof.}
Use in this chapter. It is a reusable checked theorem cited by later proof scripts.

\noindent\textbf{Proof-script interpretation.}
No proof block was collected for this item.

\end{formaltheorem}

\begin{formaltheorem}[EasyCrypt line 578]
\noindent\textbf{EasyCrypt name.}
\begin{lstlisting}[language=EasyCrypt,basicstyle=\ttfamily\scriptsize]
programming_site_prequeried_cons_preserve
\end{lstlisting}
\noindent\textbf{Checked EasyCrypt statement.}
\begin{lstlisting}[language=EasyCrypt,basicstyle=\ttfamily\scriptsize]
lemma programming_site_prequeried_cons_preserve q qs site :
  programming_site_prequeried qs site =>
  programming_site_prequeried (q :: qs) site.
\end{lstlisting}
\noindent\textbf{Formal mathematical reading.}
Mathematical theorem. This is an implication, normally turning one invariant, validity predicate, or proof obligation into another.

\noindent\textbf{Role in the proof.}
Use in this chapter. It preserves an invariant across an oracle call, adversary call, or game procedure.

\noindent\textbf{Proof-script interpretation.}
No proof block was collected for this item.

\end{formaltheorem}

\begin{formaltheorem}[EasyCrypt line 583]
\noindent\textbf{EasyCrypt name.}
\begin{lstlisting}[language=EasyCrypt,basicstyle=\ttfamily\scriptsize]
programming_site_prequeried_output_irrelevant
\end{lstlisting}
\noindent\textbf{Checked EasyCrypt statement.}
\begin{lstlisting}[language=EasyCrypt,basicstyle=\ttfamily\scriptsize]
lemma programming_site_prequeried_output_irrelevant qs q y1 y2 :
  programming_site_prequeried qs (q, y1) =
  programming_site_prequeried qs (q, y2).
\end{lstlisting}
\noindent\textbf{Formal mathematical reading.}
Mathematical theorem. This is an equality or definitional expansion used for rewriting and normalization.

\noindent\textbf{Role in the proof.}
Use in this chapter. It contributes directly to an advantage bound, bad-event bound, or real-arithmetic composition step.

\noindent\textbf{Proof-script interpretation.}
No proof block was collected for this item.

\end{formaltheorem}

\begin{formaltheorem}[EasyCrypt line 588]
\noindent\textbf{EasyCrypt name.}
\begin{lstlisting}[language=EasyCrypt,basicstyle=\ttfamily\scriptsize]
programming_site_prequeried_witness
\end{lstlisting}
\noindent\textbf{Checked EasyCrypt statement.}
\begin{lstlisting}[language=EasyCrypt,basicstyle=\ttfamily\scriptsize]
lemma programming_site_prequeried_witness qs q y :
  q \in qs =>
  programming_site_prequeried qs (q, y).
\end{lstlisting}
\noindent\textbf{Formal mathematical reading.}
Mathematical theorem. This is an implication, normally turning one invariant, validity predicate, or proof obligation into another.

\noindent\textbf{Role in the proof.}
Use in this chapter. It is a reusable checked theorem cited by later proof scripts.

\noindent\textbf{Proof-script interpretation.}
No proof block was collected for this item.

\end{formaltheorem}

\begin{formaltheorem}[EasyCrypt line 593]
\noindent\textbf{EasyCrypt name.}
\begin{lstlisting}[language=EasyCrypt,basicstyle=\ttfamily\scriptsize]
honest_signing_programming_site_queryE
\end{lstlisting}
\noindent\textbf{Checked EasyCrypt statement.}
\begin{lstlisting}[language=EasyCrypt,basicstyle=\ttfamily\scriptsize]
lemma honest_signing_programming_site_queryE md sk m ctx coins :
  honest_signing_programming_site_query md sk m ctx coins =
  challenge_hash_query md
    (commitment_highbits md sk m ctx coins)
    (commitment_lowbits md sk m ctx coins)
    (message_hash (public_key_of_secret md sk) ctx m).
\end{lstlisting}
\noindent\textbf{Formal mathematical reading.}
Mathematical theorem. This is an equality or definitional expansion used for rewriting and normalization.

\noindent\textbf{Role in the proof.}
Use in this chapter. It is a reusable checked theorem cited by later proof scripts.

\noindent\textbf{Proof-script interpretation.}
Proof-script reading. The machine proof decomposes the theorem into rewrites, program-logic obligations, relational equivalences, and arithmetic side conditions.
\emph{Main tactics visible in the proof script:} \ec{rewrite}.
\emph{Named identifiers syntactically cited in the proof block:}
\begin{lstlisting}[language=EasyCrypt,basicstyle=\ttfamily\scriptsize]
actual_signature_programming_site
honest_signing_programming_site
honest_signing_programming_site_query
programming_site_query
sig_commitment_highbits
sig_commitment_lowbits
sig_message_hash
sign_internal
signature_programming_site_of_transcript
transcript_from_honest_signing
transcript_of_signature
\end{lstlisting}

\end{formaltheorem}

\begin{formaltheorem}[EasyCrypt line 617]
\noindent\textbf{EasyCrypt name.}
\begin{lstlisting}[language=EasyCrypt,basicstyle=\ttfamily\scriptsize]
honest_signing_programming_site_query_entropy_token
\end{lstlisting}
\noindent\textbf{Checked EasyCrypt statement.}
\begin{lstlisting}[language=EasyCrypt,basicstyle=\ttfamily\scriptsize]
lemma honest_signing_programming_site_query_entropy_token
  md sk m ctx coins :
  programmed_site_entropy_token_from_query
    (honest_signing_programming_site_query md sk m ctx coins) =
  signing_entropy_token_of_coins coins.
\end{lstlisting}
\noindent\textbf{Formal mathematical reading.}
Mathematical theorem. This is an equality or definitional expansion used for rewriting and normalization.

\noindent\textbf{Role in the proof.}
Use in this chapter. It is a reusable checked theorem cited by later proof scripts.

\noindent\textbf{Proof-script interpretation.}
Proof-script reading. The machine proof decomposes the theorem into rewrites, program-logic obligations, relational equivalences, and arithmetic side conditions.
\emph{Main tactics visible in the proof script:} \ec{rewrite}, \ec{apply}.
\emph{Named identifiers syntactically cited in the proof block:}
\begin{lstlisting}[language=EasyCrypt,basicstyle=\ttfamily\scriptsize]
challenge_hash_query
commitment_lowbits_entropy_token
honest_signing_programming_site_queryE
programmed_site_entropy_token_from_query
\end{lstlisting}

\end{formaltheorem}

\begin{formaltheorem}[EasyCrypt line 629]
\noindent\textbf{EasyCrypt name.}
\begin{lstlisting}[language=EasyCrypt,basicstyle=\ttfamily\scriptsize]
honest_signing_programming_site_query_token_injective
\end{lstlisting}
\noindent\textbf{Checked EasyCrypt statement.}
\begin{lstlisting}[language=EasyCrypt,basicstyle=\ttfamily\scriptsize]
lemma honest_signing_programming_site_query_token_injective
  md sk m ctx x q :
  signing_entropy_token_in_range x =>
  honest_signing_programming_site_query md sk m ctx
    (signing_entropy_token_to_coins x) = q =>
  x = programmed_site_entropy_token_from_query q.
\end{lstlisting}
\noindent\textbf{Formal mathematical reading.}
Mathematical theorem. This is an implication, normally turning one invariant, validity predicate, or proof obligation into another.

\noindent\textbf{Role in the proof.}
Use in this chapter. It is a reusable checked theorem cited by later proof scripts.

\noindent\textbf{Proof-script interpretation.}
Proof-script reading. The machine proof decomposes the theorem into rewrites, program-logic obligations, relational equivalences, and arithmetic side conditions.
\emph{Main tactics visible in the proof script:} \ec{have}, \ec{rewrite}.
\emph{Named identifiers syntactically cited in the proof block:}
\begin{lstlisting}[language=EasyCrypt,basicstyle=\ttfamily\scriptsize]
ctx
honest_signing_programming_site_query_entropy_token
md
query_eq
signing_entropy_token_to_coins
signing_entropy_token_to_coins_tokenK
sk
xrange
\end{lstlisting}

\end{formaltheorem}

\begin{formaltheorem}[EasyCrypt line 644]
\noindent\textbf{EasyCrypt name.}
\begin{lstlisting}[language=EasyCrypt,basicstyle=\ttfamily\scriptsize]
signing_distribution_programming_site_query_spread
\end{lstlisting}
\noindent\textbf{Checked EasyCrypt statement.}
\begin{lstlisting}[language=EasyCrypt,basicstyle=\ttfamily\scriptsize]
lemma signing_distribution_programming_site_query_spread
  (d : random_coins distr) md sk m ctx q :
  signing_randomness_token_spread d =>
  mu d
    (fun coins =>
      honest_signing_programming_site_query md sk m ctx coins = q) <=
  1%r / challenge_support_cardinality_lower_bound.
\end{lstlisting}
\noindent\textbf{Formal mathematical reading.}
Mathematical theorem. This is an inequality, usually an arithmetic loss bound, event inclusion bound, or monotonicity fact used to compose probabilities.

\noindent\textbf{Role in the proof.}
Use in this chapter. It is a reusable checked theorem cited by later proof scripts.

\noindent\textbf{Proof-script interpretation.}
Proof-script reading. The machine proof decomposes the theorem into rewrites, program-logic obligations, relational equivalences, and arithmetic side conditions.
\emph{Main tactics visible in the proof script:} \ec{apply}, \ec{have}, \ec{rewrite}, \ec{smt}.
\emph{Named identifiers syntactically cited in the proof block:}
\begin{lstlisting}[language=EasyCrypt,basicstyle=\ttfamily\scriptsize]
challenge_support_cardinality_lower_bound
coins
ctx
honest_signing_programming_site_query_entropy_token
ler_trans
md
mu
mu_le
programmed_site_entropy_token_from_query
query_eq
signing_entropy_token_cardinality
signing_entropy_token_of_coins
signing_randomness_point_bound
signing_randomness_token_spread
sk
token_spread
\end{lstlisting}

\end{formaltheorem}

\begin{formaltheorem}[EasyCrypt line 673]
\noindent\textbf{EasyCrypt name.}
\begin{lstlisting}[language=EasyCrypt,basicstyle=\ttfamily\scriptsize]
honest_signing_programming_site_query_spread
\end{lstlisting}
\noindent\textbf{Checked EasyCrypt statement.}
\begin{lstlisting}[language=EasyCrypt,basicstyle=\ttfamily\scriptsize]
lemma honest_signing_programming_site_query_spread
  md sk m ctx q :
  mu drandom_coins
    (fun coins =>
      honest_signing_programming_site_query md sk m ctx coins = q) <=
  1%r / challenge_support_cardinality_lower_bound.
\end{lstlisting}
\noindent\textbf{Formal mathematical reading.}
Mathematical theorem. This is an inequality, usually an arithmetic loss bound, event inclusion bound, or monotonicity fact used to compose probabilities.

\noindent\textbf{Role in the proof.}
Use in this chapter. It is a reusable checked theorem cited by later proof scripts.

\noindent\textbf{Proof-script interpretation.}
Proof-script reading. The machine proof decomposes the theorem into rewrites, program-logic obligations, relational equivalences, and arithmetic side conditions.
\emph{Main tactics visible in the proof script:} \ec{apply}.
\emph{Named identifiers syntactically cited in the proof block:}
\begin{lstlisting}[language=EasyCrypt,basicstyle=\ttfamily\scriptsize]
signing_coin_distribution_token_spread
signing_distribution_programming_site_query_spread
\end{lstlisting}

\end{formaltheorem}

\begin{formaltheorem}[EasyCrypt line 684]
\noindent\textbf{EasyCrypt name.}
\begin{lstlisting}[language=EasyCrypt,basicstyle=\ttfamily\scriptsize]
honest_signing_programming_site_query_spread_bound
\end{lstlisting}
\noindent\textbf{Checked EasyCrypt statement.}
\begin{lstlisting}[language=EasyCrypt,basicstyle=\ttfamily\scriptsize]
lemma honest_signing_programming_site_query_spread_bound
  md sk m ctx qs :
  programmed_site_query_min_entropy_bound md sk m ctx qs
    (1%r / challenge_support_cardinality_lower_bound).
\end{lstlisting}
\noindent\textbf{Formal mathematical reading.}
Mathematical theorem. This lemma is a universally quantified proposition over the variables shown in the EasyCrypt statement.

\noindent\textbf{Role in the proof.}
Use in this chapter. It contributes directly to an advantage bound, bad-event bound, or real-arithmetic composition step.

\noindent\textbf{Proof-script interpretation.}
Proof-script reading. The machine proof decomposes the theorem into rewrites, program-logic obligations, relational equivalences, and arithmetic side conditions.
\emph{Main tactics visible in the proof script:} \ec{rewrite}, \ec{apply}.
\emph{Named identifiers syntactically cited in the proof block:}
\begin{lstlisting}[language=EasyCrypt,basicstyle=\ttfamily\scriptsize]
honest_signing_programming_site_query_spread
programmed_site_query_min_entropy_bound
\end{lstlisting}

\end{formaltheorem}

\begin{formaltheorem}[EasyCrypt line 694]
\noindent\textbf{EasyCrypt name.}
\begin{lstlisting}[language=EasyCrypt,basicstyle=\ttfamily\scriptsize]
honest_signing_programming_site_query_spread_bound_for
\end{lstlisting}
\noindent\textbf{Checked EasyCrypt statement.}
\begin{lstlisting}[language=EasyCrypt,basicstyle=\ttfamily\scriptsize]
lemma honest_signing_programming_site_query_spread_bound_for
  (d : random_coins distr) md sk m ctx qs :
  signing_randomness_token_spread d =>
  programmed_site_query_min_entropy_bound_for d md sk m ctx qs
    (1%r / challenge_support_cardinality_lower_bound).
\end{lstlisting}
\noindent\textbf{Formal mathematical reading.}
Mathematical theorem. This is an implication, normally turning one invariant, validity predicate, or proof obligation into another.

\noindent\textbf{Role in the proof.}
Use in this chapter. It contributes directly to an advantage bound, bad-event bound, or real-arithmetic composition step.

\noindent\textbf{Proof-script interpretation.}
Proof-script reading. The machine proof decomposes the theorem into rewrites, program-logic obligations, relational equivalences, and arithmetic side conditions.
\emph{Main tactics visible in the proof script:} \ec{apply}.
\emph{Named identifiers syntactically cited in the proof block:}
\begin{lstlisting}[language=EasyCrypt,basicstyle=\ttfamily\scriptsize]
signing_distribution_programming_site_query_spread
spread
\end{lstlisting}

\end{formaltheorem}

\begin{formaltheorem}[EasyCrypt line 704]
\noindent\textbf{EasyCrypt name.}
\begin{lstlisting}[language=EasyCrypt,basicstyle=\ttfamily\scriptsize]
honest_signing_programming_site_query_is_challenge
\end{lstlisting}
\noindent\textbf{Checked EasyCrypt statement.}
\begin{lstlisting}[language=EasyCrypt,basicstyle=\ttfamily\scriptsize]
lemma honest_signing_programming_site_query_is_challenge
  md sk m ctx coins :
  ro_query_is_challenge
    (honest_signing_programming_site_query md sk m ctx coins).
\end{lstlisting}
\noindent\textbf{Formal mathematical reading.}
Mathematical theorem. This lemma is a universally quantified proposition over the variables shown in the EasyCrypt statement.

\noindent\textbf{Role in the proof.}
Use in this chapter. It contributes directly to an advantage bound, bad-event bound, or real-arithmetic composition step.

\noindent\textbf{Proof-script interpretation.}
Proof-script reading. The machine proof decomposes the theorem into rewrites, program-logic obligations, relational equivalences, and arithmetic side conditions.
\emph{Main tactics visible in the proof script:} \ec{rewrite}.
\emph{Named identifiers syntactically cited in the proof block:}
\begin{lstlisting}[language=EasyCrypt,basicstyle=\ttfamily\scriptsize]
challenge_hash_query
honest_signing_programming_site_queryE
ro_query_is_challenge
\end{lstlisting}

\end{formaltheorem}

\begin{formaltheorem}[EasyCrypt line 713]
\noindent\textbf{EasyCrypt name.}
\begin{lstlisting}[language=EasyCrypt,basicstyle=\ttfamily\scriptsize]
programming_site_prequeried_honest_challenge_filter
\end{lstlisting}
\noindent\textbf{Checked EasyCrypt statement.}
\begin{lstlisting}[language=EasyCrypt,basicstyle=\ttfamily\scriptsize]
lemma programming_site_prequeried_honest_challenge_filter
  qs md sk m ctx coins :
  programming_site_prequeried qs
    (honest_signing_programming_site md sk m ctx coins) =
  programming_site_prequeried (filter ro_query_is_challenge qs)
    (honest_signing_programming_site md sk m ctx coins).
\end{lstlisting}
\noindent\textbf{Formal mathematical reading.}
Mathematical theorem. This is an equality or definitional expansion used for rewriting and normalization.

\noindent\textbf{Role in the proof.}
Use in this chapter. It contributes directly to an advantage bound, bad-event bound, or real-arithmetic composition step.

\noindent\textbf{Proof-script interpretation.}
Proof-script reading. The machine proof decomposes the theorem into rewrites, program-logic obligations, relational equivalences, and arithmetic side conditions.
\emph{Main tactics visible in the proof script:} \ec{rewrite}.
\emph{Named identifiers syntactically cited in the proof block:}
\begin{lstlisting}[language=EasyCrypt,basicstyle=\ttfamily\scriptsize]
honest_signing_programming_site_query
honest_signing_programming_site_query_is_challenge
mem_filter
programming_site_prequeried
\end{lstlisting}

\end{formaltheorem}

\begin{formaltheorem}[EasyCrypt line 726]
\noindent\textbf{EasyCrypt name.}
\begin{lstlisting}[language=EasyCrypt,basicstyle=\ttfamily\scriptsize]
programming_site_prequeried_honest_message_consE
\end{lstlisting}
\noindent\textbf{Checked EasyCrypt statement.}
\begin{lstlisting}[language=EasyCrypt,basicstyle=\ttfamily\scriptsize]
lemma programming_site_prequeried_honest_message_consE
  pk ctx m qs md sk sm sctx coins :
  programming_site_prequeried (message_hash_query pk ctx m :: qs)
    (honest_signing_programming_site md sk sm sctx coins) =
  programming_site_prequeried qs
    (honest_signing_programming_site md sk sm sctx coins).
\end{lstlisting}
\noindent\textbf{Formal mathematical reading.}
Mathematical theorem. This is an equality or definitional expansion used for rewriting and normalization.

\noindent\textbf{Role in the proof.}
Use in this chapter. It is a reusable checked theorem cited by later proof scripts.

\noindent\textbf{Proof-script interpretation.}
Proof-script reading. The machine proof decomposes the theorem into rewrites, program-logic obligations, relational equivalences, and arithmetic side conditions.
\emph{Main tactics visible in the proof script:} \ec{rewrite}, \ec{smt}.
\emph{Named identifiers syntactically cited in the proof block:}
\begin{lstlisting}[language=EasyCrypt,basicstyle=\ttfamily\scriptsize]
actual_signature_programming_site
challenge_hash_query
honest_signing_programming_site
message_hash_query
programming_site_prequeried
programming_site_query
signature_programming_site_of_transcript
transcript_from_honest_signing
transcript_of_signature
transcript_signature_challenge_query
\end{lstlisting}

\end{formaltheorem}

\begin{formaltheorem}[EasyCrypt line 745]
\noindent\textbf{EasyCrypt name.}
\begin{lstlisting}[language=EasyCrypt,basicstyle=\ttfamily\scriptsize]
honest_signing_programming_site_outputE
\end{lstlisting}
\noindent\textbf{Checked EasyCrypt statement.}
\begin{lstlisting}[language=EasyCrypt,basicstyle=\ttfamily\scriptsize]
lemma honest_signing_programming_site_outputE md sk m ctx coins :
  programming_site_output
    (honest_signing_programming_site md sk m ctx coins) =
  ro_output_of_challenge
    (challenge_hash md
      (commitment_highbits md sk m ctx coins)
      (commitment_lowbits md sk m ctx coins)
      (message_hash (public_key_of_secret md sk) ctx m)).
\end{lstlisting}
\noindent\textbf{Formal mathematical reading.}
Mathematical theorem. This is an equality or definitional expansion used for rewriting and normalization.

\noindent\textbf{Role in the proof.}
Use in this chapter. It is a reusable checked theorem cited by later proof scripts.

\noindent\textbf{Proof-script interpretation.}
Proof-script reading. The machine proof decomposes the theorem into rewrites, program-logic obligations, relational equivalences, and arithmetic side conditions.
\emph{Main tactics visible in the proof script:} \ec{rewrite}.
\emph{Named identifiers syntactically cited in the proof block:}
\begin{lstlisting}[language=EasyCrypt,basicstyle=\ttfamily\scriptsize]
actual_signature_programming_site
honest_signing_programming_site
programming_site_output
sig_challenge
sign_internal
signature_programming_site_of_transcript
transcript_from_honest_signing
transcript_of_signature
transcript_signature_challenge_output
\end{lstlisting}

\end{formaltheorem}

\begin{formaltheorem}[EasyCrypt line 764]
\noindent\textbf{EasyCrypt name.}
\begin{lstlisting}[language=EasyCrypt,basicstyle=\ttfamily\scriptsize]
honest_signing_programming_site_functional_output
\end{lstlisting}
\noindent\textbf{Checked EasyCrypt statement.}
\begin{lstlisting}[language=EasyCrypt,basicstyle=\ttfamily\scriptsize]
lemma honest_signing_programming_site_functional_output
  md sk1 m1 ctx1 coins1 sk2 m2 ctx2 coins2 :
  honest_signing_programming_site_query md sk1 m1 ctx1 coins1 =
  honest_signing_programming_site_query md sk2 m2 ctx2 coins2 =>
  programming_site_output
    (honest_signing_programming_site md sk1 m1 ctx1 coins1) =
  programming_site_output
    (honest_signing_programming_site md sk2 m2 ctx2 coins2).
\end{lstlisting}
\noindent\textbf{Formal mathematical reading.}
Mathematical theorem. This is an implication, normally turning one invariant, validity predicate, or proof obligation into another.

\noindent\textbf{Role in the proof.}
Use in this chapter. It is a reusable checked theorem cited by later proof scripts.

\noindent\textbf{Proof-script interpretation.}
Proof-script reading. The machine proof decomposes the theorem into rewrites, program-logic obligations, relational equivalences, and arithmetic side conditions.
\emph{Main tactics visible in the proof script:} \ec{rewrite}, \ec{smt}.
\emph{Named identifiers syntactically cited in the proof block:}
\begin{lstlisting}[language=EasyCrypt,basicstyle=\ttfamily\scriptsize]
honest_signing_programming_site_outputE
honest_signing_programming_site_queryE
\end{lstlisting}

\end{formaltheorem}

\begin{formaltheorem}[EasyCrypt line 778]
\noindent\textbf{EasyCrypt name.}
\begin{lstlisting}[language=EasyCrypt,basicstyle=\ttfamily\scriptsize]
honest_signing_programming_sites_no_conflict
\end{lstlisting}
\noindent\textbf{Checked EasyCrypt statement.}
\begin{lstlisting}[language=EasyCrypt,basicstyle=\ttfamily\scriptsize]
lemma honest_signing_programming_sites_no_conflict
  md sk1 m1 ctx1 coins1 sk2 m2 ctx2 coins2 :
  ! programming_sites_conflict
      (honest_signing_programming_site md sk1 m1 ctx1 coins1)
      (honest_signing_programming_site md sk2 m2 ctx2 coins2).
\end{lstlisting}
\noindent\textbf{Formal mathematical reading.}
Mathematical theorem. This lemma is a universally quantified proposition over the variables shown in the EasyCrypt statement.

\noindent\textbf{Role in the proof.}
Use in this chapter. It is a reusable checked theorem cited by later proof scripts.

\noindent\textbf{Proof-script interpretation.}
Proof-script reading. The machine proof decomposes the theorem into rewrites, program-logic obligations, relational equivalences, and arithmetic side conditions.
\emph{Main tactics visible in the proof script:} \ec{rewrite}, \ec{case}, \ec{apply}.
\emph{Named identifiers syntactically cited in the proof block:}
\begin{lstlisting}[language=EasyCrypt,basicstyle=\ttfamily\scriptsize]
coins1
coins2
ctx1
ctx2
honest_signing_programming_site
honest_signing_programming_site_functional_output
m1
m2
md
programming_site_query
programming_sites_conflict
programming_sites_same_query
query_eq
sk1
sk2
\end{lstlisting}

\end{formaltheorem}

\begin{formaltheorem}[EasyCrypt line 795]
\noindent\textbf{EasyCrypt name.}
\begin{lstlisting}[language=EasyCrypt,basicstyle=\ttfamily\scriptsize]
actual_signature_programming_site_sign_internalE
\end{lstlisting}
\noindent\textbf{Checked EasyCrypt statement.}
\begin{lstlisting}[language=EasyCrypt,basicstyle=\ttfamily\scriptsize]
lemma actual_signature_programming_site_sign_internalE
  md pk sk m ctx coins :
  pk = public_key_of_secret md sk =>
  actual_signature_programming_site md
    (transcript_of_signature md pk m ctx
      (sign_internal md sk m ctx coins)) =
  honest_signing_programming_site md sk m ctx coins.
\end{lstlisting}
\noindent\textbf{Formal mathematical reading.}
Mathematical theorem. This is an implication, normally turning one invariant, validity predicate, or proof obligation into another.

\noindent\textbf{Role in the proof.}
Use in this chapter. It contributes directly to an advantage bound, bad-event bound, or real-arithmetic composition step.

\noindent\textbf{Proof-script interpretation.}
Proof-script reading. The machine proof decomposes the theorem into rewrites, program-logic obligations, relational equivalences, and arithmetic side conditions.
\emph{Main tactics visible in the proof script:} \ec{rewrite}.
\emph{Named identifiers syntactically cited in the proof block:}
\begin{lstlisting}[language=EasyCrypt,basicstyle=\ttfamily\scriptsize]
honest_signing_programming_site
pkE
transcript_from_honest_signing
\end{lstlisting}

\end{formaltheorem}

\begin{formaltheorem}[EasyCrypt line 808]
\noindent\textbf{EasyCrypt name.}
\begin{lstlisting}[language=EasyCrypt,basicstyle=\ttfamily\scriptsize]
actual_signature_programming_site_sign_internal_any_pkE
\end{lstlisting}
\noindent\textbf{Checked EasyCrypt statement.}
\begin{lstlisting}[language=EasyCrypt,basicstyle=\ttfamily\scriptsize]
lemma actual_signature_programming_site_sign_internal_any_pkE
  md pk sk m ctx coins :
  actual_signature_programming_site md
    (transcript_of_signature md pk m ctx
      (sign_internal md sk m ctx coins)) =
  honest_signing_programming_site md sk m ctx coins.
\end{lstlisting}
\noindent\textbf{Formal mathematical reading.}
Mathematical theorem. This is an equality or definitional expansion used for rewriting and normalization.

\noindent\textbf{Role in the proof.}
Use in this chapter. It is a reusable checked theorem cited by later proof scripts.

\noindent\textbf{Proof-script interpretation.}
Proof-script reading. The machine proof decomposes the theorem into rewrites, program-logic obligations, relational equivalences, and arithmetic side conditions.
\emph{Main tactics visible in the proof script:} \ec{rewrite}.
\emph{Named identifiers syntactically cited in the proof block:}
\begin{lstlisting}[language=EasyCrypt,basicstyle=\ttfamily\scriptsize]
actual_signature_programming_site
honest_signing_programming_site
signature_programming_site_of_transcript
transcript_from_honest_signing
transcript_of_signature
transcript_signature_challenge_output
transcript_signature_challenge_query
\end{lstlisting}

\end{formaltheorem}

\begin{formaltheorem}[EasyCrypt line 823]
\noindent\textbf{EasyCrypt name.}
\begin{lstlisting}[language=EasyCrypt,basicstyle=\ttfamily\scriptsize]
actual_signature_programming_site_prequeried_after_cons
\end{lstlisting}
\noindent\textbf{Checked EasyCrypt statement.}
\begin{lstlisting}[language=EasyCrypt,basicstyle=\ttfamily\scriptsize]
lemma actual_signature_programming_site_prequeried_after_cons
  md pk sk m ctx coins q qs :
  pk = public_key_of_secret md sk =>
  programming_site_prequeried qs
    (honest_signing_programming_site md sk m ctx coins) =>
  programming_site_prequeried (q :: qs)
    (actual_signature_programming_site md
      (transcript_of_signature md pk m ctx
        (sign_internal md sk m ctx coins))).
\end{lstlisting}
\noindent\textbf{Formal mathematical reading.}
Mathematical theorem. This is an implication, normally turning one invariant, validity predicate, or proof obligation into another.

\noindent\textbf{Role in the proof.}
Use in this chapter. It is a reusable checked theorem cited by later proof scripts.

\noindent\textbf{Proof-script interpretation.}
Proof-script reading. The machine proof decomposes the theorem into rewrites, program-logic obligations, relational equivalences, and arithmetic side conditions.
\emph{Main tactics visible in the proof script:} \ec{rewrite}, \ec{apply}.
\emph{Named identifiers syntactically cited in the proof block:}
\begin{lstlisting}[language=EasyCrypt,basicstyle=\ttfamily\scriptsize]
actual_signature_programming_site_sign_internalE
coins
ctx
md
pk
pkE
pre
programming_site_prequeried_cons_preserve
sk
\end{lstlisting}

\end{formaltheorem}

\begin{formaltheorem}[EasyCrypt line 839]
\noindent\textbf{EasyCrypt name.}
\begin{lstlisting}[language=EasyCrypt,basicstyle=\ttfamily\scriptsize]
actual_signature_programming_site_prequeried_after_message_hash
\end{lstlisting}
\noindent\textbf{Checked EasyCrypt statement.}
\begin{lstlisting}[language=EasyCrypt,basicstyle=\ttfamily\scriptsize]
lemma actual_signature_programming_site_prequeried_after_message_hash
  md pk sk m ctx coins qpk qctx qm qs :
  pk = public_key_of_secret md sk =>
  programming_site_prequeried qs
    (honest_signing_programming_site md sk m ctx coins) =>
  programming_site_prequeried (message_hash_query qpk qctx qm :: qs)
    (actual_signature_programming_site md
      (transcript_of_signature md pk m ctx
        (sign_internal md sk m ctx coins))).
\end{lstlisting}
\noindent\textbf{Formal mathematical reading.}
Mathematical theorem. This is an implication, normally turning one invariant, validity predicate, or proof obligation into another.

\noindent\textbf{Role in the proof.}
Use in this chapter. It is a reusable checked theorem cited by later proof scripts.

\noindent\textbf{Proof-script interpretation.}
Proof-script reading. The machine proof decomposes the theorem into rewrites, program-logic obligations, relational equivalences, and arithmetic side conditions.
\emph{Main tactics visible in the proof script:} \ec{apply}.
\emph{Named identifiers syntactically cited in the proof block:}
\begin{lstlisting}[language=EasyCrypt,basicstyle=\ttfamily\scriptsize]
actual_signature_programming_site_prequeried_after_cons
coins
ctx
md
message_hash_query
pk
pkE
pre
qctx
qm
qpk
qs
sk
\end{lstlisting}

\end{formaltheorem}

\begin{formaltheorem}[EasyCrypt line 854]
\noindent\textbf{EasyCrypt name.}
\begin{lstlisting}[language=EasyCrypt,basicstyle=\ttfamily\scriptsize]
actual_signature_programming_site_prequeried_message_hashE
\end{lstlisting}
\noindent\textbf{Checked EasyCrypt statement.}
\begin{lstlisting}[language=EasyCrypt,basicstyle=\ttfamily\scriptsize]
lemma actual_signature_programming_site_prequeried_message_hashE
  md pk sk m ctx coins qpk qctx qm qs :
  pk = public_key_of_secret md sk =>
  programming_site_prequeried (message_hash_query qpk qctx qm :: qs)
    (actual_signature_programming_site md
      (transcript_of_signature md pk m ctx
        (sign_internal md sk m ctx coins))) =
  programming_site_prequeried qs
    (honest_signing_programming_site md sk m ctx coins).
\end{lstlisting}
\noindent\textbf{Formal mathematical reading.}
Mathematical theorem. This is an implication, normally turning one invariant, validity predicate, or proof obligation into another.

\noindent\textbf{Role in the proof.}
Use in this chapter. It is a reusable checked theorem cited by later proof scripts.

\noindent\textbf{Proof-script interpretation.}
Proof-script reading. The machine proof decomposes the theorem into rewrites, program-logic obligations, relational equivalences, and arithmetic side conditions.
\emph{Main tactics visible in the proof script:} \ec{rewrite}.
\emph{Named identifiers syntactically cited in the proof block:}
\begin{lstlisting}[language=EasyCrypt,basicstyle=\ttfamily\scriptsize]
actual_signature_programming_site_sign_internalE
coins
ctx
md
pk
pkE
programming_site_prequeried_honest_message_consE
sk
\end{lstlisting}

\end{formaltheorem}

\begin{formaltheorem}[EasyCrypt line 870]
\noindent\textbf{EasyCrypt name.}
\begin{lstlisting}[language=EasyCrypt,basicstyle=\ttfamily\scriptsize]
actual_signature_programming_site_prequeried_message_hash_any_pkE
\end{lstlisting}
\noindent\textbf{Checked EasyCrypt statement.}
\begin{lstlisting}[language=EasyCrypt,basicstyle=\ttfamily\scriptsize]
lemma actual_signature_programming_site_prequeried_message_hash_any_pkE
  md pk sk m ctx coins qpk qctx qm qs :
  programming_site_prequeried (message_hash_query qpk qctx qm :: qs)
    (actual_signature_programming_site md
      (transcript_of_signature md pk m ctx
        (sign_internal md sk m ctx coins))) =
  programming_site_prequeried qs
    (honest_signing_programming_site md sk m ctx coins).
\end{lstlisting}
\noindent\textbf{Formal mathematical reading.}
Mathematical theorem. This is an equality or definitional expansion used for rewriting and normalization.

\noindent\textbf{Role in the proof.}
Use in this chapter. It is a reusable checked theorem cited by later proof scripts.

\noindent\textbf{Proof-script interpretation.}
Proof-script reading. The machine proof decomposes the theorem into rewrites, program-logic obligations, relational equivalences, and arithmetic side conditions.
\emph{Main tactics visible in the proof script:} \ec{rewrite}.
\emph{Named identifiers syntactically cited in the proof block:}
\begin{lstlisting}[language=EasyCrypt,basicstyle=\ttfamily\scriptsize]
actual_signature_programming_site_sign_internal_any_pkE
coins
ctx
md
pk
programming_site_prequeried_honest_message_consE
sk
\end{lstlisting}

\end{formaltheorem}

\begin{formaltheorem}[EasyCrypt line 894]
\noindent\textbf{EasyCrypt name.}
\begin{lstlisting}[language=EasyCrypt,basicstyle=\ttfamily\scriptsize]
programming_site_log_conflict_free_for_honest_nil
\end{lstlisting}
\noindent\textbf{Checked EasyCrypt statement.}
\begin{lstlisting}[language=EasyCrypt,basicstyle=\ttfamily\scriptsize]
lemma programming_site_log_conflict_free_for_honest_nil md sk :
  programming_site_log_conflict_free_for_honest md sk [].
\end{lstlisting}
\noindent\textbf{Formal mathematical reading.}
Mathematical theorem. This lemma is a universally quantified proposition over the variables shown in the EasyCrypt statement.

\noindent\textbf{Role in the proof.}
Use in this chapter. It is a reusable checked theorem cited by later proof scripts.

\noindent\textbf{Proof-script interpretation.}
No proof block was collected for this item.

\end{formaltheorem}

\begin{formaltheorem}[EasyCrypt line 900]
\noindent\textbf{EasyCrypt name.}
\begin{lstlisting}[language=EasyCrypt,basicstyle=\ttfamily\scriptsize]
programming_site_log_conflict_free_for_honest_cons
\end{lstlisting}
\noindent\textbf{Checked EasyCrypt statement.}
\begin{lstlisting}[language=EasyCrypt,basicstyle=\ttfamily\scriptsize]
lemma programming_site_log_conflict_free_for_honest_cons
  md sk m ctx coins sites :
  programming_site_log_conflict_free_for_honest md sk sites =>
  programming_site_log_conflict_free_for_honest md sk
    (honest_signing_programming_site md sk m ctx coins :: sites).
\end{lstlisting}
\noindent\textbf{Formal mathematical reading.}
Mathematical theorem. This is an implication, normally turning one invariant, validity predicate, or proof obligation into another.

\noindent\textbf{Role in the proof.}
Use in this chapter. It is a reusable checked theorem cited by later proof scripts.

\noindent\textbf{Proof-script interpretation.}
Proof-script reading. The machine proof decomposes the theorem into rewrites, program-logic obligations, relational equivalences, and arithmetic side conditions.
\emph{Main tactics visible in the proof script:} \ec{rewrite}, \ec{split}, \ec{apply}.
\emph{Named identifiers syntactically cited in the proof block:}
\begin{lstlisting}[language=EasyCrypt,basicstyle=\ttfamily\scriptsize]
coins
ctx
honest_signing_programming_sites_no_conflict
log
programming_site_conflict_free
programming_site_log_conflict_free_for_honest
\end{lstlisting}

\end{formaltheorem}

\begin{formaltheorem}[EasyCrypt line 914]
\noindent\textbf{EasyCrypt name.}
\begin{lstlisting}[language=EasyCrypt,basicstyle=\ttfamily\scriptsize]
programming_site_conflict_free_no_reprograms
\end{lstlisting}
\noindent\textbf{Checked EasyCrypt statement.}
\begin{lstlisting}[language=EasyCrypt,basicstyle=\ttfamily\scriptsize]
lemma programming_site_conflict_free_no_reprograms site sites :
  programming_site_conflict_free site sites =>
  ! programming_site_reprograms sites site.
\end{lstlisting}
\noindent\textbf{Formal mathematical reading.}
Mathematical theorem. This is an implication, normally turning one invariant, validity predicate, or proof obligation into another.

\noindent\textbf{Role in the proof.}
Use in this chapter. It is a reusable checked theorem cited by later proof scripts.

\noindent\textbf{Proof-script interpretation.}
Proof-script reading. The machine proof decomposes the theorem into rewrites, program-logic obligations, relational equivalences, and arithmetic side conditions.
\emph{Main tactics visible in the proof script:} \ec{rewrite}, \ec{apply}.
\emph{Named identifiers syntactically cited in the proof block:}
\begin{lstlisting}[language=EasyCrypt,basicstyle=\ttfamily\scriptsize]
allP
hasPn
no_conflict
old
old_mem
programming_site_conflict_free
programming_site_reprograms
\end{lstlisting}

\end{formaltheorem}

\begin{formaltheorem}[EasyCrypt line 924]
\noindent\textbf{EasyCrypt name.}
\begin{lstlisting}[language=EasyCrypt,basicstyle=\ttfamily\scriptsize]
actual_signature_programming_site_sign_internal_conflict_step
\end{lstlisting}
\noindent\textbf{Checked EasyCrypt statement.}
\begin{lstlisting}[language=EasyCrypt,basicstyle=\ttfamily\scriptsize]
lemma actual_signature_programming_site_sign_internal_conflict_step
  md pk sk m ctx coins sites :
  pk = public_key_of_secret md sk =>
  programming_site_log_conflict_free_for_honest md sk sites =>
  ! programming_site_reprograms sites
      (actual_signature_programming_site md
        (transcript_of_signature md pk m ctx
          (sign_internal md sk m ctx coins))) /\
  programming_site_log_conflict_free_for_honest md sk
    (actual_signature_programming_site md
      (transcript_of_signature md pk m ctx
        (sign_internal md sk m ctx coins)) :: sites).
\end{lstlisting}
\noindent\textbf{Formal mathematical reading.}
Mathematical theorem. This is an implication, normally turning one invariant, validity predicate, or proof obligation into another.

\noindent\textbf{Role in the proof.}
Use in this chapter. It is a reusable checked theorem cited by later proof scripts.

\noindent\textbf{Proof-script interpretation.}
Proof-script reading. The machine proof decomposes the theorem into rewrites, program-logic obligations, relational equivalences, and arithmetic side conditions.
\emph{Main tactics visible in the proof script:} \ec{have}, \ec{split}, \ec{rewrite}, \ec{apply}.
\emph{Named identifiers syntactically cited in the proof block:}
\begin{lstlisting}[language=EasyCrypt,basicstyle=\ttfamily\scriptsize]
actual_signature_programming_site_sign_internalE
coins
ctx
log
md
pk
pkE
programming_site_conflict_free_no_reprograms
programming_site_log_conflict_free_for_honest_cons
siteE
sk
\end{lstlisting}

\end{formaltheorem}

\begin{formaltheorem}[EasyCrypt line 949]
\noindent\textbf{EasyCrypt name.}
\begin{lstlisting}[language=EasyCrypt,basicstyle=\ttfamily\scriptsize]
programmed_site_reprogram_one_step_conflict_free_zero
\end{lstlisting}
\noindent\textbf{Checked EasyCrypt statement.}
\begin{lstlisting}[language=EasyCrypt,basicstyle=\ttfamily\scriptsize]
lemma programmed_site_reprogram_one_step_conflict_free_zero
  md sk m ctx sites :
  (forall coins,
    programming_site_conflict_free
      (honest_signing_programming_site md sk m ctx coins) sites) =>
  mu drandom_coins
    (fun coins =>
      programming_site_reprograms sites
        (honest_signing_programming_site md sk m ctx coins)) =
  0%r.
\end{lstlisting}
\noindent\textbf{Formal mathematical reading.}
Mathematical theorem. This is an implication, normally turning one invariant, validity predicate, or proof obligation into another.

\noindent\textbf{Role in the proof.}
Use in this chapter. It is a reusable checked theorem cited by later proof scripts.

\noindent\textbf{Proof-script interpretation.}
Proof-script reading. The machine proof decomposes the theorem into rewrites, program-logic obligations, relational equivalences, and arithmetic side conditions.
\emph{Main tactics visible in the proof script:} \ec{rewrite}.
\emph{Named identifiers syntactically cited in the proof block:}
\begin{lstlisting}[language=EasyCrypt,basicstyle=\ttfamily\scriptsize]
coins
conflict_free
ctx
honest_signing_programming_site
md
mu0
mu_eq
programming_site_conflict_free_no_reprograms
random_coins
sites
sk
\end{lstlisting}

\end{formaltheorem}

\begin{formaltheorem}[EasyCrypt line 969]
\noindent\textbf{EasyCrypt name.}
\begin{lstlisting}[language=EasyCrypt,basicstyle=\ttfamily\scriptsize]
programmed_site_prequery_one_step_fset_bound_for
\end{lstlisting}
\noindent\textbf{Checked EasyCrypt statement.}
\begin{lstlisting}[language=EasyCrypt,basicstyle=\ttfamily\scriptsize]
lemma programmed_site_prequery_one_step_fset_bound_for
  (d : random_coins distr) md sk m ctx qs bd :
  programmed_site_query_min_entropy_bound_for d md sk m ctx qs bd =>
  mu d
    (fun coins =>
      programming_site_prequeried qs
        (honest_signing_programming_site md sk m ctx coins)) <=
  (card (oflist qs))%r * bd.
\end{lstlisting}
\noindent\textbf{Formal mathematical reading.}
Mathematical theorem. This is an inequality, usually an arithmetic loss bound, event inclusion bound, or monotonicity fact used to compose probabilities.

\noindent\textbf{Role in the proof.}
Use in this chapter. It contributes directly to an advantage bound, bad-event bound, or real-arithmetic composition step.

\noindent\textbf{Proof-script interpretation.}
Proof-script reading. The machine proof decomposes the theorem into rewrites, program-logic obligations, relational equivalences, and arithmetic side conditions.
\emph{Main tactics visible in the proof script:} \ec{rewrite}, \ec{smt}, \ec{apply}.
\emph{Named identifiers syntactically cited in the proof block:}
\begin{lstlisting}[language=EasyCrypt,basicstyle=\ttfamily\scriptsize]
bd
coins
ctx
honest_signing_programming_site
honest_signing_programming_site_query
md
mem_oflist
mu_eq
mu_mem_le_gen
oflist
programmed_site_query_min_entropy_bound_for
programming_site_prequeried
programming_site_query
q_mem
qs
sk
spread
\end{lstlisting}

\end{formaltheorem}

\begin{formaltheorem}[EasyCrypt line 999]
\noindent\textbf{EasyCrypt name.}
\begin{lstlisting}[language=EasyCrypt,basicstyle=\ttfamily\scriptsize]
programmed_site_prequery_one_step_fset_bound
\end{lstlisting}
\noindent\textbf{Checked EasyCrypt statement.}
\begin{lstlisting}[language=EasyCrypt,basicstyle=\ttfamily\scriptsize]
lemma programmed_site_prequery_one_step_fset_bound
  md sk m ctx qs bd :
  programmed_site_query_min_entropy_bound md sk m ctx qs bd =>
  mu drandom_coins
    (fun coins =>
      programming_site_prequeried qs
        (honest_signing_programming_site md sk m ctx coins)) <=
  (card (oflist qs))%r * bd.
\end{lstlisting}
\noindent\textbf{Formal mathematical reading.}
Mathematical theorem. This is an inequality, usually an arithmetic loss bound, event inclusion bound, or monotonicity fact used to compose probabilities.

\noindent\textbf{Role in the proof.}
Use in this chapter. It contributes directly to an advantage bound, bad-event bound, or real-arithmetic composition step.

\noindent\textbf{Proof-script interpretation.}
Proof-script reading. The machine proof decomposes the theorem into rewrites, program-logic obligations, relational equivalences, and arithmetic side conditions.
\emph{Main tactics visible in the proof script:} \ec{rewrite}, \ec{apply}.
\emph{Named identifiers syntactically cited in the proof block:}
\begin{lstlisting}[language=EasyCrypt,basicstyle=\ttfamily\scriptsize]
programmed_site_prequery_one_step_fset_bound_for
programmed_site_query_min_entropy_bound
\end{lstlisting}

\end{formaltheorem}

\begin{formaltheorem}[EasyCrypt line 1012]
\noindent\textbf{EasyCrypt name.}
\begin{lstlisting}[language=EasyCrypt,basicstyle=\ttfamily\scriptsize]
programmed_site_prequery_one_step_budget_bound_for
\end{lstlisting}
\noindent\textbf{Checked EasyCrypt statement.}
\begin{lstlisting}[language=EasyCrypt,basicstyle=\ttfamily\scriptsize]
lemma programmed_site_prequery_one_step_budget_bound_for
  (d : random_coins distr) md sk m ctx qs :
  (card (oflist qs))%r <= rom_hash_query_budget + 1%r =>
  programmed_site_query_min_entropy_bound_for d md sk m ctx qs
    (1%r / challenge_support_cardinality_lower_bound) =>
  mu d
    (fun coins =>
      programming_site_prequeried qs
        (honest_signing_programming_site md sk m ctx coins)) <=
  (rom_hash_query_budget + 1%r) /
    challenge_support_cardinality_lower_bound.
\end{lstlisting}
\noindent\textbf{Formal mathematical reading.}
Mathematical theorem. This is an inequality, usually an arithmetic loss bound, event inclusion bound, or monotonicity fact used to compose probabilities.

\noindent\textbf{Role in the proof.}
Use in this chapter. It contributes directly to an advantage bound, bad-event bound, or real-arithmetic composition step.

\noindent\textbf{Proof-script interpretation.}
Proof-script reading. The machine proof decomposes the theorem into rewrites, program-logic obligations, relational equivalences, and arithmetic side conditions.
\emph{Main tactics visible in the proof script:} \ec{have}, \ec{apply}, \ec{rewrite}, \ec{smt}.
\emph{Named identifiers syntactically cited in the proof block:}
\begin{lstlisting}[language=EasyCrypt,basicstyle=\ttfamily\scriptsize]
card
card_budget
challenge_support_cardinality_lower_bound
challenge_support_cardinality_lower_bound_nonnegative
ctx
inv
inv_nonneg
invr_ge0
ler_trans
ler_wpmul2r
md
oflist
programmed_site_prequery_one_step_fset_bound_for
qs
rom_hash_query_budget
sk
spread
\end{lstlisting}

\end{formaltheorem}

\begin{formaltheorem}[EasyCrypt line 1044]
\noindent\textbf{EasyCrypt name.}
\begin{lstlisting}[language=EasyCrypt,basicstyle=\ttfamily\scriptsize]
programmed_site_prequery_one_step_budget_bound
\end{lstlisting}
\noindent\textbf{Checked EasyCrypt statement.}
\begin{lstlisting}[language=EasyCrypt,basicstyle=\ttfamily\scriptsize]
lemma programmed_site_prequery_one_step_budget_bound
  md sk m ctx qs :
  (card (oflist qs))%r <= rom_hash_query_budget + 1%r =>
  programmed_site_query_min_entropy_bound md sk m ctx qs
    (1%r / challenge_support_cardinality_lower_bound) =>
  mu drandom_coins
    (fun coins =>
      programming_site_prequeried qs
        (honest_signing_programming_site md sk m ctx coins)) <=
  (rom_hash_query_budget + 1%r) /
    challenge_support_cardinality_lower_bound.
\end{lstlisting}
\noindent\textbf{Formal mathematical reading.}
Mathematical theorem. This is an inequality, usually an arithmetic loss bound, event inclusion bound, or monotonicity fact used to compose probabilities.

\noindent\textbf{Role in the proof.}
Use in this chapter. It contributes directly to an advantage bound, bad-event bound, or real-arithmetic composition step.

\noindent\textbf{Proof-script interpretation.}
Proof-script reading. The machine proof decomposes the theorem into rewrites, program-logic obligations, relational equivalences, and arithmetic side conditions.
\emph{Main tactics visible in the proof script:} \ec{rewrite}, \ec{apply}.
\emph{Named identifiers syntactically cited in the proof block:}
\begin{lstlisting}[language=EasyCrypt,basicstyle=\ttfamily\scriptsize]
card_budget
programmed_site_prequery_one_step_budget_bound_for
programmed_site_query_min_entropy_bound
\end{lstlisting}

\end{formaltheorem}

\begin{formaltheorem}[EasyCrypt line 1061]
\noindent\textbf{EasyCrypt name.}
\begin{lstlisting}[language=EasyCrypt,basicstyle=\ttfamily\scriptsize]
programmed_site_prequery_one_step_challenge_fset_bound_for
\end{lstlisting}
\noindent\textbf{Checked EasyCrypt statement.}
\begin{lstlisting}[language=EasyCrypt,basicstyle=\ttfamily\scriptsize]
lemma programmed_site_prequery_one_step_challenge_fset_bound_for
  (d : random_coins distr) md sk m ctx qs bd :
  programmed_site_query_min_entropy_bound_for d md sk m ctx
    (filter ro_query_is_challenge qs) bd =>
  mu d
    (fun coins =>
      programming_site_prequeried qs
        (honest_signing_programming_site md sk m ctx coins)) <=
  (card (oflist (filter ro_query_is_challenge qs)))%r * bd.
\end{lstlisting}
\noindent\textbf{Formal mathematical reading.}
Mathematical theorem. This is an inequality, usually an arithmetic loss bound, event inclusion bound, or monotonicity fact used to compose probabilities.

\noindent\textbf{Role in the proof.}
Use in this chapter. It contributes directly to an advantage bound, bad-event bound, or real-arithmetic composition step.

\noindent\textbf{Proof-script interpretation.}
Proof-script reading. The machine proof decomposes the theorem into rewrites, program-logic obligations, relational equivalences, and arithmetic side conditions.
\emph{Main tactics visible in the proof script:} \ec{rewrite}, \ec{apply}.
\emph{Named identifiers syntactically cited in the proof block:}
\begin{lstlisting}[language=EasyCrypt,basicstyle=\ttfamily\scriptsize]
coins
ctx
filter
honest_signing_programming_site
md
mu_eq
programmed_site_prequery_one_step_fset_bound_for
programming_site_prequeried
programming_site_prequeried_honest_challenge_filter
qs
ro_query_is_challenge
sk
spread
\end{lstlisting}

\end{formaltheorem}

\begin{formaltheorem}[EasyCrypt line 1081]
\noindent\textbf{EasyCrypt name.}
\begin{lstlisting}[language=EasyCrypt,basicstyle=\ttfamily\scriptsize]
programmed_site_prequery_one_step_challenge_fset_bound
\end{lstlisting}
\noindent\textbf{Checked EasyCrypt statement.}
\begin{lstlisting}[language=EasyCrypt,basicstyle=\ttfamily\scriptsize]
lemma programmed_site_prequery_one_step_challenge_fset_bound
  md sk m ctx qs bd :
  programmed_site_query_min_entropy_bound md sk m ctx
    (filter ro_query_is_challenge qs) bd =>
  mu drandom_coins
    (fun coins =>
      programming_site_prequeried qs
        (honest_signing_programming_site md sk m ctx coins)) <=
  (card (oflist (filter ro_query_is_challenge qs)))%r * bd.
\end{lstlisting}
\noindent\textbf{Formal mathematical reading.}
Mathematical theorem. This is an inequality, usually an arithmetic loss bound, event inclusion bound, or monotonicity fact used to compose probabilities.

\noindent\textbf{Role in the proof.}
Use in this chapter. It contributes directly to an advantage bound, bad-event bound, or real-arithmetic composition step.

\noindent\textbf{Proof-script interpretation.}
Proof-script reading. The machine proof decomposes the theorem into rewrites, program-logic obligations, relational equivalences, and arithmetic side conditions.
\emph{Main tactics visible in the proof script:} \ec{rewrite}, \ec{apply}.
\emph{Named identifiers syntactically cited in the proof block:}
\begin{lstlisting}[language=EasyCrypt,basicstyle=\ttfamily\scriptsize]
programmed_site_prequery_one_step_challenge_fset_bound_for
programmed_site_query_min_entropy_bound
\end{lstlisting}

\end{formaltheorem}

\begin{formaltheorem}[EasyCrypt line 1095]
\noindent\textbf{EasyCrypt name.}
\begin{lstlisting}[language=EasyCrypt,basicstyle=\ttfamily\scriptsize]
programmed_site_prequery_one_step_challenge_budget_bound_for
\end{lstlisting}
\noindent\textbf{Checked EasyCrypt statement.}
\begin{lstlisting}[language=EasyCrypt,basicstyle=\ttfamily\scriptsize]
lemma programmed_site_prequery_one_step_challenge_budget_bound_for
  (d : random_coins distr) md sk m ctx qs :
  (card (oflist (filter ro_query_is_challenge qs)))%r <=
    rom_hash_query_budget + 1%r =>
  programmed_site_query_min_entropy_bound_for d md sk m ctx
    (filter ro_query_is_challenge qs)
    (1%r / challenge_support_cardinality_lower_bound) =>
  mu d
    (fun coins =>
      programming_site_prequeried qs
        (honest_signing_programming_site md sk m ctx coins)) <=
  (rom_hash_query_budget + 1%r) /
    challenge_support_cardinality_lower_bound.
\end{lstlisting}
\noindent\textbf{Formal mathematical reading.}
Mathematical theorem. This is an inequality, usually an arithmetic loss bound, event inclusion bound, or monotonicity fact used to compose probabilities.

\noindent\textbf{Role in the proof.}
Use in this chapter. It contributes directly to an advantage bound, bad-event bound, or real-arithmetic composition step.

\noindent\textbf{Proof-script interpretation.}
Proof-script reading. The machine proof decomposes the theorem into rewrites, program-logic obligations, relational equivalences, and arithmetic side conditions.
\emph{Main tactics visible in the proof script:} \ec{have}, \ec{apply}, \ec{rewrite}, \ec{smt}.
\emph{Named identifiers syntactically cited in the proof block:}
\begin{lstlisting}[language=EasyCrypt,basicstyle=\ttfamily\scriptsize]
card
card_budget
challenge_support_cardinality_lower_bound
challenge_support_cardinality_lower_bound_nonnegative
ctx
filter
inv
inv_nonneg
invr_ge0
ler_trans
ler_wpmul2r
md
oflist
programmed_site_prequery_one_step_challenge_fset_bound_for
qs
ro_query_is_challenge
rom_hash_query_budget
sk
spread
\end{lstlisting}

\end{formaltheorem}

\begin{formaltheorem}[EasyCrypt line 1129]
\noindent\textbf{EasyCrypt name.}
\begin{lstlisting}[language=EasyCrypt,basicstyle=\ttfamily\scriptsize]
programmed_site_prequery_one_step_challenge_budget_bound
\end{lstlisting}
\noindent\textbf{Checked EasyCrypt statement.}
\begin{lstlisting}[language=EasyCrypt,basicstyle=\ttfamily\scriptsize]
lemma programmed_site_prequery_one_step_challenge_budget_bound
  md sk m ctx qs :
  (card (oflist (filter ro_query_is_challenge qs)))%r <=
    rom_hash_query_budget + 1%r =>
  programmed_site_query_min_entropy_bound md sk m ctx
    (filter ro_query_is_challenge qs)
    (1%r / challenge_support_cardinality_lower_bound) =>
  mu drandom_coins
    (fun coins =>
      programming_site_prequeried qs
        (honest_signing_programming_site md sk m ctx coins)) <=
  (rom_hash_query_budget + 1%r) /
    challenge_support_cardinality_lower_bound.
\end{lstlisting}
\noindent\textbf{Formal mathematical reading.}
Mathematical theorem. This is an inequality, usually an arithmetic loss bound, event inclusion bound, or monotonicity fact used to compose probabilities.

\noindent\textbf{Role in the proof.}
Use in this chapter. It contributes directly to an advantage bound, bad-event bound, or real-arithmetic composition step.

\noindent\textbf{Proof-script interpretation.}
Proof-script reading. The machine proof decomposes the theorem into rewrites, program-logic obligations, relational equivalences, and arithmetic side conditions.
\emph{Main tactics visible in the proof script:} \ec{rewrite}, \ec{apply}.
\emph{Named identifiers syntactically cited in the proof block:}
\begin{lstlisting}[language=EasyCrypt,basicstyle=\ttfamily\scriptsize]
card_budget
programmed_site_prequery_one_step_challenge_budget_bound_for
programmed_site_query_min_entropy_bound
\end{lstlisting}

\end{formaltheorem}

\begin{formaltheorem}[EasyCrypt line 1148]
\noindent\textbf{EasyCrypt name.}
\begin{lstlisting}[language=EasyCrypt,basicstyle=\ttfamily\scriptsize]
programmed_site_prequery_one_step_challenge_concrete_bound
\end{lstlisting}
\noindent\textbf{Checked EasyCrypt statement.}
\begin{lstlisting}[language=EasyCrypt,basicstyle=\ttfamily\scriptsize]
lemma programmed_site_prequery_one_step_challenge_concrete_bound
  md sk m ctx qs :
  (card (oflist (filter ro_query_is_challenge qs)))%r <=
    rom_hash_query_budget + 1%r =>
  mu drandom_coins
    (fun coins =>
      programming_site_prequeried qs
        (honest_signing_programming_site md sk m ctx coins)) <=
  (rom_hash_query_budget + 1%r) /
    challenge_support_cardinality_lower_bound.
\end{lstlisting}
\noindent\textbf{Formal mathematical reading.}
Mathematical theorem. This is an inequality, usually an arithmetic loss bound, event inclusion bound, or monotonicity fact used to compose probabilities.

\noindent\textbf{Role in the proof.}
Use in this chapter. It contributes directly to an advantage bound, bad-event bound, or real-arithmetic composition step.

\noindent\textbf{Proof-script interpretation.}
Proof-script reading. The machine proof decomposes the theorem into rewrites, program-logic obligations, relational equivalences, and arithmetic side conditions.
\emph{Main tactics visible in the proof script:} \ec{apply}.
\emph{Named identifiers syntactically cited in the proof block:}
\begin{lstlisting}[language=EasyCrypt,basicstyle=\ttfamily\scriptsize]
card_budget
honest_signing_programming_site_query_spread_bound
programmed_site_prequery_one_step_challenge_budget_bound
\end{lstlisting}

\end{formaltheorem}

\begin{formaltheorem}[EasyCrypt line 1165]
\noindent\textbf{EasyCrypt name.}
\begin{lstlisting}[language=EasyCrypt,basicstyle=\ttfamily\scriptsize]
programmed_site_prequery_one_step_challenge_concrete_bound_for
\end{lstlisting}
\noindent\textbf{Checked EasyCrypt statement.}
\begin{lstlisting}[language=EasyCrypt,basicstyle=\ttfamily\scriptsize]
lemma programmed_site_prequery_one_step_challenge_concrete_bound_for
  (d : random_coins distr) md sk m ctx qs :
  signing_randomness_token_spread d =>
  (card (oflist (filter ro_query_is_challenge qs)))%r <=
    rom_hash_query_budget + 1%r =>
  mu d
    (fun coins =>
      programming_site_prequeried qs
        (honest_signing_programming_site md sk m ctx coins)) <=
  (rom_hash_query_budget + 1%r) /
    challenge_support_cardinality_lower_bound.
\end{lstlisting}
\noindent\textbf{Formal mathematical reading.}
Mathematical theorem. This is an inequality, usually an arithmetic loss bound, event inclusion bound, or monotonicity fact used to compose probabilities.

\noindent\textbf{Role in the proof.}
Use in this chapter. It contributes directly to an advantage bound, bad-event bound, or real-arithmetic composition step.

\noindent\textbf{Proof-script interpretation.}
Proof-script reading. The machine proof decomposes the theorem into rewrites, program-logic obligations, relational equivalences, and arithmetic side conditions.
\emph{Main tactics visible in the proof script:} \ec{apply}.
\emph{Named identifiers syntactically cited in the proof block:}
\begin{lstlisting}[language=EasyCrypt,basicstyle=\ttfamily\scriptsize]
card_budget
honest_signing_programming_site_query_spread_bound_for
programmed_site_prequery_one_step_challenge_budget_bound_for
spread
\end{lstlisting}

\end{formaltheorem}

\begin{formaltheorem}[EasyCrypt line 1183]
\noindent\textbf{EasyCrypt name.}
\begin{lstlisting}[language=EasyCrypt,basicstyle=\ttfamily\scriptsize]
challenge_query_log_count_budget_card_bound
\end{lstlisting}
\noindent\textbf{Checked EasyCrypt statement.}
\begin{lstlisting}[language=EasyCrypt,basicstyle=\ttfamily\scriptsize]
lemma challenge_query_log_count_budget_card_bound hc qs :
  0 <= hc =>
  hc <= hash_query_budget_count =>
  challenge_query_log_count qs <= hc + 1 =>
  (card (oflist (filter ro_query_is_challenge qs)))%r <=
    rom_hash_query_budget + 1%r.
\end{lstlisting}
\noindent\textbf{Formal mathematical reading.}
Mathematical theorem. This is an inequality, usually an arithmetic loss bound, event inclusion bound, or monotonicity fact used to compose probabilities.

\noindent\textbf{Role in the proof.}
Use in this chapter. It contributes directly to an advantage bound, bad-event bound, or real-arithmetic composition step.

\noindent\textbf{Proof-script interpretation.}
Proof-script reading. The machine proof decomposes the theorem into rewrites, program-logic obligations, relational equivalences, and arithmetic side conditions.
\emph{Main tactics visible in the proof script:} \ec{rewrite}, \ec{smt}.
\emph{Named identifiers syntactically cited in the proof block:}
\begin{lstlisting}[language=EasyCrypt,basicstyle=\ttfamily\scriptsize]
challenge_query_log_count
hash_query_budget_count
hash_query_count
rom_hash_query_budget
\end{lstlisting}

\end{formaltheorem}

\begin{formaltheorem}[EasyCrypt line 1195]
\noindent\textbf{EasyCrypt name.}
\begin{lstlisting}[language=EasyCrypt,basicstyle=\ttfamily\scriptsize]
programmed_site_prequery_one_step_challenge_counter_bound
\end{lstlisting}
\noindent\textbf{Checked EasyCrypt statement.}
\begin{lstlisting}[language=EasyCrypt,basicstyle=\ttfamily\scriptsize]
lemma programmed_site_prequery_one_step_challenge_counter_bound
  md sk m ctx qs hc :
  0 <= hc =>
  hc <= hash_query_budget_count =>
  challenge_query_log_count qs <= hc + 1 =>
  mu drandom_coins
    (fun coins =>
      programming_site_prequeried qs
        (honest_signing_programming_site md sk m ctx coins)) <=
  (rom_hash_query_budget + 1%r) /
    challenge_support_cardinality_lower_bound.
\end{lstlisting}
\noindent\textbf{Formal mathematical reading.}
Mathematical theorem. This is an inequality, usually an arithmetic loss bound, event inclusion bound, or monotonicity fact used to compose probabilities.

\noindent\textbf{Role in the proof.}
Use in this chapter. It contributes directly to an advantage bound, bad-event bound, or real-arithmetic composition step.

\noindent\textbf{Proof-script interpretation.}
Proof-script reading. The machine proof decomposes the theorem into rewrites, program-logic obligations, relational equivalences, and arithmetic side conditions.
\emph{Main tactics visible in the proof script:} \ec{apply}.
\emph{Named identifiers syntactically cited in the proof block:}
\begin{lstlisting}[language=EasyCrypt,basicstyle=\ttfamily\scriptsize]
challenge_query_log_count_budget_card_bound
hc
hc_budget
hc_ge0
programmed_site_prequery_one_step_challenge_concrete_bound
qs
query_count
\end{lstlisting}

\end{formaltheorem}

\begin{formaltheorem}[EasyCrypt line 1212]
\noindent\textbf{EasyCrypt name.}
\begin{lstlisting}[language=EasyCrypt,basicstyle=\ttfamily\scriptsize]
programmed_site_prequery_one_step_challenge_counter_bound_for
\end{lstlisting}
\noindent\textbf{Checked EasyCrypt statement.}
\begin{lstlisting}[language=EasyCrypt,basicstyle=\ttfamily\scriptsize]
lemma programmed_site_prequery_one_step_challenge_counter_bound_for
  (d : random_coins distr) md sk m ctx qs hc :
  signing_randomness_token_spread d =>
  0 <= hc =>
  hc <= hash_query_budget_count =>
  challenge_query_log_count qs <= hc + 1 =>
  mu d
    (fun coins =>
      programming_site_prequeried qs
        (honest_signing_programming_site md sk m ctx coins)) <=
  (rom_hash_query_budget + 1%r) /
    challenge_support_cardinality_lower_bound.
\end{lstlisting}
\noindent\textbf{Formal mathematical reading.}
Mathematical theorem. This is an inequality, usually an arithmetic loss bound, event inclusion bound, or monotonicity fact used to compose probabilities.

\noindent\textbf{Role in the proof.}
Use in this chapter. It contributes directly to an advantage bound, bad-event bound, or real-arithmetic composition step.

\noindent\textbf{Proof-script interpretation.}
Proof-script reading. The machine proof decomposes the theorem into rewrites, program-logic obligations, relational equivalences, and arithmetic side conditions.
\emph{Main tactics visible in the proof script:} \ec{apply}.
\emph{Named identifiers syntactically cited in the proof block:}
\begin{lstlisting}[language=EasyCrypt,basicstyle=\ttfamily\scriptsize]
challenge_query_log_count_budget_card_bound
hc
hc_budget
hc_ge0
programmed_site_prequery_one_step_challenge_concrete_bound_for
qs
query_count
spread
\end{lstlisting}

\end{formaltheorem}

\begin{formaltheorem}[EasyCrypt line 1231]
\noindent\textbf{EasyCrypt name.}
\begin{lstlisting}[language=EasyCrypt,basicstyle=\ttfamily\scriptsize]
programmed_site_prequery_one_step_message_hash_counter_bound
\end{lstlisting}
\noindent\textbf{Checked EasyCrypt statement.}
\begin{lstlisting}[language=EasyCrypt,basicstyle=\ttfamily\scriptsize]
lemma programmed_site_prequery_one_step_message_hash_counter_bound
  md pk sk m ctx qpk qctx qm qs hc :
  0 <= hc =>
  hc <= hash_query_budget_count =>
  challenge_query_log_count qs <= hc + 1 =>
  mu drandom_coins
    (fun coins =>
      programming_site_prequeried (message_hash_query qpk qctx qm :: qs)
        (actual_signature_programming_site md
          (transcript_of_signature md pk m ctx
            (sign_internal md sk m ctx coins)))) <=
  (rom_hash_query_budget + 1%r) /
    challenge_support_cardinality_lower_bound.
\end{lstlisting}
\noindent\textbf{Formal mathematical reading.}
Mathematical theorem. This is an inequality, usually an arithmetic loss bound, event inclusion bound, or monotonicity fact used to compose probabilities.

\noindent\textbf{Role in the proof.}
Use in this chapter. It contributes directly to an advantage bound, bad-event bound, or real-arithmetic composition step.

\noindent\textbf{Proof-script interpretation.}
Proof-script reading. The machine proof decomposes the theorem into rewrites, program-logic obligations, relational equivalences, and arithmetic side conditions.
\emph{Main tactics visible in the proof script:} \ec{rewrite}, \ec{apply}.
\emph{Named identifiers syntactically cited in the proof block:}
\begin{lstlisting}[language=EasyCrypt,basicstyle=\ttfamily\scriptsize]
actual_signature_programming_site_prequeried_message_hash_any_pkE
coins
ctx
hc
hc_budget
hc_ge0
honest_signing_programming_site
md
mu_eq
pk
programmed_site_prequery_one_step_challenge_counter_bound
programming_site_prequeried
qctx
qm
qpk
qs
query_count
sk
\end{lstlisting}

\end{formaltheorem}

\begin{formaltheorem}[EasyCrypt line 1257]
\noindent\textbf{EasyCrypt name.}
\begin{lstlisting}[language=EasyCrypt,basicstyle=\ttfamily\scriptsize]
direct_signing_coin_sampler_prequery_bound
\end{lstlisting}
\noindent\textbf{Checked EasyCrypt statement.}
\begin{lstlisting}[language=EasyCrypt,basicstyle=\ttfamily\scriptsize]
lemma direct_signing_coin_sampler_prequery_bound
  md sk m ctx qs hc :
  0 <= hc =>
  hc <= hash_query_budget_count =>
  challenge_query_log_count qs <= hc + 1 =>
  phoare[DirectSigningCoinSampler.sample :
    true ==>
    programming_site_prequeried qs
      (honest_signing_programming_site md sk m ctx res)] <=
  ((rom_hash_query_budget + 1%r) /
    challenge_support_cardinality_lower_bound).
\end{lstlisting}
\noindent\textbf{Formal mathematical reading.}
Mathematical theorem. This is a relational program theorem: two games or procedures are coupled so that a precondition on their initial states implies the stated postcondition on final states and return values.

\noindent\textbf{Role in the proof.}
Use in this chapter. It contributes directly to an advantage bound, bad-event bound, or real-arithmetic composition step.

\noindent\textbf{Proof-script interpretation.}
Proof-script reading. The machine proof decomposes the theorem into rewrites, program-logic obligations, relational equivalences, and arithmetic side conditions.
\emph{Main tactics visible in the proof script:} \ec{byphoare}, \ec{proc}, \ec{rnd}, \ec{apply}.
\emph{Named identifiers syntactically cited in the proof block:}
\begin{lstlisting}[language=EasyCrypt,basicstyle=\ttfamily\scriptsize]
bypr
coins
ctx
hc
hc_budget
hc_ge0
honest_signing_programming_site
md
programmed_site_prequery_one_step_challenge_counter_bound
programming_site_prequeried
qs
query_count
sk
\end{lstlisting}

\end{formaltheorem}

\begin{formaltheorem}[EasyCrypt line 1283]
\noindent\textbf{EasyCrypt name.}
\begin{lstlisting}[language=EasyCrypt,basicstyle=\ttfamily\scriptsize]
direct_signing_coin_sampler_budgeted_prequery_bound
\end{lstlisting}
\noindent\textbf{Checked EasyCrypt statement.}
\begin{lstlisting}[language=EasyCrypt,basicstyle=\ttfamily\scriptsize]
lemma direct_signing_coin_sampler_budgeted_prequery_bound
  md sk m ctx qs :
  challenge_query_log_count qs <= hash_query_budget_count + 1 =>
  phoare[DirectSigningCoinSampler.sample :
    true ==>
    programming_site_prequeried qs
      (honest_signing_programming_site md sk m ctx res)] <=
  ((rom_hash_query_budget + 1%r) /
    challenge_support_cardinality_lower_bound).
\end{lstlisting}
\noindent\textbf{Formal mathematical reading.}
Mathematical theorem. This is a relational program theorem: two games or procedures are coupled so that a precondition on their initial states implies the stated postcondition on final states and return values.

\noindent\textbf{Role in the proof.}
Use in this chapter. It contributes directly to an advantage bound, bad-event bound, or real-arithmetic composition step.

\noindent\textbf{Proof-script interpretation.}
Proof-script reading. The machine proof decomposes the theorem into rewrites, program-logic obligations, relational equivalences, and arithmetic side conditions.
\emph{Main tactics visible in the proof script:} \ec{apply}, \ec{rewrite}, \ec{smt}.
\emph{Named identifiers syntactically cited in the proof block:}
\begin{lstlisting}[language=EasyCrypt,basicstyle=\ttfamily\scriptsize]
ctx
direct_signing_coin_sampler_prequery_bound
hash_query_budget_count
md
qs
query_count
sk
\end{lstlisting}

\end{formaltheorem}

\begin{formaltheorem}[EasyCrypt line 1301]
\noindent\textbf{EasyCrypt name.}
\begin{lstlisting}[language=EasyCrypt,basicstyle=\ttfamily\scriptsize]
direct_signing_coin_sampler_message_hash_prequery_bound
\end{lstlisting}
\noindent\textbf{Checked EasyCrypt statement.}
\begin{lstlisting}[language=EasyCrypt,basicstyle=\ttfamily\scriptsize]
lemma direct_signing_coin_sampler_message_hash_prequery_bound
  md pk sk m ctx qpk qctx qm qs hc :
  0 <= hc =>
  hc <= hash_query_budget_count =>
  challenge_query_log_count qs <= hc + 1 =>
  phoare[DirectSigningCoinSampler.sample :
    true ==>
    programming_site_prequeried (message_hash_query qpk qctx qm :: qs)
      (actual_signature_programming_site md
        (transcript_of_signature md pk m ctx
          (sign_internal md sk m ctx res)))] <=
	  ((rom_hash_query_budget + 1%r) /
	    challenge_support_cardinality_lower_bound).
\end{lstlisting}
\noindent\textbf{Formal mathematical reading.}
Mathematical theorem. This is a relational program theorem: two games or procedures are coupled so that a precondition on their initial states implies the stated postcondition on final states and return values.

\noindent\textbf{Role in the proof.}
Use in this chapter. It contributes directly to an advantage bound, bad-event bound, or real-arithmetic composition step.

\noindent\textbf{Proof-script interpretation.}
Proof-script reading. The machine proof decomposes the theorem into rewrites, program-logic obligations, relational equivalences, and arithmetic side conditions.
\emph{Main tactics visible in the proof script:} \ec{byphoare}, \ec{proc}, \ec{rnd}, \ec{apply}.
\emph{Named identifiers syntactically cited in the proof block:}
\begin{lstlisting}[language=EasyCrypt,basicstyle=\ttfamily\scriptsize]
actual_signature_programming_site
bypr
coins
ctx
hc
hc_budget
hc_ge0
m0
md
message_hash_query
pk
programmed_site_prequery_one_step_message_hash_counter_bound
programming_site_prequeried
qctx
qm
qpk
qs
query_count
sign_internal
sk
transcript_of_signature
\end{lstlisting}

\end{formaltheorem}

\begin{formaltheorem}[EasyCrypt line 1333]
\noindent\textbf{EasyCrypt name.}
\begin{lstlisting}[language=EasyCrypt,basicstyle=\ttfamily\scriptsize]
direct_signing_coin_sampler_message_hash_budgeted_prequery_bound
\end{lstlisting}
\noindent\textbf{Checked EasyCrypt statement.}
\begin{lstlisting}[language=EasyCrypt,basicstyle=\ttfamily\scriptsize]
lemma direct_signing_coin_sampler_message_hash_budgeted_prequery_bound
  md pk sk m ctx qpk qctx qm qs :
  challenge_query_log_count qs <= hash_query_budget_count + 1 =>
  phoare[DirectSigningCoinSampler.sample :
    true ==>
    programming_site_prequeried (message_hash_query qpk qctx qm :: qs)
      (actual_signature_programming_site md
        (transcript_of_signature md pk m ctx
          (sign_internal md sk m ctx res)))] <=
	  ((rom_hash_query_budget + 1%r) /
	    challenge_support_cardinality_lower_bound).
\end{lstlisting}
\noindent\textbf{Formal mathematical reading.}
Mathematical theorem. This is a relational program theorem: two games or procedures are coupled so that a precondition on their initial states implies the stated postcondition on final states and return values.

\noindent\textbf{Role in the proof.}
Use in this chapter. It contributes directly to an advantage bound, bad-event bound, or real-arithmetic composition step.

\noindent\textbf{Proof-script interpretation.}
Proof-script reading. The machine proof decomposes the theorem into rewrites, program-logic obligations, relational equivalences, and arithmetic side conditions.
\emph{Main tactics visible in the proof script:} \ec{apply}, \ec{rewrite}, \ec{smt}.
\emph{Named identifiers syntactically cited in the proof block:}
\begin{lstlisting}[language=EasyCrypt,basicstyle=\ttfamily\scriptsize]
ctx
direct_signing_coin_sampler_message_hash_prequery_bound
hash_query_budget_count
md
pk
qctx
qm
qpk
qs
query_count
sk
\end{lstlisting}

\end{formaltheorem}

\begin{formaltheorem}[EasyCrypt line 1353]
\noindent\textbf{EasyCrypt name.}
\begin{lstlisting}[language=EasyCrypt,basicstyle=\ttfamily\scriptsize]
programmed_site_prequery_one_step_concrete_bound
\end{lstlisting}
\noindent\textbf{Checked EasyCrypt statement.}
\begin{lstlisting}[language=EasyCrypt,basicstyle=\ttfamily\scriptsize]
lemma programmed_site_prequery_one_step_concrete_bound
  md sk m ctx qs :
  (card (oflist qs))%r <= rom_hash_query_budget + 1%r =>
  mu drandom_coins
    (fun coins =>
      programming_site_prequeried qs
        (honest_signing_programming_site md sk m ctx coins)) <=
  (rom_hash_query_budget + 1%r) /
    challenge_support_cardinality_lower_bound.
\end{lstlisting}
\noindent\textbf{Formal mathematical reading.}
Mathematical theorem. This is an inequality, usually an arithmetic loss bound, event inclusion bound, or monotonicity fact used to compose probabilities.

\noindent\textbf{Role in the proof.}
Use in this chapter. It contributes directly to an advantage bound, bad-event bound, or real-arithmetic composition step.

\noindent\textbf{Proof-script interpretation.}
Proof-script reading. The machine proof decomposes the theorem into rewrites, program-logic obligations, relational equivalences, and arithmetic side conditions.
\emph{Main tactics visible in the proof script:} \ec{apply}.
\emph{Named identifiers syntactically cited in the proof block:}
\begin{lstlisting}[language=EasyCrypt,basicstyle=\ttfamily\scriptsize]
card_budget
honest_signing_programming_site_query_spread_bound
programmed_site_prequery_one_step_budget_bound
\end{lstlisting}

\end{formaltheorem}

\begin{formaltheorem}[EasyCrypt line 1369]
\noindent\textbf{EasyCrypt name.}
\begin{lstlisting}[language=EasyCrypt,basicstyle=\ttfamily\scriptsize]
programmed_site_prequery_reprogram_one_step_concrete_bound
\end{lstlisting}
\noindent\textbf{Checked EasyCrypt statement.}
\begin{lstlisting}[language=EasyCrypt,basicstyle=\ttfamily\scriptsize]
lemma programmed_site_prequery_reprogram_one_step_concrete_bound
  md sk m ctx qs sites :
  (forall coins,
    programming_site_conflict_free
      (honest_signing_programming_site md sk m ctx coins) sites) =>
  (card (oflist qs))%r <= rom_hash_query_budget + 1%r =>
  mu drandom_coins
    (fun coins =>
      programming_site_prequeried qs
        (honest_signing_programming_site md sk m ctx coins) \/
      programming_site_reprograms sites
        (honest_signing_programming_site md sk m ctx coins)) <=
  (rom_hash_query_budget + 1%r) /
    challenge_support_cardinality_lower_bound.
\end{lstlisting}
\noindent\textbf{Formal mathematical reading.}
Mathematical theorem. This is an inequality, usually an arithmetic loss bound, event inclusion bound, or monotonicity fact used to compose probabilities.

\noindent\textbf{Role in the proof.}
Use in this chapter. It contributes directly to an advantage bound, bad-event bound, or real-arithmetic composition step.

\noindent\textbf{Proof-script interpretation.}
Proof-script reading. The machine proof decomposes the theorem into rewrites, program-logic obligations, relational equivalences, and arithmetic side conditions.
\emph{Main tactics visible in the proof script:} \ec{apply}, \ec{case}, \ec{have}, \ec{rewrite}.
\emph{Named identifiers syntactically cited in the proof block:}
\begin{lstlisting}[language=EasyCrypt,basicstyle=\ttfamily\scriptsize]
card_budget
coins
conflict_free
ctx
drandom_coins
honest_signing_programming_site
ler_trans
md
mu
mu_le
no_reprog
preq
programmed_site_prequery_one_step_concrete_bound
programming_site_conflict_free_no_reprograms
programming_site_prequeried
programming_site_reprograms
qs
reprog
sites
sk
\end{lstlisting}

\end{formaltheorem}

\begin{formaltheorem}[EasyCrypt line 1404]
\noindent\textbf{EasyCrypt name.}
\begin{lstlisting}[language=EasyCrypt,basicstyle=\ttfamily\scriptsize]
programmed_site_query_min_entropy_bound_constant_impossible
\end{lstlisting}
\noindent\textbf{Checked EasyCrypt statement.}
\begin{lstlisting}[language=EasyCrypt,basicstyle=\ttfamily\scriptsize]
lemma programmed_site_query_min_entropy_bound_constant_impossible
  md sk m ctx qs bd q :
  q \in qs =>
  (forall coins,
    honest_signing_programming_site_query md sk m ctx coins = q) =>
  bd < 1%r =>
  programmed_site_query_min_entropy_bound md sk m ctx qs bd => false.
\end{lstlisting}
\noindent\textbf{Formal mathematical reading.}
Mathematical theorem. This is an implication, normally turning one invariant, validity predicate, or proof obligation into another.

\noindent\textbf{Role in the proof.}
Use in this chapter. It contributes directly to an advantage bound, bad-event bound, or real-arithmetic composition step.

\noindent\textbf{Proof-script interpretation.}
Proof-script reading. The machine proof decomposes the theorem into rewrites, program-logic obligations, relational equivalences, and arithmetic side conditions.
\emph{Main tactics visible in the proof script:} \ec{have}, \ec{rewrite}, \ec{apply}, \ec{smt}.
\emph{Named identifiers syntactically cited in the proof block:}
\begin{lstlisting}[language=EasyCrypt,basicstyle=\ttfamily\scriptsize]
bd_lt1
coins
ctx
drandom_coins
eq1_mu
eq_sym
hit1
hit_le_bd
honest_signing_programming_site_query
md
mu
q_mem
query_constant
signing_coin_distribution_lossless
sk
spread
\end{lstlisting}

\end{formaltheorem}

\begin{formaltheorem}[EasyCrypt line 1426]
\noindent\textbf{EasyCrypt name.}
\begin{lstlisting}[language=EasyCrypt,basicstyle=\ttfamily\scriptsize]
programming_site_reprograms_conflict_cons
\end{lstlisting}
\noindent\textbf{Checked EasyCrypt statement.}
\begin{lstlisting}[language=EasyCrypt,basicstyle=\ttfamily\scriptsize]
lemma programming_site_reprograms_conflict_cons site old programmed :
  programming_sites_conflict site old =>
  programming_site_reprograms (old :: programmed) site.
\end{lstlisting}
\noindent\textbf{Formal mathematical reading.}
Mathematical theorem. This is an implication, normally turning one invariant, validity predicate, or proof obligation into another.

\noindent\textbf{Role in the proof.}
Use in this chapter. It is a reusable checked theorem cited by later proof scripts.

\noindent\textbf{Proof-script interpretation.}
Proof-script reading. The machine proof decomposes the theorem into rewrites, program-logic obligations, relational equivalences, and arithmetic side conditions.
\emph{Main tactics visible in the proof script:} \ec{rewrite}.
\emph{Named identifiers syntactically cited in the proof block:}
\begin{lstlisting}[language=EasyCrypt,basicstyle=\ttfamily\scriptsize]
conflict
programming_site_reprograms
\end{lstlisting}

\end{formaltheorem}

\begin{formaltheorem}[EasyCrypt line 1434]
\noindent\textbf{EasyCrypt name.}
\begin{lstlisting}[language=EasyCrypt,basicstyle=\ttfamily\scriptsize]
rom_counted_programming_bad_min_entropy
\end{lstlisting}
\noindent\textbf{Checked EasyCrypt statement.}
\begin{lstlisting}[language=EasyCrypt,basicstyle=\ttfamily\scriptsize]
lemma rom_counted_programming_bad_min_entropy queries programmed :
  rom_counted_programming_bad queries programmed true.
\end{lstlisting}
\noindent\textbf{Formal mathematical reading.}
Mathematical theorem. This lemma is a universally quantified proposition over the variables shown in the EasyCrypt statement.

\noindent\textbf{Role in the proof.}
Use in this chapter. It is a reusable checked theorem cited by later proof scripts.

\noindent\textbf{Proof-script interpretation.}
No proof block was collected for this item.

\end{formaltheorem}

\begin{formaltheorem}[EasyCrypt line 1438]
\noindent\textbf{EasyCrypt name.}
\begin{lstlisting}[language=EasyCrypt,basicstyle=\ttfamily\scriptsize]
rom_counted_programming_bad_prequery_cons
\end{lstlisting}
\noindent\textbf{Checked EasyCrypt statement.}
\begin{lstlisting}[language=EasyCrypt,basicstyle=\ttfamily\scriptsize]
lemma rom_counted_programming_bad_prequery_cons queries programmed site
  entropy_bad :
  programming_site_prequeried queries site =>
  rom_counted_programming_bad queries (site :: programmed) entropy_bad.
\end{lstlisting}
\noindent\textbf{Formal mathematical reading.}
Mathematical theorem. This is an implication, normally turning one invariant, validity predicate, or proof obligation into another.

\noindent\textbf{Role in the proof.}
Use in this chapter. It is a reusable checked theorem cited by later proof scripts.

\noindent\textbf{Proof-script interpretation.}
Proof-script reading. The machine proof decomposes the theorem into rewrites, program-logic obligations, relational equivalences, and arithmetic side conditions.
\emph{Main tactics visible in the proof script:} \ec{rewrite}, \ec{case}.
\emph{Named identifiers syntactically cited in the proof block:}
\begin{lstlisting}[language=EasyCrypt,basicstyle=\ttfamily\scriptsize]
entropy_bad
preq
rom_counted_programming_bad
\end{lstlisting}

\end{formaltheorem}

\begin{formaltheorem}[EasyCrypt line 1448]
\noindent\textbf{EasyCrypt name.}
\begin{lstlisting}[language=EasyCrypt,basicstyle=\ttfamily\scriptsize]
counted_rom_programming_bound_obligation_holds
\end{lstlisting}
\noindent\textbf{Checked EasyCrypt statement.}
\begin{lstlisting}[language=EasyCrypt,basicstyle=\ttfamily\scriptsize]
lemma counted_rom_programming_bound_obligation_holds :
  counted_rom_programming_bound_obligation.
\end{lstlisting}
\noindent\textbf{Formal mathematical reading.}
Mathematical theorem. This lemma is a universally quantified proposition over the variables shown in the EasyCrypt statement.

\noindent\textbf{Role in the proof.}
Use in this chapter. It contributes directly to an advantage bound, bad-event bound, or real-arithmetic composition step.

\noindent\textbf{Proof-script interpretation.}
Proof-script reading. The machine proof decomposes the theorem into rewrites, program-logic obligations, relational equivalences, and arithmetic side conditions.
\emph{Main tactics visible in the proof script:} \ec{rewrite}.
\emph{Named identifiers syntactically cited in the proof block:}
\begin{lstlisting}[language=EasyCrypt,basicstyle=\ttfamily\scriptsize]
counted_rom_programming_bound_obligation
counted_rom_programming_loss_term_nonnegative
\end{lstlisting}

\end{formaltheorem}

\begin{formaltheorem}[EasyCrypt line 1463]
\noindent\textbf{EasyCrypt name.}
\begin{lstlisting}[language=EasyCrypt,basicstyle=\ttfamily\scriptsize]
counted_lazy_rom_bad_flags_universal_bound
\end{lstlisting}
\noindent\textbf{Checked EasyCrypt statement.}
\begin{lstlisting}[language=EasyCrypt,basicstyle=\ttfamily\scriptsize]
lemma counted_lazy_rom_bad_flags_universal_bound &m :
  1%r >= Pr[CountedLazyROMBadGame(O, A).main() @ &m : res].
\end{lstlisting}
\noindent\textbf{Formal mathematical reading.}
Mathematical theorem. This theorem has the form \(\Pr[G:E] = B\) is a game-probability statement.  It is one of the exact hops, bad-event bounds, or final advantage inequalities used in the reduction.

\noindent\textbf{Role in the proof.}
Use in this chapter. It contributes directly to an advantage bound, bad-event bound, or real-arithmetic composition step.

\noindent\textbf{Proof-script interpretation.}
Proof-script reading. The machine proof decomposes the theorem into rewrites, program-logic obligations, relational equivalences, and arithmetic side conditions.
\emph{Main tactics visible in the proof script:} \ec{byphoare}.

\end{formaltheorem}

\begin{formaltheorem}[EasyCrypt line 1469]
\noindent\textbf{EasyCrypt name.}
\begin{lstlisting}[language=EasyCrypt,basicstyle=\ttfamily\scriptsize]
counted_lazy_rom_bad_flags_query_budget_bound
\end{lstlisting}
\noindent\textbf{Checked EasyCrypt statement.}
\begin{lstlisting}[language=EasyCrypt,basicstyle=\ttfamily\scriptsize]
lemma counted_lazy_rom_bad_flags_query_budget_bound &m :
  counted_lazy_rom_bad_flags_budget_sound
    (Pr[CountedLazyROMBadGame(O, A).main() @ &m : res]) =>
  Pr[CountedLazyROMBadGame(O, A).main() @ &m : res] <=
    counted_rom_programming_loss_term.
\end{lstlisting}
\noindent\textbf{Formal mathematical reading.}
Mathematical theorem. This theorem has the form \(\Pr[G:E] \le B\) is a game-probability statement.  It is one of the exact hops, bad-event bounds, or final advantage inequalities used in the reduction.

\noindent\textbf{Role in the proof.}
Use in this chapter. It contributes directly to an advantage bound, bad-event bound, or real-arithmetic composition step.

\noindent\textbf{Proof-script interpretation.}
Proof-script reading. The machine proof decomposes the theorem into rewrites, program-logic obligations, relational equivalences, and arithmetic side conditions.
\emph{Main tactics visible in the proof script:} \ec{rewrite}.
\emph{Named identifiers syntactically cited in the proof block:}
\begin{lstlisting}[language=EasyCrypt,basicstyle=\ttfamily\scriptsize]
counted_lazy_rom_bad_flags_budget_sound
\end{lstlisting}

\end{formaltheorem}

\begin{formaltheorem}[EasyCrypt line 1478]
\noindent\textbf{EasyCrypt name.}
\begin{lstlisting}[language=EasyCrypt,basicstyle=\ttfamily\scriptsize]
counted_lazy_rom_bad_flags_query_budget_bound_subunit
\end{lstlisting}
\noindent\textbf{Checked EasyCrypt statement.}
\begin{lstlisting}[language=EasyCrypt,basicstyle=\ttfamily\scriptsize]
lemma counted_lazy_rom_bad_flags_query_budget_bound_subunit &m :
  counted_lazy_rom_bad_flags_budget_sound
    (Pr[CountedLazyROMBadGame(O, A).main() @ &m : res]) =>
  Pr[CountedLazyROMBadGame(O, A).main() @ &m : res] < 1%r.
\end{lstlisting}
\noindent\textbf{Formal mathematical reading.}
Mathematical theorem. This theorem has the form \(\Pr[G:E] = B\) is a game-probability statement.  It is one of the exact hops, bad-event bounds, or final advantage inequalities used in the reduction.

\noindent\textbf{Role in the proof.}
Use in this chapter. It contributes directly to an advantage bound, bad-event bound, or real-arithmetic composition step.

\noindent\textbf{Proof-script interpretation.}
Proof-script reading. The machine proof decomposes the theorem into rewrites, program-logic obligations, relational equivalences, and arithmetic side conditions.
\emph{Main tactics visible in the proof script:} \ec{apply}, \ec{rewrite}.
\emph{Named identifiers syntactically cited in the proof block:}
\begin{lstlisting}[language=EasyCrypt,basicstyle=\ttfamily\scriptsize]
budget_sound
counted_lazy_rom_bad_flags_budget_sound
counted_rom_programming_loss_term
counted_rom_programming_loss_term_subunit
ler_lt_trans
\end{lstlisting}

\end{formaltheorem}

\begin{formaltheorem}[EasyCrypt line 1491]
\noindent\textbf{EasyCrypt name.}
\begin{lstlisting}[language=EasyCrypt,basicstyle=\ttfamily\scriptsize]
fs_with_aborts_no_programming_failure
\end{lstlisting}
\noindent\textbf{Checked EasyCrypt statement.}
\begin{lstlisting}[language=EasyCrypt,basicstyle=\ttfamily\scriptsize]
lemma fs_with_aborts_no_programming_failure md tr y :
  transcript_challenge_matches tr y =>
  fs_with_aborts_failure md tr (programming_site_of_transcript md tr y) =
  min_entropy_failure md tr.
\end{lstlisting}
\noindent\textbf{Formal mathematical reading.}
Mathematical theorem. This is an implication, normally turning one invariant, validity predicate, or proof obligation into another.

\noindent\textbf{Role in the proof.}
Use in this chapter. It is a reusable checked theorem cited by later proof scripts.

\noindent\textbf{Proof-script interpretation.}
Proof-script reading. The machine proof decomposes the theorem into rewrites, program-logic obligations, relational equivalences, and arithmetic side conditions.
\emph{Main tactics visible in the proof script:} \ec{rewrite}.
\emph{Named identifiers syntactically cited in the proof block:}
\begin{lstlisting}[language=EasyCrypt,basicstyle=\ttfamily\scriptsize]
fs_with_aborts_failure
match_y
md
programming_site_self_no_reprogramming_failure
tr
\end{lstlisting}

\end{formaltheorem}

\begin{formaltheorem}[EasyCrypt line 1501]
\noindent\textbf{EasyCrypt name.}
\begin{lstlisting}[language=EasyCrypt,basicstyle=\ttfamily\scriptsize]
transcript_signature_challenge_output_matches
\end{lstlisting}
\noindent\textbf{Checked EasyCrypt statement.}
\begin{lstlisting}[language=EasyCrypt,basicstyle=\ttfamily\scriptsize]
lemma transcript_signature_challenge_output_matches tr :
  transcript_signature_challenge_matches tr
    (transcript_signature_challenge_output tr).
\end{lstlisting}
\noindent\textbf{Formal mathematical reading.}
Mathematical theorem. This lemma is a universally quantified proposition over the variables shown in the EasyCrypt statement.

\noindent\textbf{Role in the proof.}
Use in this chapter. It contributes directly to an advantage bound, bad-event bound, or real-arithmetic composition step.

\noindent\textbf{Proof-script interpretation.}
Proof-script reading. The machine proof decomposes the theorem into rewrites, program-logic obligations, relational equivalences, and arithmetic side conditions.
\emph{Main tactics visible in the proof script:} \ec{rewrite}.
\emph{Named identifiers syntactically cited in the proof block:}
\begin{lstlisting}[language=EasyCrypt,basicstyle=\ttfamily\scriptsize]
ro_challenge_hash
ro_output_of_challenge
ro_output_to_challenge
transcript_signature_challenge_matches
transcript_signature_challenge_output
\end{lstlisting}

\end{formaltheorem}

\begin{formaltheorem}[EasyCrypt line 1511]
\noindent\textbf{EasyCrypt name.}
\begin{lstlisting}[language=EasyCrypt,basicstyle=\ttfamily\scriptsize]
honest_transcript_challenge_entropy_support
\end{lstlisting}
\noindent\textbf{Checked EasyCrypt statement.}
\begin{lstlisting}[language=EasyCrypt,basicstyle=\ttfamily\scriptsize]
lemma honest_transcript_challenge_entropy_support md sk m ctx coins :
  challenge_entropy_support_ok md
    (transcript_from_honest_signing md sk m ctx coins).
\end{lstlisting}
\noindent\textbf{Formal mathematical reading.}
Mathematical theorem. This lemma is a universally quantified proposition over the variables shown in the EasyCrypt statement.

\noindent\textbf{Role in the proof.}
Use in this chapter. It contributes directly to an advantage bound, bad-event bound, or real-arithmetic composition step.

\noindent\textbf{Proof-script interpretation.}
Proof-script reading. The machine proof decomposes the theorem into rewrites, program-logic obligations, relational equivalences, and arithmetic side conditions.
\emph{Main tactics visible in the proof script:} \ec{rewrite}, \ec{apply}.
\emph{Named identifiers syntactically cited in the proof block:}
\begin{lstlisting}[language=EasyCrypt,basicstyle=\ttfamily\scriptsize]
challenge_entropy_support_ok
challenge_from_seed_sparse
challenge_hash
sig_challenge
sign_internal
signature_challenge
transcript_challenge
transcript_from_honest_signing
transcript_of_signature
\end{lstlisting}

\end{formaltheorem}

\begin{formaltheorem}[EasyCrypt line 1521]
\noindent\textbf{EasyCrypt name.}
\begin{lstlisting}[language=EasyCrypt,basicstyle=\ttfamily\scriptsize]
honest_transcript_no_min_entropy_failure
\end{lstlisting}
\noindent\textbf{Checked EasyCrypt statement.}
\begin{lstlisting}[language=EasyCrypt,basicstyle=\ttfamily\scriptsize]
lemma honest_transcript_no_min_entropy_failure md sk m ctx coins :
  ! min_entropy_failure md
      (transcript_from_honest_signing md sk m ctx coins).
\end{lstlisting}
\noindent\textbf{Formal mathematical reading.}
Mathematical theorem. This lemma is a universally quantified proposition over the variables shown in the EasyCrypt statement.

\noindent\textbf{Role in the proof.}
Use in this chapter. It is a reusable checked theorem cited by later proof scripts.

\noindent\textbf{Proof-script interpretation.}
Proof-script reading. The machine proof decomposes the theorem into rewrites, program-logic obligations, relational equivalences, and arithmetic side conditions.
\emph{Main tactics visible in the proof script:} \ec{rewrite}.
\emph{Named identifiers syntactically cited in the proof block:}
\begin{lstlisting}[language=EasyCrypt,basicstyle=\ttfamily\scriptsize]
honest_transcript_challenge_entropy_support
min_entropy_failure
\end{lstlisting}

\end{formaltheorem}

\begin{formaltheorem}[EasyCrypt line 1529]
\noindent\textbf{EasyCrypt name.}
\begin{lstlisting}[language=EasyCrypt,basicstyle=\ttfamily\scriptsize]
signing_transcript_relation_challenge_entropy_support
\end{lstlisting}
\noindent\textbf{Checked EasyCrypt statement.}
\begin{lstlisting}[language=EasyCrypt,basicstyle=\ttfamily\scriptsize]
lemma signing_transcript_relation_challenge_entropy_support
  md pk m ctx sig tr :
  signing_transcript_relation md pk m ctx sig tr =>
  challenge_entropy_support_ok md tr.
\end{lstlisting}
\noindent\textbf{Formal mathematical reading.}
Mathematical theorem. This is an implication, normally turning one invariant, validity predicate, or proof obligation into another.

\noindent\textbf{Role in the proof.}
Use in this chapter. It contributes directly to an advantage bound, bad-event bound, or real-arithmetic composition step.

\noindent\textbf{Proof-script interpretation.}
Proof-script reading. The machine proof decomposes the theorem into rewrites, program-logic obligations, relational equivalences, and arithmetic side conditions.
\emph{Main tactics visible in the proof script:} \ec{rewrite}, \ec{case}, \ec{elim}, \ec{apply}.
\emph{Named identifiers syntactically cited in the proof block:}
\begin{lstlisting}[language=EasyCrypt,basicstyle=\ttfamily\scriptsize]
attempt
challenge_entropy_support_ok
challenge_from_seed_sparse
challenge_hash
coins
hsim
rel
rest
sig_challenge
sign_internal
signature_challenge
signing_attempt_relation
signing_simulation_relation
signing_transcript_relation
sk
trE
transcript_challenge
transcript_of_signature
witness
\end{lstlisting}

\end{formaltheorem}

\begin{formaltheorem}[EasyCrypt line 1549]
\noindent\textbf{EasyCrypt name.}
\begin{lstlisting}[language=EasyCrypt,basicstyle=\ttfamily\scriptsize]
signing_transcript_relation_no_min_entropy_failure
\end{lstlisting}
\noindent\textbf{Checked EasyCrypt statement.}
\begin{lstlisting}[language=EasyCrypt,basicstyle=\ttfamily\scriptsize]
lemma signing_transcript_relation_no_min_entropy_failure
  md pk m ctx sig tr :
  signing_transcript_relation md pk m ctx sig tr =>
  ! min_entropy_failure md tr.
\end{lstlisting}
\noindent\textbf{Formal mathematical reading.}
Mathematical theorem. This is an implication, normally turning one invariant, validity predicate, or proof obligation into another.

\noindent\textbf{Role in the proof.}
Use in this chapter. It is a reusable checked theorem cited by later proof scripts.

\noindent\textbf{Proof-script interpretation.}
Proof-script reading. The machine proof decomposes the theorem into rewrites, program-logic obligations, relational equivalences, and arithmetic side conditions.
\emph{Main tactics visible in the proof script:} \ec{rewrite}.
\emph{Named identifiers syntactically cited in the proof block:}
\begin{lstlisting}[language=EasyCrypt,basicstyle=\ttfamily\scriptsize]
ctx
md
min_entropy_failure
pk
rel
sig
signing_transcript_relation_challenge_entropy_support
tr
\end{lstlisting}

\end{formaltheorem}

\begin{formaltheorem}[EasyCrypt line 1560]
\noindent\textbf{EasyCrypt name.}
\begin{lstlisting}[language=EasyCrypt,basicstyle=\ttfamily\scriptsize]
signing_record_relation_no_min_entropy_failure
\end{lstlisting}
\noindent\textbf{Checked EasyCrypt statement.}
\begin{lstlisting}[language=EasyCrypt,basicstyle=\ttfamily\scriptsize]
lemma signing_record_relation_no_min_entropy_failure md pk r :
  signing_record_relation md pk r =>
  ! min_entropy_failure md (signing_record_transcript r).
\end{lstlisting}
\noindent\textbf{Formal mathematical reading.}
Mathematical theorem. This is an implication, normally turning one invariant, validity predicate, or proof obligation into another.

\noindent\textbf{Role in the proof.}
Use in this chapter. It is a reusable checked theorem cited by later proof scripts.

\noindent\textbf{Proof-script interpretation.}
Proof-script reading. The machine proof decomposes the theorem into rewrites, program-logic obligations, relational equivalences, and arithmetic side conditions.
\emph{Main tactics visible in the proof script:} \ec{rewrite}, \ec{apply}.
\emph{Named identifiers syntactically cited in the proof block:}
\begin{lstlisting}[language=EasyCrypt,basicstyle=\ttfamily\scriptsize]
md
pk
rel
signing_record_context
signing_record_message
signing_record_relation
signing_record_signature
signing_transcript_relation_no_min_entropy_failure
\end{lstlisting}

\end{formaltheorem}

\begin{formaltheorem}[EasyCrypt line 1572]
\noindent\textbf{EasyCrypt name.}
\begin{lstlisting}[language=EasyCrypt,basicstyle=\ttfamily\scriptsize]
signing_record_log_sound_no_min_entropy
\end{lstlisting}
\noindent\textbf{Checked EasyCrypt statement.}
\begin{lstlisting}[language=EasyCrypt,basicstyle=\ttfamily\scriptsize]
lemma signing_record_log_sound_no_min_entropy md pk rs :
  signing_record_log_sound md pk rs =>
  signing_record_log_min_entropy_clear md rs.
\end{lstlisting}
\noindent\textbf{Formal mathematical reading.}
Mathematical theorem. This is an implication, normally turning one invariant, validity predicate, or proof obligation into another.

\noindent\textbf{Role in the proof.}
Use in this chapter. It proves a structural correctness fact needed by later extraction or verification arguments.

\noindent\textbf{Proof-script interpretation.}
Proof-script reading. The machine proof decomposes the theorem into rewrites, program-logic obligations, relational equivalences, and arithmetic side conditions.
\emph{Main tactics visible in the proof script:} \ec{rewrite}, \ec{elim}, \ec{split}, \ec{apply}.
\emph{Named identifiers syntactically cited in the proof block:}
\begin{lstlisting}[language=EasyCrypt,basicstyle=\ttfamily\scriptsize]
ih
md
pk
rel
rs
signing_record_log_min_entropy_clear
signing_record_log_sound
signing_record_relation_no_min_entropy_failure
sound
\end{lstlisting}

\end{formaltheorem}

\begin{formaltheorem}[EasyCrypt line 1584]
\noindent\textbf{EasyCrypt name.}
\begin{lstlisting}[language=EasyCrypt,basicstyle=\ttfamily\scriptsize]
signing_record_log_sound_transcripts_no_min_entropy
\end{lstlisting}
\noindent\textbf{Checked EasyCrypt statement.}
\begin{lstlisting}[language=EasyCrypt,basicstyle=\ttfamily\scriptsize]
lemma signing_record_log_sound_transcripts_no_min_entropy md pk rs :
  signing_record_log_sound md pk rs =>
  transcript_log_min_entropy_clear md (signing_record_transcripts rs).
\end{lstlisting}
\noindent\textbf{Formal mathematical reading.}
Mathematical theorem. This is an implication, normally turning one invariant, validity predicate, or proof obligation into another.

\noindent\textbf{Role in the proof.}
Use in this chapter. It proves a structural correctness fact needed by later extraction or verification arguments.

\noindent\textbf{Proof-script interpretation.}
Proof-script reading. The machine proof decomposes the theorem into rewrites, program-logic obligations, relational equivalences, and arithmetic side conditions.
\emph{Main tactics visible in the proof script:} \ec{rewrite}, \ec{elim}, \ec{split}, \ec{apply}.
\emph{Named identifiers syntactically cited in the proof block:}
\begin{lstlisting}[language=EasyCrypt,basicstyle=\ttfamily\scriptsize]
ih
md
pk
rel
rs
signing_record_log_sound
signing_record_relation_no_min_entropy_failure
signing_record_transcripts
sound
transcript_log_min_entropy_clear
\end{lstlisting}

\end{formaltheorem}

\begin{formaltheorem}[EasyCrypt line 1597]
\noindent\textbf{EasyCrypt name.}
\begin{lstlisting}[language=EasyCrypt,basicstyle=\ttfamily\scriptsize]
fs_with_aborts_no_failure_for_signing_relation
\end{lstlisting}
\noindent\textbf{Checked EasyCrypt statement.}
\begin{lstlisting}[language=EasyCrypt,basicstyle=\ttfamily\scriptsize]
lemma fs_with_aborts_no_failure_for_signing_relation
  md pk m ctx sig tr site :
  signing_transcript_relation md pk m ctx sig tr =>
  programming_site_matches_transcript md tr site =>
  ! fs_with_aborts_failure md tr site.
\end{lstlisting}
\noindent\textbf{Formal mathematical reading.}
Mathematical theorem. This is an implication, normally turning one invariant, validity predicate, or proof obligation into another.

\noindent\textbf{Role in the proof.}
Use in this chapter. It is a reusable checked theorem cited by later proof scripts.

\noindent\textbf{Proof-script interpretation.}
Proof-script reading. The machine proof decomposes the theorem into rewrites, program-logic obligations, relational equivalences, and arithmetic side conditions.
\emph{Main tactics visible in the proof script:} \ec{rewrite}, \ec{apply}.
\emph{Named identifiers syntactically cited in the proof block:}
\begin{lstlisting}[language=EasyCrypt,basicstyle=\ttfamily\scriptsize]
ctx
fs_with_aborts_failure
match_site
md
pk
rel
reprogramming_failure
sig
signing_transcript_relation_no_min_entropy_failure
tr
\end{lstlisting}

\end{formaltheorem}

\begin{formaltheorem}[EasyCrypt line 1609]
\noindent\textbf{EasyCrypt name.}
\begin{lstlisting}[language=EasyCrypt,basicstyle=\ttfamily\scriptsize]
fs_with_aborts_no_failure_for_signing_record
\end{lstlisting}
\noindent\textbf{Checked EasyCrypt statement.}
\begin{lstlisting}[language=EasyCrypt,basicstyle=\ttfamily\scriptsize]
lemma fs_with_aborts_no_failure_for_signing_record md pk r site :
  signing_record_relation md pk r =>
  programming_site_matches_transcript
    md (signing_record_transcript r) site =>
  ! fs_with_aborts_failure md (signing_record_transcript r) site.
\end{lstlisting}
\noindent\textbf{Formal mathematical reading.}
Mathematical theorem. This is an implication, normally turning one invariant, validity predicate, or proof obligation into another.

\noindent\textbf{Role in the proof.}
Use in this chapter. It is a reusable checked theorem cited by later proof scripts.

\noindent\textbf{Proof-script interpretation.}
Proof-script reading. The machine proof decomposes the theorem into rewrites, program-logic obligations, relational equivalences, and arithmetic side conditions.
\emph{Main tactics visible in the proof script:} \ec{rewrite}, \ec{apply}.
\emph{Named identifiers syntactically cited in the proof block:}
\begin{lstlisting}[language=EasyCrypt,basicstyle=\ttfamily\scriptsize]
fs_with_aborts_no_failure_for_signing_relation
match_site
md
pk
rel
signing_record_context
signing_record_message
signing_record_relation
signing_record_signature
signing_record_transcript
site
\end{lstlisting}

\end{formaltheorem}

\begin{formaltheorem}[EasyCrypt line 1624]
\noindent\textbf{EasyCrypt name.}
\begin{lstlisting}[language=EasyCrypt,basicstyle=\ttfamily\scriptsize]
fs_with_aborts_no_failure_for_honest_programming
\end{lstlisting}
\noindent\textbf{Checked EasyCrypt statement.}
\begin{lstlisting}[language=EasyCrypt,basicstyle=\ttfamily\scriptsize]
lemma fs_with_aborts_no_failure_for_honest_programming md sk m ctx coins y :
  transcript_challenge_matches
    (transcript_from_honest_signing md sk m ctx coins) y =>
  ! fs_with_aborts_failure md
      (transcript_from_honest_signing md sk m ctx coins)
      (programming_site_of_transcript md
        (transcript_from_honest_signing md sk m ctx coins) y).
\end{lstlisting}
\noindent\textbf{Formal mathematical reading.}
Mathematical theorem. This is an implication, normally turning one invariant, validity predicate, or proof obligation into another.

\noindent\textbf{Role in the proof.}
Use in this chapter. It is a reusable checked theorem cited by later proof scripts.

\noindent\textbf{Proof-script interpretation.}
Proof-script reading. The machine proof decomposes the theorem into rewrites, program-logic obligations, relational equivalences, and arithmetic side conditions.
\emph{Main tactics visible in the proof script:} \ec{rewrite}.
\emph{Named identifiers syntactically cited in the proof block:}
\begin{lstlisting}[language=EasyCrypt,basicstyle=\ttfamily\scriptsize]
coins
ctx
fs_with_aborts_failure
honest_transcript_no_min_entropy_failure
match_y
md
programming_site_self_no_reprogramming_failure
sk
transcript_from_honest_signing
\end{lstlisting}

\end{formaltheorem}

\begin{formaltheorem}[EasyCrypt line 1639]
\noindent\textbf{EasyCrypt name.}
\begin{lstlisting}[language=EasyCrypt,basicstyle=\ttfamily\scriptsize]
transcript_valid_programming_site_signature_query
\end{lstlisting}
\noindent\textbf{Checked EasyCrypt statement.}
\begin{lstlisting}[language=EasyCrypt,basicstyle=\ttfamily\scriptsize]
lemma transcript_valid_programming_site_signature_query md tr y :
  transcript_valid md tr =>
  programming_site_query (programming_site_of_transcript md tr y) =
  programming_site_query (signature_programming_site_of_transcript md tr y).
\end{lstlisting}
\noindent\textbf{Formal mathematical reading.}
Mathematical theorem. This is an implication, normally turning one invariant, validity predicate, or proof obligation into another.

\noindent\textbf{Role in the proof.}
Use in this chapter. It proves a structural correctness fact needed by later extraction or verification arguments.

\noindent\textbf{Proof-script interpretation.}
Proof-script reading. The machine proof decomposes the theorem into rewrites, program-logic obligations, relational equivalences, and arithmetic side conditions.
\emph{Main tactics visible in the proof script:} \ec{rewrite}, \ec{apply}.
\emph{Named identifiers syntactically cited in the proof block:}
\begin{lstlisting}[language=EasyCrypt,basicstyle=\ttfamily\scriptsize]
programming_site_of_transcript
programming_site_query
signature_programming_site_of_transcript
transcript_valid_signature_challenge_query
valid
\end{lstlisting}

\end{formaltheorem}

\begin{formaltheorem}[EasyCrypt line 1650]
\noindent\textbf{EasyCrypt name.}
\begin{lstlisting}[language=EasyCrypt,basicstyle=\ttfamily\scriptsize]
transcript_valid_signature_programming_site_matches
\end{lstlisting}
\noindent\textbf{Checked EasyCrypt statement.}
\begin{lstlisting}[language=EasyCrypt,basicstyle=\ttfamily\scriptsize]
lemma transcript_valid_signature_programming_site_matches md tr y :
  transcript_valid md tr =>
  transcript_signature_challenge_matches tr y =>
  programming_site_matches_transcript md tr
    (signature_programming_site_of_transcript md tr y).
\end{lstlisting}
\noindent\textbf{Formal mathematical reading.}
Mathematical theorem. This is an implication, normally turning one invariant, validity predicate, or proof obligation into another.

\noindent\textbf{Role in the proof.}
Use in this chapter. It proves a structural correctness fact needed by later extraction or verification arguments.

\noindent\textbf{Proof-script interpretation.}
Proof-script reading. The machine proof decomposes the theorem into rewrites, program-logic obligations, relational equivalences, and arithmetic side conditions.
\emph{Main tactics visible in the proof script:} \ec{rewrite}.
\emph{Named identifiers syntactically cited in the proof block:}
\begin{lstlisting}[language=EasyCrypt,basicstyle=\ttfamily\scriptsize]
match_y
md
programming_site_matches_transcript
programming_site_output
programming_site_query
signature_programming_site_of_transcript
tr
transcript_valid_signature_challenge_matches
transcript_valid_signature_challenge_query
valid
\end{lstlisting}

\end{formaltheorem}

\begin{formaltheorem}[EasyCrypt line 1666]
\noindent\textbf{EasyCrypt name.}
\begin{lstlisting}[language=EasyCrypt,basicstyle=\ttfamily\scriptsize]
transcript_valid_signature_programming_site_no_reprogramming
\end{lstlisting}
\noindent\textbf{Checked EasyCrypt statement.}
\begin{lstlisting}[language=EasyCrypt,basicstyle=\ttfamily\scriptsize]
lemma transcript_valid_signature_programming_site_no_reprogramming
  md tr y :
  transcript_valid md tr =>
  transcript_signature_challenge_matches tr y =>
  ! reprogramming_failure md tr
      (signature_programming_site_of_transcript md tr y).
\end{lstlisting}
\noindent\textbf{Formal mathematical reading.}
Mathematical theorem. This is an implication, normally turning one invariant, validity predicate, or proof obligation into another.

\noindent\textbf{Role in the proof.}
Use in this chapter. It proves a structural correctness fact needed by later extraction or verification arguments.

\noindent\textbf{Proof-script interpretation.}
Proof-script reading. The machine proof decomposes the theorem into rewrites, program-logic obligations, relational equivalences, and arithmetic side conditions.
\emph{Main tactics visible in the proof script:} \ec{rewrite}.
\emph{Named identifiers syntactically cited in the proof block:}
\begin{lstlisting}[language=EasyCrypt,basicstyle=\ttfamily\scriptsize]
match_y
md
reprogramming_failure
tr
transcript_valid_signature_programming_site_matches
valid
\end{lstlisting}

\end{formaltheorem}

\begin{formaltheorem}[EasyCrypt line 1679]
\noindent\textbf{EasyCrypt name.}
\begin{lstlisting}[language=EasyCrypt,basicstyle=\ttfamily\scriptsize]
transcript_valid_actual_signature_programming_site_matches
\end{lstlisting}
\noindent\textbf{Checked EasyCrypt statement.}
\begin{lstlisting}[language=EasyCrypt,basicstyle=\ttfamily\scriptsize]
lemma transcript_valid_actual_signature_programming_site_matches md tr :
  transcript_valid md tr =>
  programming_site_matches_transcript md tr
    (actual_signature_programming_site md tr).
\end{lstlisting}
\noindent\textbf{Formal mathematical reading.}
Mathematical theorem. This is an implication, normally turning one invariant, validity predicate, or proof obligation into another.

\noindent\textbf{Role in the proof.}
Use in this chapter. It proves a structural correctness fact needed by later extraction or verification arguments.

\noindent\textbf{Proof-script interpretation.}
Proof-script reading. The machine proof decomposes the theorem into rewrites, program-logic obligations, relational equivalences, and arithmetic side conditions.
\emph{Main tactics visible in the proof script:} \ec{rewrite}, \ec{apply}.
\emph{Named identifiers syntactically cited in the proof block:}
\begin{lstlisting}[language=EasyCrypt,basicstyle=\ttfamily\scriptsize]
actual_signature_programming_site
md
tr
transcript_signature_challenge_output
transcript_signature_challenge_output_matches
transcript_valid_signature_programming_site_matches
valid
\end{lstlisting}

\end{formaltheorem}

\begin{formaltheorem}[EasyCrypt line 1691]
\noindent\textbf{EasyCrypt name.}
\begin{lstlisting}[language=EasyCrypt,basicstyle=\ttfamily\scriptsize]
transcript_valid_actual_signature_programming_site_no_reprogramming
\end{lstlisting}
\noindent\textbf{Checked EasyCrypt statement.}
\begin{lstlisting}[language=EasyCrypt,basicstyle=\ttfamily\scriptsize]
lemma transcript_valid_actual_signature_programming_site_no_reprogramming
  md tr :
  transcript_valid md tr =>
  ! reprogramming_failure md tr (actual_signature_programming_site md tr).
\end{lstlisting}
\noindent\textbf{Formal mathematical reading.}
Mathematical theorem. This is an implication, normally turning one invariant, validity predicate, or proof obligation into another.

\noindent\textbf{Role in the proof.}
Use in this chapter. It proves a structural correctness fact needed by later extraction or verification arguments.

\noindent\textbf{Proof-script interpretation.}
Proof-script reading. The machine proof decomposes the theorem into rewrites, program-logic obligations, relational equivalences, and arithmetic side conditions.
\emph{Main tactics visible in the proof script:} \ec{rewrite}.
\emph{Named identifiers syntactically cited in the proof block:}
\begin{lstlisting}[language=EasyCrypt,basicstyle=\ttfamily\scriptsize]
md
reprogramming_failure
tr
transcript_valid_actual_signature_programming_site_matches
valid
\end{lstlisting}

\end{formaltheorem}

\begin{formaltheorem}[EasyCrypt line 1702]
\noindent\textbf{EasyCrypt name.}
\begin{lstlisting}[language=EasyCrypt,basicstyle=\ttfamily\scriptsize]
valid_forking_pair_signature_programming_site_query
\end{lstlisting}
\noindent\textbf{Checked EasyCrypt statement.}
\begin{lstlisting}[language=EasyCrypt,basicstyle=\ttfamily\scriptsize]
lemma valid_forking_pair_signature_programming_site_query md tr1 tr2 y1 y2 :
  valid_forking_pair md tr1 tr2 =>
  programming_site_query (signature_programming_site_of_transcript md tr1 y1) =
  programming_site_query (signature_programming_site_of_transcript md tr2 y2).
\end{lstlisting}
\noindent\textbf{Formal mathematical reading.}
Mathematical theorem. This is an implication, normally turning one invariant, validity predicate, or proof obligation into another.

\noindent\textbf{Role in the proof.}
Use in this chapter. It proves a structural correctness fact needed by later extraction or verification arguments.

\noindent\textbf{Proof-script interpretation.}
Proof-script reading. The machine proof decomposes the theorem into rewrites, program-logic obligations, relational equivalences, and arithmetic side conditions.
\emph{Main tactics visible in the proof script:} \ec{rewrite}, \ec{apply}.
\emph{Named identifiers syntactically cited in the proof block:}
\begin{lstlisting}[language=EasyCrypt,basicstyle=\ttfamily\scriptsize]
fork_pair
programming_site_query
signature_programming_site_of_transcript
valid_forking_pair_signature_challenge_query
\end{lstlisting}

\end{formaltheorem}

\begin{formaltheorem}[EasyCrypt line 1713]
\noindent\textbf{EasyCrypt name.}
\begin{lstlisting}[language=EasyCrypt,basicstyle=\ttfamily\scriptsize]
valid_forking_pair_programming_site_query
\end{lstlisting}
\noindent\textbf{Checked EasyCrypt statement.}
\begin{lstlisting}[language=EasyCrypt,basicstyle=\ttfamily\scriptsize]
lemma valid_forking_pair_programming_site_query md tr1 tr2 y1 y2 :
  valid_forking_pair md tr1 tr2 =>
  programming_site_query (programming_site_of_transcript md tr1 y1) =
  programming_site_query (programming_site_of_transcript md tr2 y2).
\end{lstlisting}
\noindent\textbf{Formal mathematical reading.}
Mathematical theorem. This is an implication, normally turning one invariant, validity predicate, or proof obligation into another.

\noindent\textbf{Role in the proof.}
Use in this chapter. It proves a structural correctness fact needed by later extraction or verification arguments.

\noindent\textbf{Proof-script interpretation.}
Proof-script reading. The machine proof decomposes the theorem into rewrites, program-logic obligations, relational equivalences, and arithmetic side conditions.
\emph{Main tactics visible in the proof script:} \ec{rewrite}, \ec{apply}.
\emph{Named identifiers syntactically cited in the proof block:}
\begin{lstlisting}[language=EasyCrypt,basicstyle=\ttfamily\scriptsize]
fork_pair
programming_site_of_transcript
programming_site_query
valid_forking_pair_challenge_query
\end{lstlisting}

\end{formaltheorem}

\begin{formaltheorem}[EasyCrypt line 1723]
\noindent\textbf{EasyCrypt name.}
\begin{lstlisting}[language=EasyCrypt,basicstyle=\ttfamily\scriptsize]
valid_forking_pair_self_programming_sites_no_reprogramming
\end{lstlisting}
\noindent\textbf{Checked EasyCrypt statement.}
\begin{lstlisting}[language=EasyCrypt,basicstyle=\ttfamily\scriptsize]
lemma valid_forking_pair_self_programming_sites_no_reprogramming
  md tr1 tr2 y1 y2 :
  valid_forking_pair md tr1 tr2 =>
  transcript_challenge_matches tr1 y1 =>
  transcript_challenge_matches tr2 y2 =>
  ! reprogramming_failure md tr1
      (programming_site_of_transcript md tr1 y1) /\
  ! reprogramming_failure md tr2
      (programming_site_of_transcript md tr2 y2).
\end{lstlisting}
\noindent\textbf{Formal mathematical reading.}
Mathematical theorem. This is an implication, normally turning one invariant, validity predicate, or proof obligation into another.

\noindent\textbf{Role in the proof.}
Use in this chapter. It proves a structural correctness fact needed by later extraction or verification arguments.

\noindent\textbf{Proof-script interpretation.}
Proof-script reading. The machine proof decomposes the theorem into rewrites, program-logic obligations, relational equivalences, and arithmetic side conditions.
\emph{Main tactics visible in the proof script:} \ec{split}, \ec{apply}.
\emph{Named identifiers syntactically cited in the proof block:}
\begin{lstlisting}[language=EasyCrypt,basicstyle=\ttfamily\scriptsize]
match1
match2
programming_site_self_no_reprogramming_failure
\end{lstlisting}

\end{formaltheorem}

\begin{formaltheorem}[EasyCrypt line 1739]
\noindent\textbf{EasyCrypt name.}
\begin{lstlisting}[language=EasyCrypt,basicstyle=\ttfamily\scriptsize]
fs_bound_parts_nonnegative
\end{lstlisting}
\noindent\textbf{Checked EasyCrypt statement.}
\begin{lstlisting}[language=EasyCrypt,basicstyle=\ttfamily\scriptsize]
lemma fs_bound_parts_nonnegative :
  fs_with_aborts_bound_obligation =>
  0%r <= fs_with_aborts_reprogramming_term /\
  0%r <= fs_with_aborts_min_entropy_term.
\end{lstlisting}
\noindent\textbf{Formal mathematical reading.}
Mathematical theorem. This is an inequality, usually an arithmetic loss bound, event inclusion bound, or monotonicity fact used to compose probabilities.

\noindent\textbf{Role in the proof.}
Use in this chapter. It contributes directly to an advantage bound, bad-event bound, or real-arithmetic composition step.

\noindent\textbf{Proof-script interpretation.}
No proof block was collected for this item.

\end{formaltheorem}

\begin{formaltheorem}[EasyCrypt line 1745]
\noindent\textbf{EasyCrypt name.}
\begin{lstlisting}[language=EasyCrypt,basicstyle=\ttfamily\scriptsize]
fs_with_aborts_bound_obligation_holds
\end{lstlisting}
\noindent\textbf{Checked EasyCrypt statement.}
\begin{lstlisting}[language=EasyCrypt,basicstyle=\ttfamily\scriptsize]
lemma fs_with_aborts_bound_obligation_holds :
  fs_with_aborts_bound_obligation.
\end{lstlisting}
\noindent\textbf{Formal mathematical reading.}
Mathematical theorem. This lemma is a universally quantified proposition over the variables shown in the EasyCrypt statement.

\noindent\textbf{Role in the proof.}
Use in this chapter. It contributes directly to an advantage bound, bad-event bound, or real-arithmetic composition step.

\noindent\textbf{Proof-script interpretation.}
Proof-script reading. The machine proof decomposes the theorem into rewrites, program-logic obligations, relational equivalences, and arithmetic side conditions.
\emph{Main tactics visible in the proof script:} \ec{rewrite}, \ec{split}, \ec{apply}.
\emph{Named identifiers syntactically cited in the proof block:}
\begin{lstlisting}[language=EasyCrypt,basicstyle=\ttfamily\scriptsize]
fs_with_aborts_bound_obligation
fs_with_aborts_min_entropy_term_nonnegative
fs_with_aborts_reprogramming_term_nonnegative
\end{lstlisting}

\end{formaltheorem}

\medskip
\noindent\emph{End of generated formal inventory for \file{HAETAE_ROM_Programming.ec}.}
